% Contraction de texte : rédiger un résumé de texte — Français 1re Tech — Cours signé OnBuch+
% Programme officiel MINESEC. Compiler avec : tectonic N07-contraction-texte-rediger-resume-texte.tex
\newcommand{\DOCMATIERE}{Français}
\newcommand{\DOCNIVEAU}{1re Tech}
\newcommand{\DOCMODULE}{Français – classe de Première technique}
\newcommand{\DOCLECON}{N07}
\newcommand{\DOCTITRE}{Contraction de texte : rédiger un résumé de texte}
% OnBuch+ — préambule « cours signés » ENRICHI (compilé avec tectonic / XeLaTeX)
% Système visuel « Pop » d'OnBuch+ : fond crème (jamais blanc), bordures encre
% épaisses, ombres « dures » décalées, titres Archivo Black en capitales.
% Version MATHÉMATIQUES : graphes (pgfplots/tikz), boîtes pédagogiques
% (définition, propriété, méthode, attention, le savais-tu, exercice résolu,
% exercice, TP, bilan), page de garde et signature OnBuch+ ancrée en bas.
%
% Macros attendues dans main.tex AVANT \input{preamble} :
%   \DOCMATIERE  \DOCNIVEAU  \DOCMODULE  \DOCLECON (numéro)  \DOCTITRE
\documentclass[11pt]{article}

\usepackage{fontspec}
\usepackage[french]{babel}
\usepackage[a4paper,margin=2cm,top=3.4cm,bottom=2.8cm,headheight=1.4cm,footskip=1.3cm]{geometry}
\usepackage{amsmath,amssymb,amsfonts}
\usepackage{mathtools} % \xmapsto, \coloneqq... souvent utilisés par les modèles
\usepackage{cancel} % simplifications barrées (bilans de forces)
% \abs{z} (module, valeur absolue) et \norm{u} (norme), utilisés par les modèles
\providecommand{\abs}{}\let\abs\relax\DeclarePairedDelimiter{\abs}{\lvert}{\rvert}
\providecommand{\norm}{}\let\norm\relax\DeclarePairedDelimiter{\norm}{\lVert}{\rVert}
\usepackage[table]{xcolor} % [table] : fournit aussi \rowcolors (alternance de lignes)
\usepackage{pagecolor}
\usepackage{tikz}
\usetikzlibrary{arrows.meta,positioning,calc,shapes.geometric,shapes.misc,shapes.arrows,decorations.pathmorphing,decorations.pathreplacing,decorations.markings,patterns,shadows,mindmap,trees,fit,backgrounds,matrix,intersections,angles,quotes}
\usepackage{pgfplots}
\usepackage[version=4]{mhchem} % équations nucléaires \ce{^{235}_{92}U}
\usepackage[european,siunitx]{circuitikz} % schémas électriques
\pgfplotsset{compat=1.18}
\usepgfplotslibrary{fillbetween,groupplots} % aires entre courbes (calcul intégral)
\usepackage{enumitem}
\usepackage{fancyhdr}
% [most] charge toutes les bibliothèques tcolorbox (listings, minted, raster,
% documentation...) et gonfle inutilement la mémoire TeX sur un document long
% (~300 boîtes) : on ne charge que ce qui est réellement utilisé.
\usepackage[skins,breakable]{tcolorbox}
\usepackage{graphicx}
\usepackage{colortbl}
\usepackage{booktabs}
\usepackage{tabularx}
\usepackage{multirow}
\usepackage{diagbox} % case d'en-tête diagonale des tableaux à double entrée
% Colonnes extensibles centrée / gauche / droite (C, L, R), souvent utilisées par les modèles
\newcolumntype{C}{>{\centering\arraybackslash}X}
\newcolumntype{L}{>{\raggedright\arraybackslash}X}
\newcolumntype{R}{>{\raggedleft\arraybackslash}X}
\usepackage{array}
\usepackage{multicol}
\usepackage{siunitx}
% parse-numbers=false : accepte aussi \SI{2,5 \times 10^{-3}}{...}
\sisetup{locale=FR,output-decimal-marker={,},group-minimum-digits=4,parse-numbers=false}
\DeclareSIUnit\ohm{\ensuremath{\symup{\Omega}}} % U+2126 absent de la police
\DeclareSIUnit\division{div} % réglages d'oscilloscope (ms/div, V/div)
\DeclareSIUnit\tr{tr} % tours (vitesse de rotation, ex. \SI{1500}{\tr\per\minute})
\DeclareSIUnit\molar{M} % concentration molaire (ex. \SI{3}{\molar}, \SI{1}{\micro\molar})
\DeclareSIUnit\an{an} % année (le modèle confond parfois avec \year, un compteur TeX) — ex. \SI{5}{\kilo\meter\per\an}
\DeclareSIUnit\FCFA{FCFA} % franc CFA (ex. \SI{500}{\FCFA\per\kilo\gram})
\DeclareSIUnit\degreeBrix{\textdegree Brix} % teneur en sucre d'un sirop/jus (réfractométrie)
\DeclareSIUnit\month{mois} % le modèle utilise parfois \month, un compteur TeX, comme unité de durée
\DeclareSIUnit\microliter{\micro\liter} % le modèle écrit parfois \microliter en un seul token
\DeclareSIUnit\million{millions} % dénombrement (ex. \SI{15}{\million\per\mL} spermatozoïdes)
\usepackage{qrcode}
% \degree : commande fréquemment utilisée par les modèles pour un angle ou une
% température en dehors de \SI{}{} (non fournie par les paquets chargés).
\providecommand{\degree}{^{\circ}}
% \qed (fin de démonstration), utilisé par les modèles sans amsthm
\providecommand{\qed}{\hfill$\square$}
% \barre : double barre des valeurs interdites dans un tableau de variations (tkz-tab)
\providecommand{\barre}{\ensuremath{\|}}
% environnement proof (amsthm non chargé) utilisé par les modèles
\ifdefined\proof\else\newenvironment{proof}[1][Démonstration]{\par\noindent\textit{#1.}\ }{\qed\par}\fi
\usepackage{lastpage}
\usepackage{hyperref}
\hypersetup{hidelinks}

% ============================================================
% Palette « Pop » OnBuch+ — mêmes hex que la classe P (plus_ui.dart)
% ============================================================
\definecolor{poporange}{HTML}{F59321}
\definecolor{popdark}{HTML}{E07A0C}
\definecolor{popcream}{HTML}{FFF6E8}
\definecolor{popink}{HTML}{1C1714}
\definecolor{poppurple}{HTML}{7A5AE0}
\definecolor{popgreen}{HTML}{2E9E6A}
\definecolor{poppink}{HTML}{E0457B}
\definecolor{popblue}{HTML}{2F7DE1}
\definecolor{popgold}{HTML}{C98A12}
\definecolor{popmuted}{HTML}{8A7F76}
\definecolor{popline}{HTML}{EADFD2}
\definecolor{popgray}{HTML}{8A7F76}
\colorlet{popgrey}{popgray}
% Alias tolérants (le modèle nomme parfois les couleurs différemment)
\colorlet{popwhite}{white}
\colorlet{popblack}{popink}
\colorlet{popred}{poppink}
\colorlet{popyellow}{popgold}
% Teintes claires (fonds de tableaux/encadrés) — le modèle nomme parfois
% les couleurs avec un suffixe L (ex. popblueL) sans les avoir définies.
\colorlet{poporangeL}{poporange!20}
\colorlet{popdarkL}{popdark!20}
\colorlet{poppurpleL}{poppurple!20}
\colorlet{popgreenL}{popgreen!20}
\colorlet{poppinkL}{poppink!20}
\colorlet{popblueL}{popblue!20}
\colorlet{popgoldL}{popgold!20}
\colorlet{popredL}{popred!20}
\colorlet{popgrayL}{popgray!20}
% Teintes claires pour fonds de tableaux / zones
\colorlet{popblueL}{popblue!10!white}
\colorlet{poporangeL}{poporange!14!white}
\colorlet{popgreenL}{popgreen!12!white}
\colorlet{poppurpleL}{poppurple!12!white}
\colorlet{poppinkL}{poppink!12!white}
\colorlet{popgoldL}{popgold!14!white}

\pagecolor{popcream}

% ============================================================
% Typographie
% ============================================================
\defaultfontfeatures{Ligatures=TeX}
\setmainfont{PlusJakartaSans-Medium.ttf}[Path=../../../fonts/,BoldFont=PlusJakartaSans-Bold.ttf]
% Maths en sans-serif (Fira Math) avec chiffres et lettres droites de Plus
% Jakarta, pour que les formules se fondent dans le texte.
\usepackage[math-style=ISO,bold-style=ISO]{unicode-math}
\setmathfont{FiraMath-Regular.otf}[Path=../../../fonts/]
\setmathfont{PlusJakartaSans-Medium.ttf}[Path=../../../fonts/,range={up/num,up/{Latin,latin},\mathup}]
\usepackage{icomma} % virgule décimale sans espace en mode math
% Caractères mathématiques Unicode absents des polices de texte (Plus Jakarta,
% Archivo Black) : rendus par leur équivalent mathématique, en texte comme en
% formule — sinon ils s'affichent en carré vide (ex. « Arithmétique dans ℤ »).
\usepackage{newunicodechar}
\newunicodechar{ℝ}{\ensuremath{\mathbb{R}}}
\newunicodechar{ℤ}{\ensuremath{\mathbb{Z}}}
\newunicodechar{ℂ}{\ensuremath{\mathbb{C}}}
\newunicodechar{ℕ}{\ensuremath{\mathbb{N}}}
\newunicodechar{ℚ}{\ensuremath{\mathbb{Q}}}
\newunicodechar{𝔻}{\ensuremath{\mathbb{D}}}
\newunicodechar{α}{\ensuremath{\alpha}}
\newunicodechar{β}{\ensuremath{\beta}}
\newunicodechar{γ}{\ensuremath{\gamma}}
\newunicodechar{δ}{\ensuremath{\delta}}
\newunicodechar{ε}{\ensuremath{\varepsilon}}
\newunicodechar{θ}{\ensuremath{\theta}}
\newunicodechar{λ}{\ensuremath{\lambda}}
\newunicodechar{μ}{\ensuremath{\mu}}
\newunicodechar{σ}{\ensuremath{\sigma}}
\newunicodechar{τ}{\ensuremath{\tau}}
\newunicodechar{φ}{\ensuremath{\varphi}}
\newunicodechar{ψ}{\ensuremath{\psi}}
\newunicodechar{ω}{\ensuremath{\omega}}
\newunicodechar{Σ}{\ensuremath{\Sigma}}
% majuscules produites par \MakeUppercase dans les titres (α-aminés -> Α)
\newunicodechar{Α}{\ensuremath{\alpha}}
\newunicodechar{Β}{\ensuremath{\beta}}
\newunicodechar{Γ}{\ensuremath{\Gamma}}
\newunicodechar{Δ}{\ensuremath{\Delta}}
\newunicodechar{Ε}{\ensuremath{\varepsilon}}
\newunicodechar{Θ}{\ensuremath{\Theta}}
\newunicodechar{Λ}{\ensuremath{\Lambda}}
\newunicodechar{Μ}{\ensuremath{\mu}}
\newunicodechar{Τ}{\ensuremath{\tau}}
\newunicodechar{Φ}{\ensuremath{\Phi}}
\newunicodechar{Ψ}{\ensuremath{\Psi}}
\newunicodechar{Ω}{\ensuremath{\Omega}}
\newunicodechar{↦}{\ensuremath{\mapsto}}
\newunicodechar{∈}{\ensuremath{\in}}
\newunicodechar{∉}{\ensuremath{\notin}}
\newunicodechar{∧}{\ensuremath{\wedge}}
\newunicodechar{⇔}{\ensuremath{\Leftrightarrow}}
\newunicodechar{⟺}{\ensuremath{\Longleftrightarrow}}
\newunicodechar{⇒}{\ensuremath{\Rightarrow}}
\newunicodechar{⟹}{\ensuremath{\Longrightarrow}}
\newunicodechar{∘}{\ensuremath{\circ}}
\newunicodechar{∪}{\ensuremath{\cup}}
\newunicodechar{∩}{\ensuremath{\cap}}
\newunicodechar{≡}{\ensuremath{\equiv}}
\newunicodechar{⊂}{\ensuremath{\subset}}
\newunicodechar{⊥}{\ensuremath{\perp}}
\newunicodechar{∃}{\ensuremath{\exists}}
\newunicodechar{∀}{\ensuremath{\forall}}
\newunicodechar{∛}{\ensuremath{\sqrt[3]{\,}}}
\newunicodechar{∜}{\ensuremath{\sqrt[4]{\,}}}
\newunicodechar{⃗}{\ensuremath{{}^{\to}}}
\newunicodechar{ⁿ}{\ifmmode^{n}\else\textsuperscript{\ensuremath{n}}\fi}
\newunicodechar{ˣ}{\ifmmode^{x}\else\textsuperscript{\ensuremath{x}}\fi}
\newunicodechar{ᵇ}{\ifmmode^{b}\else\textsuperscript{\ensuremath{b}}\fi}
\newunicodechar{ʳ}{\ifmmode^{r}\else\textsuperscript{\ensuremath{r}}\fi}
\newunicodechar{ᵘ}{\ifmmode^{u}\else\textsuperscript{\ensuremath{u}}\fi}
\newunicodechar{ʸ}{\ifmmode^{y}\else\textsuperscript{\ensuremath{y}}\fi}
\newunicodechar{ᵃ}{\ifmmode^{a}\else\textsuperscript{\ensuremath{a}}\fi}
\newunicodechar{ᵉ}{\ifmmode^{e}\else\textsuperscript{\ensuremath{e}}\fi}
\newunicodechar{ᵏ}{\ifmmode^{k}\else\textsuperscript{\ensuremath{k}}\fi}
\newunicodechar{ᵖ}{\ifmmode^{p}\else\textsuperscript{\ensuremath{p}}\fi}
\newunicodechar{ᴺ}{\ifmmode^{N}\else\textsuperscript{\ensuremath{N}}\fi}
\newunicodechar{ᶜ}{\ifmmode^{c}\else\textsuperscript{\ensuremath{c}}\fi}
\newunicodechar{ʰ}{\ifmmode^{h}\else\textsuperscript{\ensuremath{h}}\fi}
\newunicodechar{ᵠ}{\ifmmode^{q}\else\textsuperscript{\ensuremath{q}}\fi}
\newunicodechar{ₙ}{\ifmmode_{n}\else\textsubscript{\ensuremath{n}}\fi}
\newunicodechar{ₐ}{\ifmmode_{a}\else\textsubscript{\ensuremath{a}}\fi}
\newunicodechar{ᵢ}{\ifmmode_{i}\else\textsubscript{\ensuremath{i}}\fi}
\newunicodechar{ᵣ}{\ifmmode_{r}\else\textsubscript{\ensuremath{r}}\fi}
\newunicodechar{ₖ}{\ifmmode_{k}\else\textsubscript{\ensuremath{k}}\fi}
\newunicodechar{ᵦ}{\ifmmode_{\beta}\else\textsubscript{\ensuremath{\beta}}\fi}
\newunicodechar{ₓ}{\ifmmode_{x}\else\textsubscript{\ensuremath{x}}\fi}
\newfontfamily\popdisplay{ArchivoBlack-Regular.ttf}[Path=../../../fonts/]
\newfontfamily\poplogo{SpaceGrotesk-Bold.ttf}[Path=../../../fonts/]
\newfontfamily\popsemi{PlusJakartaSans-SemiBold.ttf}[Path=../../../fonts/]
\newfontfamily\popxbold{PlusJakartaSans-ExtraBold.ttf}[Path=../../../fonts/]
\newfontfamily\popxboldls{PlusJakartaSans-ExtraBold.ttf}[Path=../../../fonts/,LetterSpace=3.0]

\setlist[enumerate]{leftmargin=1.7em,itemsep=5pt,topsep=4pt}
\setlist[itemize]{leftmargin=1.7em,itemsep=4pt,topsep=4pt,label={\color{poporange}\textbullet}}
\setlength{\parindent}{0pt}
\setlength{\parskip}{5pt}
\renewcommand{\arraystretch}{1.25}

% pgfplots : style Pop par défaut
\pgfplotsset{
  every axis/.append style={axis line style={popink,line width=1pt},tick style={popink},
    grid style={popline,line width=0.5pt},label style={font=\small},tick label style={font=\footnotesize}},
  popcurve/.style={poporange,line width=1.8pt,no markers,smooth},
}
% Même style disponible dans un \draw TikZ simple (hors pgfplots)
\tikzset{popcurve/.style={poporange,line width=1.8pt}}
\tikzset{popgrid/.style={popline,very thin}}
\colorlet{popcurve}{poporange} % le nom du style est parfois utilisé comme couleur

% ============================================================
% Pill / squiggle
% ============================================================
\newcommand{\poppill}[3]{%
  \tikz[baseline=-0.55ex]\node[draw=popink,line width=1.2pt,rounded corners=2.6pt,%
    fill=#2,inner xsep=6.5pt,inner ysep=3pt]{%
    \popxboldls\fontsize{7}{7}\selectfont\color{#3}\MakeUppercase{#1}};%
}
\newcommand{\popsquiggle}[1]{%
  \tikz[baseline=-0.4ex]{
    \draw[#1,line width=2.6pt,line cap=round]
      plot [smooth] coordinates {(0,0.09) (0.34,0.01) (0.68,0.21) (1.02,0.01) (1.36,0.10)};
  }%
}
% Mot-clé surligné (marqueur orange sous le texte)
\newcommand{\cle}[1]{{\popxbold\color{popdark}#1}}

% ============================================================
% En-tête / pied
% ============================================================
\pagestyle{fancy}
\fancyhf{}
\renewcommand{\headrulewidth}{0pt}
\renewcommand{\footrulewidth}{0pt}
\lhead{\raisebox{-0.15ex}{\poplogo\fontsize{13}{13}\selectfont\color{popink}OnBuch\color{poporange}+}}
\rhead{\poppill{\DOCMATIERE\ \textbullet\ \DOCNIVEAU\ \textbullet\ Leçon \DOCLECON}{white}{popink}}
\lfoot{\popxbold\fontsize{7.6}{7.6}\selectfont\color{popmuted}\MakeUppercase{Cours signé OnBuch+}\ \textbullet\ {\popsemi\DOCTITRE}}
\rfoot{\tikz[baseline=-0.6ex]\node[rounded corners=8.5pt,fill=popink,minimum height=17pt,minimum width=17pt,inner xsep=5pt,inner sep=0]{\popxbold\fontsize{8}{8}\selectfont\color{popcream}\thepage\,/\,\pageref*{LastPage}};}

% ============================================================
% Titres
% ============================================================
\newcommand{\chaptitre}[2]{%
  \begin{tcolorbox}[enhanced,colback=poporange,colframe=popink,boxrule=2.6pt,arc=11pt,
      drop shadow=popink,left=15pt,right=15pt,top=12pt,bottom=11pt,before skip=3pt,after skip=18pt]
    \poppill{#2}{popink}{popcream}\par\vspace{9pt}
    {\raggedright\hyphenpenalty=10000\exhyphenpenalty=10000\popdisplay\fontsize{22}{24}\selectfont\color{popcream}\MakeUppercase{#1}\par}
  \end{tcolorbox}
}
% Section numérotée : pastille orange avec le numéro + titre Archivo
\newcounter{popsec}
\newcommand{\coursec}[1]{%
  \stepcounter{popsec}\setcounter{popsub}{0}%
  \par\vspace{16pt}\noindent
  \begin{minipage}{\linewidth}
  \tikz[baseline=-0.7ex]\node[draw=popink,line width=1.6pt,fill=poporange,rounded corners=4pt,minimum size=22pt,inner sep=0]{\popdisplay\fontsize{12}{12}\selectfont\color{popcream}\thepopsec};%
  \hspace{8pt}{\popdisplay\fontsize{15}{17}\selectfont\color{popink}\MakeUppercase{#1}}\par
  \vspace{4pt}\noindent{\color{poporange}\rule{36pt}{3.6pt}}
  \end{minipage}\par\nopagebreak\vspace{8pt}
}
\newcounter{popsub}
\newcommand{\courssub}[1]{%
  \stepcounter{popsub}%
  \par\vspace{9pt}\noindent{\popxbold\fontsize{11.5}{13}\selectfont\color{popink}{\color{poporange}\thepopsec.\thepopsub}\hspace{6pt}#1}\par\nopagebreak\vspace{4pt}
}

% ============================================================
% Boîtes pédagogiques — même recette : fond blanc, bordure épaisse colorée,
% ombre dure encre, badge plein. Titre optionnel [#1].
% ============================================================
\tcbset{popbase/.style={breakable,enhanced,colback=white,boxrule=2.3pt,arc=9pt,
  shadow={3pt}{-3pt}{0pt}{popink},left=14pt,right=14pt,top=11pt,bottom=11pt,before skip=11pt,after skip=15pt,
  fontupper=\popsemi\fontsize{10.6}{14.5}\selectfont\color{popink},
  fontlower=\popsemi\fontsize{10.6}{14.5}\selectfont\color{popink}}}
% Coche dessinée (le glyphe ✓ n'existe pas dans les polices du document)
\newcommand{\popcheck}[1]{\tikz[baseline=-0.5ex]\draw[#1,line width=1.6pt,line cap=round,line join=round](-0.13,0.0)--(-0.04,-0.09)--(0.13,0.1);}
\providecommand{\coche}{\popcheck{popgreen}}
\providecommand{\resultat}[1]{\par\smallskip\noindent\fcolorbox{popgreen}{popgreenL}{\parbox{\dimexpr\linewidth-2\fboxsep-2\fboxrule\relax}{\textbf{Résultat.} #1}}\par\smallskip}
\newcommand{\popboxhead}[3]{\poppill{#1}{#2}{white}\ifx\relax#3\relax\else\hspace{6pt}{\popxbold\fontsize{10.6}{12}\selectfont\color{popink}#3}\fi\par\vspace{7pt}}

\newtcolorbox{aretenir}[1][]{popbase,colframe=popgreen,before upper={\popboxhead{À retenir}{popgreen}{#1}}}
\newtcolorbox{exemplebox}[1][]{popbase,colframe=poporange,before upper={\popboxhead{Exemple}{poporange}{#1}}}
\newtcolorbox{definition}[1][]{popbase,colframe=popblue,before upper={\popboxhead{Définition}{popblue}{#1}}}
\newtcolorbox{propriete}[1][]{popbase,colframe=poppurple,before upper={\popboxhead{Propriété}{poppurple}{#1}}}
\newtcolorbox{methode}[1][]{popbase,colframe=popgold,before upper={\popboxhead{Méthode}{popgold}{#1}}}
\newtcolorbox{attention}[1][]{popbase,colframe=poppink,before upper={\popboxhead{Attention}{poppink}{#1}}}
\newtcolorbox{experience}[1][]{popbase,colframe=popblue,colback=popblueL,before upper={\popboxhead{Expérience}{popblue}{#1}}}
% « Le savais-tu ? » : carte violette pleine (contexte, culture, Cameroun)
\newtcolorbox{savaistu}[1][]{popbase,colback=poppurple,colframe=popink,
  fontupper=\popsemi\fontsize{10.4}{14.2}\selectfont\color{white},
  before upper={\poppill{Le savais-tu\,?}{white}{poppurple}\ifx\relax#1\relax\else\hspace{6pt}{\popxbold\color{white}#1}\fi\par\vspace{7pt}}}
% Exercice résolu : énoncé en haut, \tcblower puis la solution sur fond crème
\newcounter{popexo}\newcounter{popexr}
\newtcolorbox{exoresolu}[1][]{popbase,colframe=poporange,colbacklower=popcream,
  segmentation style={popink,line width=1.2pt,dashed},
  before upper={\stepcounter{popexr}\popboxhead{Exercice résolu \thepopexr}{poporange}{#1}},
  before lower={\poppill{Solution}{popink}{popcream}\par\vspace{6pt}}}
% Exercice (non résolu en place) — la correction est dans la partie Corrigés
\newtcolorbox{exercice}[1][]{popbase,colframe=popink,
  before upper={\stepcounter{popexo}\popboxhead{Exercice \thepopexo}{popink}{#1}}}
% Corrigé
\newtcolorbox{corrige}[1][]{popbase,colframe=popgreen,colback=popgreenL,
  before upper={\popboxhead{Corrigé}{popgreen}{#1}}}
% Objectifs / prérequis / bilan
\newtcolorbox{objectifs}{popbase,colframe=popink,colback=poporangeL,
  before upper={\popboxhead{Objectifs de la leçon}{popink}{}}}
\newtcolorbox{prerequis}{popbase,colframe=popink,colback=popblueL,
  before upper={\popboxhead{Prérequis}{popblue}{}}}
\newtcolorbox{bilan}{popbase,colframe=popink,colback=white,boxrule=2.8pt,
  before upper={\popboxhead{Fiche bilan}{poporange}{L'essentiel à connaître pour l'examen}}}
% Figure encadrée avec légende
\newtcolorbox{popfigure}[1][]{enhanced,breakable=false,colback=white,colframe=popline,boxrule=1.4pt,arc=8pt,
  left=10pt,right=10pt,top=10pt,bottom=8pt,before skip=10pt,after skip=12pt,halign=center,
  lower separated=false,halign lower=center,
  fontlower=\popsemi\fontsize{9}{11.5}\selectfont\color{popmuted},
  #1}
\newcommand{\legende}[1]{\par\vspace{6pt}{\popsemi\fontsize{9}{11.5}\selectfont\color{popmuted}#1}}

% ============================================================
% Signature OnBuch+
% ============================================================
\newcommand{\poplogoinline}[1]{{\poplogo\fontsize{#1}{#1}\selectfont\color{popink}OnBuch\color{poporange}+}}
\newcommand{\signatureonbuch}{%
  \begin{tcolorbox}[enhanced,colback=popink,colframe=popink,boxrule=2.2pt,arc=10pt,
      drop shadow=poporange,left=14pt,right=14pt,top=10pt,bottom=10pt,before skip=0pt,after skip=0pt]
    \tikz[baseline=-0.7ex]\node[circle,fill=popgreen,draw=popcream,line width=1.3pt,minimum size=22pt,inner sep=0]{\popcheck{white}};%
    \hspace{9pt}\begin{minipage}[c]{0.62\linewidth}
      {\popxbold\fontsize{12}{13}\selectfont\color{popcream}Signé\ \textbullet\ {\poplogo OnBuch\color{poporange}+}}\par\vspace{2pt}
      {\raggedright\popsemi\fontsize{8}{10.5}\selectfont\color{popline}\MakeUppercase{Équipe pédagogique OnBuch+}\\\MakeUppercase{Conforme au programme officiel MINESEC}\par}
    \end{minipage}\hfill
    {\poplogo\fontsize{22}{22}\selectfont\color{popcream}OnBuch\color{poporange}+}
  \end{tcolorbox}%
}

% ============================================================
% Page de garde
% #1 titre · #2 matière · #3 niveau · #4 module · #5 n° leçon · #6 durée ·
% #7 description / objectifs (texte ou itemize)
% ============================================================
\newcommand{\pagegarde}[7]{%
  \thispagestyle{empty}
  \newgeometry{a4paper,margin=1.7cm,top=1.6cm,bottom=1.4cm}%
  \begin{tikzpicture}[remember picture,overlay]
    \fill[poporange!18] ([xshift=-3.2cm,yshift=-3.6cm]current page.north east) circle (4.6cm);
    \fill[poppurple!14] ([xshift=2.4cm,yshift=7.6cm]current page.south west) circle (3.8cm);
  \end{tikzpicture}%
  {\poplogo\fontsize{30}{30}\selectfont\color{popink}OnBuch\color{poporange}+}\hfill
  \raisebox{0.6ex}{\poppill{Cours signé \textbullet\ Premium}{popink}{popcream}}\par
  \vspace{1pt}\popsquiggle{poppurple}\par
  \vspace{8pt}
  \begin{tcolorbox}[enhanced,colback=poporange,colframe=popink,boxrule=2.8pt,arc=13pt,
      drop shadow=popink,left=18pt,right=18pt,top=12pt,bottom=13pt,after skip=10pt]
    \poppill{#2 \textbullet\ #3}{popink}{popcream}\hspace{5pt}\poppill{#4}{white}{popink}\par\vspace{8pt}
    {\popdisplay\fontsize{14}{14}\selectfont\color{popink}LEÇON #5}\par\vspace{4pt}
    {\raggedright\hyphenpenalty=10000\exhyphenpenalty=10000\popdisplay\fontsize{25}{27}\selectfont\color{popcream}\MakeUppercase{#1}\par}
    \vspace{8pt}
    {\popxbold\fontsize{9.5}{9.5}\selectfont\color{popink}\MakeUppercase{Durée conseillée : #6}}
  \end{tcolorbox}
  \begin{tcolorbox}[enhanced,colback=white,colframe=popink,boxrule=2.2pt,arc=10pt,
      drop shadow=popink,left=16pt,right=16pt,top=10pt,bottom=10pt,after skip=8pt]
    \poppill{Dans cette leçon}{poppurple}{white}\par\vspace{5pt}
    \popsemi\fontsize{9.3}{12.6}\selectfont\color{popink}#7
  \end{tcolorbox}
  \noindent\begin{minipage}[t]{0.487\linewidth}
    \begin{tcolorbox}[enhanced,colback=popgreen,colframe=popink,boxrule=2.2pt,arc=10pt,
        drop shadow=popink,left=11pt,right=11pt,top=10pt,bottom=10pt,height=3.1cm,valign=center]
      \tcbox[colback=white,colframe=popink,boxrule=1.3pt,arc=5pt,boxsep=0pt,nobeforeafter,
        left=3pt,right=3pt,top=3pt,bottom=3pt,valign=center]{\qrcode[height=1.9cm]{https://play.google.com/store/apps/details?id=com.luvvix.onbuch&hl=fr}}%
      \hspace{8pt}\begin{minipage}[c]{0.5\linewidth}
        {\popdisplay\fontsize{11}{12.5}\selectfont\color{white}\MakeUppercase{Télécharge OnBuch}}\par\vspace{4pt}
        {\raggedright\popsemi\fontsize{8.4}{11}\selectfont\color{white}Gratuit \textbullet\ Google Play\\Tuteur IA, annales, concours, résultats\par}
      \end{minipage}
    \end{tcolorbox}
  \end{minipage}\hfill
  \begin{minipage}[t]{0.487\linewidth}
    \begin{tcolorbox}[enhanced,colback=poppurple,colframe=popink,boxrule=2.2pt,arc=10pt,
        drop shadow=popink,left=11pt,right=11pt,top=10pt,bottom=10pt,height=3.1cm,valign=center]
      \tikz[baseline=-0.7ex]\node[circle,fill=white,draw=popink,line width=1.3pt,minimum size=1.6cm,inner sep=0]{\popdisplay\fontsize{24}{24}\selectfont\color{poppurple}+};%
      \hspace{8pt}\begin{minipage}[c]{0.6\linewidth}
        {\popdisplay\fontsize{11}{12.5}\selectfont\color{white}\MakeUppercase{Passe à OnBuch+}}\par\vspace{4pt}
        {\raggedright\popsemi\fontsize{8.4}{11}\selectfont\color{white}Tous les cours signés, Tuteur IA illimité, fiches PDF — onglet OnBuch+ de l'app\par}
      \end{minipage}
    \end{tcolorbox}
  \end{minipage}
  \vspace{8pt}
  \popcontenu
  \vfill
  \signatureonbuch
  \restoregeometry
  \newpage
}

% Rangée « Ce cours contient » (page de garde)
\newcommand{\poptile}[3]{% #1 couleur #2 gros texte #3 légende
  \begin{tcolorbox}[enhanced,colback=#1,colframe=popink,boxrule=2pt,arc=9pt,shadow={2.5pt}{-2.5pt}{0pt}{popink},
      left=6pt,right=6pt,top=9pt,bottom=9pt,halign=center,valign=center,height=1.75cm,width=\linewidth,nobeforeafter]
    {\popdisplay\fontsize{12.5}{13}\selectfont\color{white}\MakeUppercase{#2}}\par\vspace{3pt}
    {\centering\popsemi\fontsize{7.8}{9.5}\selectfont\color{white}#3\par}
  \end{tcolorbox}}
\newcommand{\popcontenu}{%
  \par\noindent{\popxboldls\fontsize{7.5}{7.5}\selectfont\color{popmuted}CE COURS SIGNÉ CONTIENT}\par\vspace{7pt}
  \noindent\begin{minipage}[t]{0.235\linewidth}\poptile{popblue}{Cours}{complet, illustré, pas à pas}\end{minipage}\hfill
  \begin{minipage}[t]{0.235\linewidth}\poptile{poporange}{Méthodes}{et exercices résolus}\end{minipage}\hfill
  \begin{minipage}[t]{0.235\linewidth}\poptile{poppink}{Exercices}{type examen + corrigés}\end{minipage}\hfill
  \begin{minipage}[t]{0.235\linewidth}\poptile{popgreen}{Fiche bilan}{l'essentiel à retenir}\end{minipage}\par}

% Sommaire
\newtcolorbox{sommaire}{enhanced,colback=white,colframe=popink,boxrule=2.2pt,arc=10pt,
  drop shadow=popink,left=18pt,right=18pt,top=13pt,bottom=13pt,before skip=6pt,after skip=20pt,
  before upper={\poppill{Sommaire}{poppurple}{white}\par\vspace{9pt}\popsemi\fontsize{10.6}{14.5}\selectfont\color{popink}}}

% Signature de fin de cours — poussée tout en bas de la dernière page
\newcommand{\findecours}{%
  \par\vfill\nopagebreak
  \begin{center}\popsquiggle{poporange}\end{center}\vspace{4pt}
  \signatureonbuch
}
\AtBeginDocument{\def\llbracket{\mathopen{[\mkern-3.5mu[}}\def\rrbracket{\mathclose{]\mkern-3.5mu]}}} % intervalles d'entiers (glyphe ⟦ absent de la police)
% Glyphes absents de Fira Math : redéfinis à partir de glyphes disponibles
\AtBeginDocument{%
  \def\setminus{\mathbin{\backslash}}\let\smallsetminus\setminus
  \def\vdots{\mathord{\vbox{\baselineskip3.2pt\lineskiplimit0pt\kern4pt\hbox{.}\hbox{.}\hbox{.}}}}%
  \def\ddots{\mathinner{\mkern1mu\raise6.5pt\hbox{.}\mkern2mu\raise3.6pt\hbox{.}\mkern2mu\raise0.7pt\hbox{.}\mkern1mu}}%
  \def\triangle{\mathord{\scalebox{0.9}{$\symup{\Delta}$}}}%
  \def\nabla{\mathord{\raisebox{\depth}{\scalebox{1}[-1]{$\symup{\Delta}$}}}}%
  \def\checkmark{\popcheck{popdark}}%
  \def\star{\mathbin{\ast}}\def\bigstar{\mathord{\ast}}%
  \def\diamond{\mathbin{\scalebox{0.6}{$\Diamond$}}}%
  \def\bigcup{\mathop{\mathchoice{\raisebox{-0.3em}{\scalebox{1.7}{$\cup$}}}{\scalebox{1.25}{$\cup$}}{\cup}{\cup}}\displaylimits}%
  \def\bigcap{\mathop{\mathchoice{\raisebox{-0.3em}{\scalebox{1.7}{$\cap$}}}{\scalebox{1.25}{$\cap$}}{\cap}{\cap}}\displaylimits}%
  \def\lesseqqgtr{\mathrel{\lessgtr}}\def\bigoplus{\mathop{\oplus}\displaylimits}%
}
\providecommand{\ssi}{si et seulement si} % abréviation fréquente
\AtBeginDocument{\providecommand{\wideparen}{\overparen}} % arc de cercle

%% FIGFIX-BEGIN v5 — correctifs globaux des figures (docs/onbuch-generation/tools/figfix)
\usepackage{adjustbox}\usepackage{etoolbox}\tcbuselibrary{raster}
% toute figure TikZ/pgfplots plus large que la ligne est réduite à la largeur disponible
\BeforeBeginEnvironment{tikzpicture}{\begin{adjustbox}{max width=\linewidth}}
\AfterEndEnvironment{tikzpicture}{\end{adjustbox}}
% légendes pgfplots lisibles par-dessus les courbes
\pgfplotsset{every axis/.append style={legend style={fill=white,fill opacity=0.92,text opacity=1,draw=black!15,inner sep=3pt}}}
% étiquettes d'arêtes (syntaxe edge["..."]) avec fond blanc
\tikzset{every edge quotes/.append style={fill=white,inner sep=1.2pt}}
%% FIGFIX-END
\begin{document}

\pagegarde{Contraction de texte : rédiger un résumé de texte}{Français}{1re Tech}{Français – classe de Première technique}{N07}{19 séances (fiche de progression harmonisée)}{Cette leçon outille l'élève de Première technique pour maîtriser la contraction de texte : repérage de la structure argumentative et du plan d'un texte, techniques de réduction et de reformulation, rédaction d'un résumé cohérent et personnel. Elle s'appuie sur des textes explicatifs et argumentatifs en lien avec des réalités camerounaises et sur l'étude de l'œuvre Le Lion et la perle de Wole Soyinka.
\vspace{4pt}
\begin{multicols}{2}\small
\begin{itemize}
\item À la fin de cette leçon, je sais identifier la thèse, les arguments et les exemples d'un texte argumentatif ou explicatif.
\item À la fin de cette leçon, je sais dégager le plan d'un texte en repérant les étapes de la progression (linéaire, dialectique, thématique).
\item À la fin de cette leçon, je sais appliquer les techniques de réduction : suppression, généralisation, condensation, substitution.
\item À la fin de cette leçon, je sais reformuler les idées essentielles avec mes propres mots sans trahir la pensée de l'auteur.
\item À la fin de cette leçon, je sais rédiger un résumé respectant la consigne de longueur (environ le quart du texte), la cohésion et la cohérence.
\item À la fin de cette leçon, je sais assurer la cohésion de mon résumé grâce aux connecteurs logiques, aux reprises pronominales et au champ lexical.
\end{itemize}\end{multicols}}

\begin{sommaire}
\begin{multicols}{2}
\begin{enumerate}
\item Découverte : le texte explicatif et la contraction
\item La structure argumentative et le plan du texte
\item Cohésion, cohérence et structure de la phrase
\item Techniques de réduction et de reformulation
\item Les figures d'atténuation : litote, euphémisme, prétérition
\item Rédaction, confrontation et amélioration du résumé
\item Ouverture littéraire : Le Lion et la perle — bilan de l'œuvre
\item Activité d'intégration
\item Méthodes et exercices résolus
\item Exercices
\item Corrigés des exercices
\item Fiche bilan
\end{enumerate}
\end{multicols}
\end{sommaire}

\begin{objectifs}
\begin{itemize}
\item À la fin de cette leçon, je sais identifier la thèse, les arguments et les exemples d'un texte argumentatif ou explicatif.
\item À la fin de cette leçon, je sais dégager le plan d'un texte en repérant les étapes de la progression (linéaire, dialectique, thématique).
\item À la fin de cette leçon, je sais appliquer les techniques de réduction : suppression, généralisation, condensation, substitution.
\item À la fin de cette leçon, je sais reformuler les idées essentielles avec mes propres mots sans trahir la pensée de l'auteur.
\item À la fin de cette leçon, je sais rédiger un résumé respectant la consigne de longueur (environ le quart du texte), la cohésion et la cohérence.
\item À la fin de cette leçon, je sais assurer la cohésion de mon résumé grâce aux connecteurs logiques, aux reprises pronominales et au champ lexical.
\item À la fin de cette leçon, je sais distinguer phrase simple et phrase composée pour varier et clarifier mes formulations.
\item À la fin de cette leçon, je sais reconnaître et employer les figures d'atténuation (litote, euphémisme, prétérition) dans un texte.
\item À la fin de cette leçon, je sais confronter ma production à celle d'autrui, l'évaluer selon une grille et l'améliorer.
\end{itemize}
\end{objectifs}

\begin{prerequis}
\begin{itemize}
\item Notions de phrase simple et phrase complexe vues en classe de Seconde
\item Les connecteurs logiques et les articulateurs du texte (classe de Seconde)
\item La lecture méthodique d'un texte argumentatif : thèse, arguments, exemples
\item L'étude de l'œuvre Le Lion et la perle : personnages, thèmes et contexte (leçons précédentes de l'année)
\item Les types de textes (narratif, descriptif, explicatif, argumentatif)
\end{itemize}
\end{prerequis}

\newpage

\coursec{Découverte : le texte explicatif et la contraction}

\courssub{Situation d'accroche : pourquoi apprendre à résumer ?}

À la veille des examens, Amina, élève de Première technique à Yaoundé, doit réviser un chapitre de 12 pages sur l'urbanisation du Cameroun : elle n'a que deux heures. Comment faire pour retenir l'essentiel sans tout recopier ? Au ministère de la Communication, un attaché de presse doit résumer en 15 lignes un discours de 10 pages pour le journal télévisé de 20 heures. À l'examen du Probatoire technique comme au Baccalauréat technique, l'épreuve de contraction de texte exige de restituer en un quart du texte initial les idées essentielles d'un document, sans ajouter son opinion.

Ces trois situations ont un point commun : il faut \cle{réduire} un texte sans le \emph{trahir}. Ni Amina, ni l'attaché de presse, ni le candidat au Baccalauréat n'ont le droit d'inventer, de commenter ou de critiquer : ils doivent rester fidèles à l'auteur tout en étant beaucoup plus brefs que lui.

\textbf{Problématique :} quelles techniques permettent de passer d'un texte long à un résumé fidèle, clair et cohérent ?

Pour répondre à cette question, nous allons d'abord définir la contraction de texte, puis découvrir le type de texte le plus souvent proposé à l'examen, le \cle{texte explicatif}, avant d'apprendre à repérer ses idées essentielles sur un exemple camerounais.

\courssub{Qu'est-ce que la contraction de texte ?}

\textbf{Intuition.} Imagine un entonnoir : tu verses en haut un grand volume de liquide, et il en sort en bas un filet concentré. La contraction de texte fonctionne de la même manière : le texte long entre par le haut, on retire tout ce qui est superflu (exemples, répétitions, détails), et il ne reste que le \emph{concentré} des idées.

\begin{definition}[La contraction de texte]
La \cle{contraction de texte} est une opération qui consiste à \textbf{réduire un texte à une longueur imposée} (généralement le \textbf{quart} du texte initial) \textbf{en en conservant les idées essentielles}, c'est-à-dire le thème, la thèse et les arguments principaux, \textbf{sans aucun commentaire personnel}. Le résumé obtenu doit être rédigé en phrases complètes, cohérentes et correctes.
\end{definition}

Trois contraintes sont donc indissociables :
\begin{enumerate}
\item le \textbf{respect du nombre de mots} : pour un texte de 400 mots, le résumé doit compter environ 100 mots (une marge de $\pm 10\,\%$ est tolérée) ;
\item la \textbf{fidélité au texte} : on ne dit rien que l'auteur n'ait dit, on ne change pas ses idées ;
\item la \textbf{neutralité} : on n'exprime jamais son opinion personnelle (interdit : \og je pense que \fg{}, \og à mon avis \fg{}, \og l'auteur a raison \fg{}).
\end{enumerate}

\begin{popfigure}
\begin{center}
\begin{tikzpicture}[scale=1]
% Entonnoir
\fill[popblue!25] (-3,3) -- (3,3) -- (0.6,0.6) -- (0.6,-0.6) -- (-0.6,-0.6) -- (-0.6,0.6) -- cycle;
\draw[popblue,very thick] (-3,3) -- (3,3);
\draw[popblue,very thick] (-3,3) -- (-0.6,0.6) -- (-0.6,-0.6);
\draw[popblue,very thick] (3,3) -- (0.6,0.6) -- (0.6,-0.6);
\draw[popblue,very thick] (-0.6,-0.6) -- (0.6,-0.6);
% Étiquettes
\node[align=center,font=\bfseries] at (0,3.6) {Texte initial : 400 mots};
\node[align=center,text=popdark] at (0,1.7) {\footnotesize On élimine : répétitions,\\ \footnotesize exemples, détails, citations};
\node[align=center,text=poporange,font=\bfseries] at (0,0.15) {\footnotesize Idées essentielles};
\draw[-{Stealth},very thick,poporange] (0,-0.7) -- (0,-1.4);
\node[align=center,font=\bfseries,text=popgreen!60!black] at (0,-1.9) {Résumé : 100 mots};
% Flèches latérales de ce qui est évacué
\draw[-{Stealth},thick,popmuted] (-2.2,2.2) -- (-4,2.2) node[left,align=right,font=\footnotesize] {exemples\\ chiffres secondaires};
\draw[-{Stealth},thick,popmuted] (2.2,1.4) -- (4,1.4) node[right,align=left,font=\footnotesize] {répétitions\\ reformulations};
\end{tikzpicture}
\end{center}
\legende{Le principe de l'entonnoir : d'un texte de 400 mots, on ne conserve que les idées essentielles pour rédiger un résumé d'environ 100 mots (le quart).}
\end{popfigure}

\begin{attention}[Résumé, compte rendu, synthèse : ne pas confondre]
Ces trois exercices sont proches mais distincts :
\begin{itemize}
\item le \textbf{résumé} (ou contraction de texte) réduit \emph{un seul} texte au quart de sa longueur, sans opinion personnelle ;
\item le \textbf{compte rendu} restitue de manière fidèle et objective le contenu d'un texte, d'une réunion ou d'un événement, sans contrainte stricte de longueur ;
\item la \textbf{synthèse} confronte \emph{plusieurs} documents sur un même sujet pour en dégager les convergences et les divergences.
\end{itemize}
Au Probatoire et au Baccalauréat technique, c'est le \textbf{résumé} qui est demandé : fidélité totale à l'auteur, aucune opinion, nombre de mots imposé.
\end{attention}

\courssub{Le texte explicatif : le texte type de l'épreuve}

\textbf{Intuition.} Quand ton professeur de sciences explique \emph{pourquoi} la pression augmente dans un pneu chaud, il ne raconte pas une histoire et ne cherche pas à te convaincre d'une opinion : il te fait \emph{comprendre} un phénomène. C'est exactement le rôle du texte explicatif.

\begin{definition}[Le texte explicatif]
Le \cle{texte explicatif} est un texte à \textbf{visée informative} qui fait \textbf{comprendre un phénomène} en en donnant les \textbf{causes}, les \textbf{mécanismes} et les \textbf{conséquences}. Il répond aux questions : \emph{pourquoi ?} et \emph{comment ?}
\end{definition}

\begin{propriete}[Caractéristiques du texte explicatif]
On reconnaît un texte explicatif aux indices suivants :
\begin{enumerate}
\item une \textbf{visée informative et objective} : l'auteur explique, il ne juge pas ;
\item des \textbf{connecteurs de cause} : \emph{parce que, car, puisque, en raison de, grâce à, du fait de} ;
\item des \textbf{connecteurs de conséquence} : \emph{c'est pourquoi, ainsi, donc, par conséquent, de sorte que, si bien que} ;
\item l'emploi du \textbf{présent de vérité générale} (ou présent intemporel) : \og l'urbanisation entraîne une forte demande de logements \fg{} ;
\item un vocabulaire précis, parfois technique, et des \textbf{exemples et chiffres} qui illustrent l'explication.
\end{enumerate}
\end{propriete}

\begin{exemplebox}[Repérer les connecteurs]
Phrase explicative : \og Les jeunes quittent les villages \textbf{parce que} les villes offrent plus d'emplois ; \textbf{c'est pourquoi} la population de Douala augmente rapidement. \fg{}

Ici, \emph{parce que} introduit la \textbf{cause} (l'emploi urbain) et \emph{c'est pourquoi} introduit la \textbf{conséquence} (la croissance démographique de Douala). Repérer ces connecteurs, c'est déjà repérer l'ossature logique du texte.
\end{exemplebox}

\begin{attention}[Erreur fréquente : la paraphrase]
Recopier des phrases entières du texte en changeant seulement quelques mots, ce n'est \textbf{pas} un résumé : c'est de la \cle{paraphrase}. Le correcteur du Baccalauréat le sanctionne. Un vrai résumé \textbf{reformule} les idées avec tes propres mots, en phrases plus courtes et plus denses. On y reviendra en détail dans les sections suivantes.
\end{attention}

\courssub{Repérer le thème, la thèse et la problématique}

Avant de réduire un texte, il faut comprendre \emph{de quoi il parle} et \emph{ce qu'il veut démontrer}. Trois repères guident cette première lecture :

\begin{definition}[Thème, thèse, problématique]
\begin{itemize}
\item Le \textbf{thème} est le sujet général du texte : il tient souvent en quelques mots (\og l'urbanisation du Cameroun \fg{}).
\item La \textbf{thèse} est l'idée principale que l'auteur explique ou défend (\og l'urbanisation rapide crée de graves problèmes de logement et d'emploi \fg{}).
\item La \textbf{problématique} est la question centrale à laquelle le texte répond (\og quelles sont les causes et les conséquences de l'urbanisation rapide ? \fg{}).
\end{itemize}
\end{definition}

\begin{methode}[Lire un texte avant de le résumer]
\begin{enumerate}
\item \textbf{Lire le texte une première fois} en entier, sans rien noter, pour saisir le sens général et identifier le thème.
\item \textbf{Lire une deuxième fois}, crayon en main : repérer la thèse (souvent dans l'introduction ou la conclusion) et la problématique.
\item \textbf{Souligner les idées principales} : une idée par paragraphe en général ; les connecteurs logiques (\emph{parce que, c'est pourquoi, ainsi, d'abord, ensuite}) sont des balises précieuses.
\item \textbf{Raturer ou mettre entre parenthèses} tout ce qui n'est pas essentiel : les exemples, les répétitions, les comparaisons, les citations, les chiffres secondaires et les détails anecdotiques.
\item \textbf{Vérifier} que ce qui reste souligné forme, à lui seul, un ensemble cohérent qui résume déjà le texte.
\end{enumerate}
\end{methode}

\courssub{Lecture analytique d'un texte explicatif : l'urbanisation de Douala et Yaoundé}

Mettons la méthode en application sur un texte explicatif. Lis d'abord le document en entier.

\begin{exemplebox}[Texte à étudier (environ 200 mots)]
\og Le Cameroun connaît depuis un demi-siècle une urbanisation spectaculaire. Alors qu'en 1976 à peine 20\,\% des Camerounais vivaient en ville, ce taux dépassait 35\,\% en 1987, approchait 50\,\% en 2005 et franchit aujourd'hui la barre des 55\,\%. Douala et Yaoundé, les deux métropoles du pays, concentrent à elles seules plusieurs millions d'habitants.

Comment expliquer un tel bouleversement ? D'abord, l'exode rural : les jeunes quittent les villages parce que la ville promet l'emploi, les études et les services de santé. Ensuite, la croissance démographique : les citadins ont des enfants qui, nés en ville, y restent. Ainsi, la ville gonfle de l'intérieur en même temps qu'elle attire.

Cette urbanisation rapide n'est pas sans conséquences. Parce que la construction de logements ne suit pas le rythme de l'arrivée des populations, des quartiers précaires s'étendent, comme à Ndogpassi à Douala ou à Nsam à Yaoundé. C'est pourquoi les pouvoirs publics doivent anticiper : planifier l'habitat, développer les transports et créer des emplois. Sans cette planification, la ville, au lieu d'être un moteur de développement, deviendrait un facteur de crise sociale. \fg{}
\end{exemplebox}

\textbf{Première lecture : le thème.} Le texte parle de l'\emph{urbanisation rapide du Cameroun}, et plus précisément de Douala et Yaoundé.

\textbf{Deuxième lecture : thèse et problématique.} La thèse apparaît dans le dernier paragraphe : l'urbanisation rapide, si elle n'est pas planifiée, devient un danger social ; il faut donc l'anticiper. La problématique peut se formuler ainsi : \emph{quelles sont les causes et les conséquences de l'urbanisation rapide du Cameroun, et comment y répondre ?}

\textbf{Repérage des idées principales.} Paragraphe par paragraphe, on souligne :
\begin{enumerate}
\item le taux d'urbanisation est passé d'environ 20\,\% en 1976 à plus de 55\,\% aujourd'hui (idée : \emph{une urbanisation spectaculaire}) ;
\item causes : l'exode rural et la croissance démographique (idée : \emph{deux moteurs de la croissance urbaine}) ;
\item conséquences : l'extension des quartiers précaires et la nécessité de planifier (idée : \emph{des problèmes qui exigent une anticipation}).
\end{enumerate}

\textbf{Ce qu'on rature.} Les exemples de quartiers (Ndogpassi, Nsam), la formulation imagée \og la ville gonfle de l'intérieur \fg{}, les détails (\og les études et les services de santé \fg{} peuvent être condensés en \og de meilleures conditions de vie \fg{}). Attention toutefois : les chiffres-clés (20\,\% en 1976, plus de 55\,\% aujourd'hui) sont ici \emph{essentiels} car ils prouvent l'ampleur du phénomène ; on peut les condenser en une seule mention.

\begin{popfigure}
\begin{center}
\begin{tikzpicture}
\begin{axis}[
ybar,
width=0.85\linewidth,
height=6.5cm,
bar width=0.9cm,
ymin=0,ymax=65,
ylabel={Taux d'urbanisation (en \%)},
xlabel={Année},
symbolic x coords={1976,1987,2005,2020},
xtick=data,
nodes near coords,
nodes near coords align={vertical},
every node near coord/.append style={font=\small\bfseries},
ymajorgrids=true,
grid style={popline!40},
axis line style={popdark},
tick label style={font=\small},
label style={font=\small},
]
\addplot[fill=popblue!70,draw=popblue] coordinates {(1976,20) (1987,35) (2005,50) (2020,56)};
\end{axis}
\end{tikzpicture}
\end{center}
\legende{Évolution du taux d'urbanisation du Cameroun (valeurs approximatives, en \% de la population vivant en ville) : d'environ 20\,\% en 1976 à plus de 55\,\% en 2020.}
\end{popfigure}

Le diagramme en barres confirme visuellement ce que le texte explique : la progression est \emph{continue et rapide}. Dans un résumé, on ne décrira pas le graphique barre par barre ; on en retiendra l'idée unique : \og le taux d'urbanisation a presque triplé en un demi-siècle \fg{}.

\begin{exoresolu}[Première approche de la contraction]
\textbf{Énoncé.} Voici un passage extrait du texte étudié : \og D'abord, l'exode rural : les jeunes quittent les villages parce que la ville promet l'emploi, les études et les services de santé. Ensuite, la croissance démographique : les citadins ont des enfants qui, nés en ville, y restent. Ainsi, la ville gonfle de l'intérieur en même temps qu'elle attire. \fg{} (43 mots). Réduis ce passage à environ 11 mots (le quart), sans recopier de phrase entière.
\tcblower
\textbf{Solution détaillée.}
\begin{enumerate}
\item Idées essentielles : deux causes (exode rural, croissance démographique) ; une conséquence (la ville grandit).
\item À éliminer : les détails (\og l'emploi, les études et les services de santé \fg{}), l'image (\og gonfle de l'intérieur \fg{}).
\item Reformulation : \og L'exode rural et la natalité urbaine font croître les villes. \fg{} (10 mots).
\end{enumerate}
Cette phrase conserve les deux causes et la conséquence, avec des mots personnels, sans aucune phrase recopiée : c'est une contraction, pas une paraphrase.
\end{exoresolu}

\begin{savaistu}[La contraction au Baccalauréat technique]
Au Probatoire et au Baccalauréat technique, la contraction de texte est notée selon trois critères principaux : la \textbf{fidélité au texte} (les idées essentielles sont toutes présentes, sans ajout ni opinion), la \textbf{qualité de la langue} (orthographe, grammaire, ponctuation, phrases bien construites) et le \textbf{respect du nombre de mots imposé} (en général le quart, avec une tolérance de $\pm 10\,\%$). Un résumé brillant mais hors limite de mots perd des points : compte toujours tes mots et indique le total à la fin de ta copie !
\end{savaistu}

\begin{exercice}[À toi de jouer]
Reprends le premier paragraphe du texte sur l'urbanisation (\og Le Cameroun connaît depuis un demi-siècle\ldots plusieurs millions d'habitants. \fg{}). 1) Souligne les idées principales. 2) Rature les détails et exemples. 3) Reformule le paragraphe en une phrase d'environ 15 mots. 4) Indique le nombre de mots de ta phrase.
\end{exercice}

\begin{corrige}[Exercice : réduire le premier paragraphe]
Éléments de correction : idées essentielles = urbanisation spectaculaire ; passage d'environ 20\,\% (1976) à plus de 55\,\% (aujourd'hui) ; poids de Douala et Yaoundé. À raturer : la liste intermédiaire des pourcentages (1987, 2005), qui peut être condensée. Reformulation possible (15 mots) : \og L'urbanisation camerounaise, menée par Douala et Yaoundé, est passée de 20\,\% en 1976 à plus de 55\,\% aujourd'hui. \fg{} Toute formulation équivalente, fidèle et sans phrase recopiée, est acceptée.
\end{corrige}

\begin{aretenir}[L'essentiel de la section]
\begin{itemize}
\item La \textbf{contraction de texte} réduit un texte à une longueur imposée (le quart) en conservant les idées essentielles, sans opinion personnelle.
\item Le \textbf{texte explicatif} informe et fait comprendre : causes (\emph{parce que}), conséquences (\emph{c'est pourquoi, ainsi}), présent de vérité générale.
\item Avant de résumer : \textbf{deux lectures}, repérage du \textbf{thème}, de la \textbf{thèse} et de la \textbf{problématique}, puis on \textbf{souligne} les idées principales et on \textbf{élimine} exemples, répétitions et détails.
\item Piège majeur : la \textbf{paraphrase} (recopier en changeant quelques mots) n'est pas un résumé.
\end{itemize}
\end{aretenir}

\coursec{La structure argumentative et le plan du texte}

Dans la section précédente, nous avons découvert le texte explicatif et compris pourquoi la contraction de texte est un exercice fondamental : résumer, c'est restituer fidèlement les idées essentielles d'un texte en un quart de sa longueur, sans rien ajouter de soi-même. Mais comment savoir quelles idées sont « essentielles » ? La réponse tient en un mot : la \cle{structure}. Un texte argumentatif ou explicatif n'est pas un sac de phrases jetées en désordre : c'est un édifice construit selon un plan. Celui qui voit le plan voit l'essentiel. Cette section vous apprend donc à démonter l'édifice pour en dégager l'ossature : thèse, arguments, exemples, progression, articulateurs logiques, et enfin le plan du texte, qui deviendra le plan de votre résumé.

\courssub{Thèse, argument, exemple : les trois étages du raisonnement}

Commençons par une situation simple. Un élève déclare : « Le téléphone portable est utile à l'école. » Un camarade lui demande : « Pourquoi ? » Il répond : « Parce qu'il permet de faire des recherches documentaires rapidement. » Le camarade insiste : « Par exemple ? » Il ajoute : « Hier, en cours de géographie, j'ai vérifié en deux minutes la population de la région du Centre. » En trois répliques, cet élève a produit, sans le savoir, les trois étages de tout raisonnement argumentatif.

\begin{definition}[Thèse, argument, exemple]
\begin{itemize}
\item La \cle{thèse} est l'\textbf{idée défendue} par l'auteur, le point de vue qu'il cherche à faire partager au lecteur. Elle répond à la question : « Que veut-il démontrer ? »
\item L'\cle{argument} est la \textbf{raison} avancée pour soutenir la thèse. Il répond à la question : « Pourquoi faut-il l'admettre ? »
\item L'\cle{exemple} est l'\textbf{illustration concrète} qui rend l'argument visible et crédible. Il répond à la question : « Peut-on le voir en acte ? »
\end{itemize}
\end{definition}

Ces trois notions forment une hiérarchie : la thèse domine, les arguments la soutiennent, les exemples soutiennent les arguments. C'est exactement la structure d'un arbre : le tronc porte les branches, les branches portent les feuilles.

\begin{popfigure}
\begin{tikzpicture}[
  these/.style={rectangle, rounded corners, draw=popink, fill=popblueL, thick, text width=5.2cm, align=center, font=\small\bfseries},
  arg/.style={rectangle, rounded corners, draw=popdark, fill=poporangeL, thick, text width=3.4cm, align=center, font=\small},
  ex/.style={rectangle, rounded corners, draw=popmuted, fill=popgreenL, text width=2.9cm, align=center, font=\footnotesize},
  lien/.style={thick, popdark}
]
\node[these, align=center] (T) at (0,0){THÈSE\\ Le téléphone portable est utile à l'école};
\node[arg, align=center] (A1) at (-5.2,-2.6){Argument 1\\ recherches documentaires rapides};
\node[arg, align=center] (A2) at (0,-2.6){Argument 2\\ calculs et applications éducatives};
\node[arg, align=center] (A3) at (5.2,-2.6){Argument 3\\ communication avec les parents};
\node[ex] (E1) at (-5.2,-5.2) {Ex. : vérifier la population d'une région en géographie};
\node[ex] (E2) at (0,-5.2) {Ex. : utiliser la calculatrice scientifique en physique};
\node[ex] (E3) at (5.2,-5.2) {Ex. : prévenir en cas de retard du transport};
\draw[lien] (T) -- (A1);
\draw[lien] (T) -- (A2);
\draw[lien] (T) -- (A3);
\draw[lien] (A1) -- (E1);
\draw[lien] (A2) -- (E2);
\draw[lien] (A3) -- (E3);
\end{tikzpicture}
\legende{Arbre de la structure argumentative : la thèse au sommet (tronc), trois arguments en branches, un exemple en feuille sous chaque argument. Dans un résumé, on conserve le tronc et les branches ; on coupe les feuilles.}
\end{popfigure}

\begin{attention}[Argument ou exemple : ne pas confondre]
L'erreur la plus fréquente à l'examen est de confondre un \textbf{exemple} (à supprimer dans le résumé) avec un \textbf{argument} (à conserver). Test pratique : demandez-vous si la phrase donne une \emph{raison générale} (argument : « le portable facilite la recherche ») ou un \emph{cas particulier daté et situé} (exemple : « hier, en cours de géographie, j'ai vérifié... »). Les marques « par exemple », « ainsi », « c'est le cas de », « comme » introduisent presque toujours des exemples. Un résumé qui garde les exemples dépasse la longueur autorisée et trahit la hiérarchie des idées.
\end{attention}

\courssub{La notion de progression : les trois types}

\begin{definition}[Progression]
La \cle{progression} est l'\textbf{organisation ordonnée des idées} dans un texte, du début vers la conclusion. Elle désigne le chemin que l'auteur fait parcourir au lecteur pour l'amener de la question posée à la réponse défendue.
\end{definition}

L'intuition est celle d'un voyage : pour aller de Yaoundé à Douala, on peut prendre la route directe, faire un détour pour visiter une ville, ou s'arrêter à chaque étape pour observer un paysage. De même, un auteur peut avancer droit vers sa conclusion, passer par une étape d'opposition, ou explorer son sujet thème par thème. On distingue trois types de progression.

\begin{definition}[Les trois types de progression]
\begin{itemize}
\item La \cle{progression linéaire} (ou déductive) : l'auteur annonce sa thèse dès le début, puis enchaîne les arguments qui la justifient, et conclut. Schéma : \textbf{thèse $\rightarrow$ arguments $\rightarrow$ conclusion}. C'est la progression du texte explicatif et de l'éditorial engagé.
\item La \cle{progression dialectique} : l'auteur examine d'abord une position (thèse), puis la position opposée (antithèse), puis dépasse l'opposition dans une synthèse. Schéma : \textbf{thèse $\rightarrow$ antithèse $\rightarrow$ synthèse}. C'est la progression du débat et de la dissertation.
\item La \cle{progression thématique} : l'auteur explore un sujet par aspects successifs, sans opposition frontale : causes, effets, solutions ; ou passé, présent, avenir. Chaque partie traite un \textbf{thème} ou un angle du sujet.
\end{itemize}
\end{definition}

\begin{popfigure}
\begin{tikzpicture}[
  etape/.style={rectangle, rounded corners, draw=popdark, thick, text width=2.5cm, align=center, font=\footnotesize, minimum height=0.9cm},
  titre/.style={font=\small\bfseries, popink},
  fl/.style={-{Stealth[length=2.5mm]}, very thick, popink}
]
% Progression linéaire
\node[titre] at (0,0) {Progression linéaire};
\node[etape, fill=popblueL] (L1) at (-3.4,-1) {Thèse annoncée};
\node[etape, fill=poporangeL] (L2) at (0,-1) {Arguments 1, 2, 3};
\node[etape, fill=popgreenL] (L3) at (3.4,-1) {Conclusion};
\draw[fl] (L1) -- (L2);
\draw[fl] (L2) -- (L3);
% Progression dialectique
\node[titre] at (0,-2.8) {Progression dialectique};
\node[etape, fill=popblueL] (D1) at (-3.4,-3.8) {Thèse (pour)};
\node[etape, fill=poppinkL] (D2) at (0,-3.8) {Antithèse (contre)};
\node[etape, fill=popgreenL] (D3) at (3.4,-3.8) {Synthèse (position de l'auteur)};
\draw[fl] (D1) -- (D2);
\draw[fl] (D2) -- (D3);
% Progression thématique
\node[titre] at (0,-5.6) {Progression thématique};
\node[etape, fill=popgoldL, align=center] (T1) at (-3.4,-6.6){Aspect 1\\ (causes)};
\node[etape, fill=popgoldL, align=center] (T2) at (0,-6.6){Aspect 2\\ (effets)};
\node[etape, fill=popgoldL, align=center] (T3) at (3.4,-6.6){Aspect 3\\ (solutions)};
\draw[fl] (T1) -- (T2);
\draw[fl] (T2) -- (T3);
\end{tikzpicture}
\legende{Les trois types de progression : 1) linéaire, la thèse est donnée d'emblée puis justifiée ; 2) dialectique, deux positions s'affrontent avant la synthèse ; 3) thématique, le sujet est parcouru par aspects successifs.}
\end{popfigure}

\begin{propriete}[Où se trouve la thèse de l'auteur ?]
Dans une \textbf{progression linéaire}, la thèse de l'auteur se trouve dans l'introduction et est reprise dans la conclusion. En revanche, dans une \textbf{progression dialectique}, la thèse réellement défendue par l'auteur se trouve le plus souvent dans la \textbf{synthèse finale}, et non dans la première partie : la première thèse exposée est une position examinée, parfois même combattue. Conséquence pratique pour le résumé : ne jamais attribuer à l'auteur la première opinion venue ; lire le texte jusqu'au bout avant de formuler sa position.
\end{propriete}

\courssub{Les articulateurs logiques : les panneaux indicateurs du plan}

Comment le lecteur repère-t-il la progression ? Grâce aux \cle{articulateurs logiques} (ou connecteurs), qui fonctionnent comme des panneaux indicateurs le long de la route. Les connaître, c'est pouvoir lire le plan directement dans le texte.

\begin{center}
\begin{tabularx}{\linewidth}{|l|X|X|}\hline
\rowcolor{popblueL}
\textbf{Fonction} & \textbf{Articulateurs fréquents} & \textbf{Ce qu'ils signalent} \\ \hline
Énumération & d'abord, ensuite, puis, enfin, en premier lieu & le découpage en parties ou en arguments \\ \hline
Explication & en effet, car, parce que, c'est pourquoi & une justification, une cause \\ \hline
Concession & certes, il est vrai que, sans doute & une idée admise provisoirement \\ \hline
Opposition & mais, cependant, pourtant, néanmoins, en revanche & un retournement, souvent vers la vraie position de l'auteur \\ \hline
Illustration & par exemple, ainsi, c'est le cas de & un exemple (à écarter du résumé) \\ \hline
Conclusion & donc, ainsi, finalement, en définitive, c'est pourquoi & la synthèse ou la conclusion \\ \hline
\end{tabularx}
\end{center}

Observez le couple \textbf{certes... mais...} : il est la signature de la progression dialectique en miniature. « \emph{Certes} les jeunes manquent d'expérience, \emph{mais} ils apportent l'énergie et l'innovation. » L'auteur concède d'abord (antithèse), puis retourne (thèse réelle). De même, « \emph{en effet} » confirme ce qui précède, tandis que « \emph{cependant} » annonce ce qui le contredit. Surligner ces mots dans le texte, c'est déjà dessiner son plan.

\courssub{Construire le plan du texte : méthode pas à pas}

\begin{methode}[Dégager le plan d'un texte en quatre étapes]
\begin{enumerate}
\item \textbf{Lire le texte entier deux fois}, sans écrire, pour saisir le sujet et la position globale de l'auteur (attention aux progressions dialectiques : lire jusqu'au bout).
\item \textbf{Repérer et surligner les connecteurs} (d'abord, ensuite, certes, mais, enfin...) : ils marquent les frontières entre les parties.
\item \textbf{Regrouper les paragraphes par idée} : plusieurs paragraphes qui développent la même idée forment une seule grande partie. On obtient en général 3 ou 4 parties.
\item \textbf{Formuler un titre pour chaque partie}, sous forme d'un groupe nominal bref (ex. : « Les difficultés des jeunes », « Leur rôle dans le développement »). L'ensemble des titres, dans l'ordre, constitue le plan du texte.
\end{enumerate}
\end{methode}

\begin{exoresolu}[Le plan d'un éditorial du \emph{Cameroon Tribune}]
Énoncé. Voici, résumé fidèlement, le contenu d'un éditorial intitulé « Les jeunes, moteur du Cameroun de demain » (texte de 600 mots environ) :

\emph{Paragraphes 1-2 : le constat.} Les jeunes de moins de 35 ans représentent la majorité de la population camerounaise, pourtant ils restent largement absents des postes de décision et beaucoup sont frappés par le chômage. \emph{Paragraphes 3-4 : certes...} Certes, on objecte que les jeunes manquent d'expérience et de formation adaptée ; les programmes scolaires ne répondent pas toujours aux besoins des entreprises. \emph{Paragraphes 5-7 : mais...} Mais les jeunes possèdent des atouts décisifs : la maîtrise du numérique, l'esprit d'initiative, l'énergie. L'agropastoral, l'artisanat et les start-up montrent déjà, par exemple à Bafoussam ou à Douala, ce que des jeunes formés peuvent accomplir. \emph{Paragraphe 8 : conclusion.} En définitive, investir dans la formation et l'insertion des jeunes, c'est investir dans le développement du pays tout entier.

Dégagez la progression et le plan de ce texte.

\tcblower
\textbf{Solution détaillée.}

\emph{Étape 1 — Position de l'auteur.} Le titre et la conclusion concordent : l'auteur défend l'idée que les jeunes sont l'avenir du pays et qu'il faut investir en eux.

\emph{Étape 2 — Connecteurs repérés.} « Certes » (début de concession), « Mais » (retournement), « par exemple » (illustrations à écarter), « En définitive » (conclusion).

\emph{Étape 3 — Progression.} Le schéma « certes... mais... » révèle une \textbf{progression dialectique} : constat/thèse adverse $\rightarrow$ objection $\rightarrow$ réfutation $\rightarrow$ synthèse. La vraie thèse de l'auteur est donc dans la dernière partie, conformément à la propriété vue plus haut.

\emph{Étape 4 — Plan (4 parties).}
\begin{enumerate}
\item Un constat paradoxal : une jeunesse majoritaire mais marginalisée.
\item L'objection : le manque d'expérience et de formation des jeunes.
\item La réfutation : les atouts et les réussites déjà visibles de la jeunesse.
\item Conclusion : investir dans les jeunes, condition du développement national.
\end{enumerate}
Chaque titre est un groupe nominal bref qui résume l'idée de la partie, sans reprendre les exemples (Bafoussam, Douala), qui seront supprimés du résumé.
\end{exoresolu}

\courssub{Du plan du texte au plan du résumé}

Voici maintenant le point décisif pour l'examen : \textbf{le plan du texte devient le plan du résumé}. La règle est simple : chaque grande partie du texte donne \textbf{un paragraphe} du résumé. Si le texte compte quatre parties, le résumé comptera quatre paragraphes courts, reliés entre eux par des connecteurs qui reproduisent la progression de l'auteur (d'abord, certes, mais, enfin).

\begin{exemplebox}[Un exemple chiffré pour bien mesurer]
Un texte de \textbf{600 mots} structuré en \textbf{4 parties} donnera un résumé d'environ \textbf{150 mots} (le quart de la longueur), rédigé en \textbf{4 paragraphes courts} de 30 à 45 mots chacun. Répartition indicative : 35 mots pour le constat, 30 mots pour l'objection, 50 mots pour la réfutation (c'est la partie la plus chargée en arguments), 35 mots pour la conclusion. Ce découpage garantit deux choses : aucune idée essentielle n'est oubliée, et la proportion des parties respecte celle du texte original.
\end{exemplebox}

\begin{aretenir}[L'essentiel de la section]
\begin{itemize}
\item Toute démonstration repose sur trois étages : la \cle{thèse} (idée défendue), les \cle{arguments} (raisons), les \cle{exemples} (illustrations à écarter du résumé).
\item La \cle{progression} est le chemin des idées : \textbf{linéaire} (thèse $\rightarrow$ arguments $\rightarrow$ conclusion), \textbf{dialectique} (thèse $\rightarrow$ antithèse $\rightarrow$ synthèse, la vraie position étant dans la synthèse), \textbf{thématique} (par aspects).
\item Les \cle{articulateurs logiques} (d'abord, en effet, certes, mais, cependant, enfin) sont les panneaux qui balisent le plan.
\item Pour construire le plan : lire deux fois, surligner les connecteurs, regrouper les paragraphes par idée, titrer chaque partie.
\item Chaque partie du plan = un paragraphe du résumé ; la longueur du résumé = le quart du texte.
\end{itemize}
\end{aretenir}

\begin{exercice}[À vous de jouer]
Lisez un éditorial ou une lettre ouverte parue dans un journal camerounais (\emph{Cameroon Tribune}, \emph{Mutations}...) sur un sujet de société (l'éducation, l'environnement, le sport). 1) Soulignez tous les articulateurs logiques. 2) Identifiez le type de progression. 3) Rédigez le plan du texte en trois ou quatre titres. 4) Indiquez, pour chaque partie, si elle contient des exemples à écarter du futur résumé.
\end{exercice}

\begin{corrige}[Exercice — éléments de correction]
La correction dépend du texte choisi, mais la démarche attendue est constante : 1) les articulateurs doivent être cités exactement (d'abord, certes, mais, en effet, enfin...) ; 2) la présence du couple « certes... mais... » ou d'une opposition franche signe une progression dialectique, l'enchaînement « d'abord / ensuite / enfin » sous une thèse annoncée signe une progression linéaire ; 3) chaque titre doit être un groupe nominal bref couvrant une seule idée ; 4) les passages introduits par « par exemple », « ainsi », ou contenant des noms propres, des dates et des cas précis sont des exemples : ils seront supprimés au moment de la rédaction du résumé, que nous aborderons dans les sections suivantes.
\end{corrige}

Nous savons désormais lire l'ossature d'un texte. Reste à comprendre ce qui tient cet édifice debout de l'intérieur : la cohésion et la cohérence, c'est-à-dire les liens invisibles entre les phrases, que la prochaine section examinera avec la structure de la phrase simple et de la phrase composée.

\coursec{Cohésion, cohérence et structure de la phrase}

Dans la section précédente, nous avons appris à dégager la \cle{structure argumentative} d'un texte et à en construire le \cle{plan}. Mais un bon plan ne suffit pas : un résumé peut contenir toutes les idées essentielles et pourtant être illisible, parce que ses phrases se succèdent sans lien, comme des pierres posées les unes à côté des autres sans ciment. Cette section vous donne le ciment : la \cle{cohésion} (les liens matériels entre les phrases), la \cle{cohérence} (la logique globale du sens) et la maîtrise de la \cle{phrase simple} et de la \cle{phrase composée}, outils indispensables pour condenser ou alléger vos formulations.

\courssub{La cohérence : l'unité de sens du texte}

\textbf{Intuition.}\par
Lisez ces trois phrases : « Le Cameroun exporte du bois. Mon frère aime le foot. Il pleut à Douala. » Chaque phrase est correcte, pourtant l'ensemble ne fait pas un texte : rien ne relie ces idées, aucun sens global ne se dégage. Un texte n'est pas une somme de phrases, c'est un \emph{tissu} d'idées qui avancent ensemble vers un même but.

\begin{definition}[La cohérence]
La \cle{cohérence} est la qualité d'un texte dont les idées s'enchaînent logiquement et forment un sens global. Elle repose sur trois piliers :
\begin{itemize}
\item l'\cle{unité de sens} : toutes les phrases parlent du même sujet et servent la même thèse ;
\item l'\cle{enchaînement logique} : chaque idée découle de la précédente (cause, conséquence, opposition, illustration) ;
\item le \cle{fil conducteur sémantique} : un même thème (souvent porté par un champ lexical) traverse le texte du début à la fin.
\end{itemize}
\end{definition}

Dans un résumé, la cohérence est doublement exigée : vous devez \emph{respecter} la logique du texte de départ (ne jamais inverser une cause et une conséquence) et \emph{reproduire} cette logique dans vos propres phrases, sans l'auteur pour vous guider.

\begin{attention}[Piège fréquent]
Un résumé peut être fidèle idée par idée et pourtant incohérent si l'ordre des idées est bouleversé. L'auteur écrit : « Les pluies diminuent, \emph{par conséquent} les récoltes baissent. » Si vous résumez : « Les récoltes baissent, car les pluies diminuent », le sens est conservé ; mais si vous écrivez « Les récoltes baissent. Les pluies diminuent. » sans lien, le rapport de cause à effet disparaît : la cohérence est perdue.
\end{attention}

\courssub{La cohésion : les liens matériels du texte}

\textbf{Intuition.}\par
Si la cohérence est l'architecture invisible du texte, la cohésion en est la maçonnerie visible : ce sont les petits mots et procédés qui cousent les phrases entre elles et évitent les répétitions lourdes.

\begin{definition}[La cohésion]
La \cle{cohésion} est l'ensemble des procédés grammaticaux et lexicaux qui relient les éléments du texte entre eux : reprises pronominales, substituts lexicaux, champs lexicaux et connecteurs logiques.
\end{definition}

\textbf{1. Les reprises pronominales.} Un pronom reprend un nom déjà mentionné (son \emph{antécédent}) et évite de le répéter : « La forêt du bassin du Congo est menacée. \emph{Elle} couvre pourtant des millions d'hectares. » Les pronoms personnels (\emph{il, elle, ils, elles}), démonstratifs (\emph{celle-ci, celui-là, cela}) et relatifs (\emph{qui, que, dont}) sont les coutures les plus fréquentes du texte.

\textbf{2. Les substituts lexicaux.} Au lieu de répéter « le Cameroun », on peut écrire « \emph{le pays} », « \emph{cette nation} », « \emph{l'État camerounais} ». Le substitut reprend la même réalité sous un autre nom : c'est un outil précieux dans le résumé, où l'on condense en variant les désignations. Veillez cependant à ce que le substitut soit toujours clair : le lecteur ne doit jamais se demander de quoi l'on parle.

\textbf{3. Les champs lexicaux.} Un \cle{champ lexical} est l'ensemble des mots qui se rapportent à une même notion : \emph{forêt, arbre, bois, essence, défrichement, canopée} appartiennent au champ lexical de la forêt. Sa présence continue assure le fil conducteur sémantique et signale au lecteur que l'on reste sur le même sujet.

\begin{center}
\begin{tabularx}{\linewidth}{|l|X|X|}
\hline
\rowcolor{popblueL} \textbf{Procédé de cohésion} & \textbf{Exemple} & \textbf{Effet} \\
\hline
Reprise pronominale & « La forêt est menacée. \emph{Elle} abrite des espèces rares. » & Évite la répétition, relie les phrases \\
\hline
Substitut lexical & « Le Cameroun... \emph{ce pays forestier}... » & Varie la désignation, condense \\
\hline
Champ lexical & \emph{forêt, bois, défrichement, essence} & Assure l'unité thématique \\
\hline
Connecteur logique & « \emph{Cependant}, la déforestation s'accélère. » & Explicite le rapport logique \\
\hline
\end{tabularx}
\end{center}

\courssub{Les connecteurs logiques et leur fonction}

Les \cle{connecteurs logiques} sont les panneaux de signalisation du texte : ils indiquent au lecteur dans quelle direction va la pensée. Un résumé sans connecteurs est une route sans panneaux.

\begin{popfigure}
\begin{center}
\begin{tikzpicture}[font=\small]
\node[draw=popblue, fill=popblueL, rounded corners, minimum width=2.6cm, minimum height=1.1cm, align=center] (add) at (0,2.6) {\textbf{Addition}\\ et, de plus,\\ aussi, en outre};
\node[draw=poporange, fill=poporangeL, rounded corners, minimum width=2.6cm, minimum height=1.1cm, align=center] (opp) at (5.2,2.6) {\textbf{Opposition}\\ mais, cependant,\\ pourtant, alors que};
\node[draw=popgreen, fill=popgreenL, rounded corners, minimum width=2.6cm, minimum height=1.1cm, align=center] (cause) at (0,0) {\textbf{Cause}\\ car, parce que,\\ puisque, en effet};
\node[draw=poppurple, fill=poppurpleL, rounded corners, minimum width=2.6cm, minimum height=1.1cm, align=center] (cons) at (5.2,0) {\textbf{Conséquence}\\ donc, par conséquent,\\ c'est pourquoi, ainsi};
\node[draw=popgold, fill=popgoldL, rounded corners, minimum width=2.6cm, minimum height=1.1cm, align=center] (conc) at (2.6,-2.4) {\textbf{Concession}\\ certes... mais,\\ bien que, même si};
\end{tikzpicture}
\end{center}
\legende{Les connecteurs logiques classés par fonction : chaque famille exprime un rapport précis entre les idées.}
\end{popfigure}

\begin{attention}[Erreur fréquente : « mais » pour « donc »]
Un connecteur mal choisi \emph{fausse le sens} du texte de départ. Écrire « Il pleut, \emph{mais} je prends mon parapluie » suggère une opposition absurde ; le texte original disait peut-être « Il pleut, \emph{donc} je prends mon parapluie » (conséquence). Au résumé, avant d'écrire un connecteur, demandez-vous toujours : quel est le rapport logique réel entre ces deux idées dans le texte de l'auteur ?
\end{attention}

\courssub{La phrase simple et la phrase composée}

\textbf{Intuition.}\par
Comptez les verbes conjugués : c'est le moyen infaillible de reconnaître le type de phrase.

\begin{definition}[Phrase simple et phrase composée]
\begin{itemize}
\item La \cle{phrase simple} contient \textbf{un seul verbe conjugué} : elle forme une seule proposition. Exemple : « Le Cameroun \emph{exporte} du bois. »
\item La \cle{phrase composée} contient \textbf{plusieurs verbes conjugués}, donc plusieurs propositions, reliées par \emph{juxtaposition} ou par \emph{coordination}.
\end{itemize}
\end{definition}

\textbf{La juxtaposition} relie deux propositions par une simple ponctuation : la virgule ou le point-virgule. « Les pluies diminuent, les récoltes baissent. »

\textbf{La coordination} relie deux propositions par un \cle{conjonction de coordination}. Retenez la liste par cœur : \emph{\cle{mais, ou, et, donc, or, ni, car}}. « Les pluies diminuent, \emph{donc} les récoltes baissent. »

\begin{popfigure}
\begin{center}
\begin{tikzpicture}[font=\small]
% Phrase simple
\node[draw=popblue, fill=popblueL, rounded corners, minimum width=6.2cm, minimum height=1.2cm, align=center] (ps) at (0,2.2) {\textbf{Phrase simple} : une proposition\\ « Le Cameroun \emph{exporte} du bois. »};
% Phrase composée
\node[draw=poporange, fill=poporangeL, rounded corners, minimum width=2.9cm, minimum height=1.1cm, align=center] (p1) at (-1.8,0) {Proposition 1\\ « Les pluies \emph{diminuent} »};
\node[draw=poporange, fill=poporangeL, rounded corners, minimum width=2.9cm, minimum height=1.1cm, align=center] (p2) at (1.8,0) {Proposition 2\\ « les récoltes \emph{baissent} »};
\draw[-{Stealth}, thick, popdark] (p1) -- (p2) node[midway, above, popdark]{\emph{donc} (ou ;)};
\node[popdark, align=center] at (0,-1.3) {\textbf{Phrase composée} : deux propositions reliées\\ par coordination (conjonction) ou juxtaposition (ponctuation)};
\end{tikzpicture}
\end{center}
\legende{Comparaison : la phrase simple est un bloc unique ; la phrase composée assemble plusieurs blocs (propositions) reliés par une conjonction de coordination ou une ponctuation.}
\end{popfigure}

\courssub{Transformer les phrases pour condenser ou alléger}

C'est ici que la grammaire devient un outil de résumé. Deux mouvements sont possibles :

\textbf{1. Condenser : passer de plusieurs phrases simples à une phrase composée ou resserrée.} On fusionne les phrases en exploitant les reprises et les expansions du nom (apposition, relative).

\begin{exemplebox}[La condensation par apposition]
Texte de départ : « Le Cameroun est un pays forestier. Il perd chaque année des milliers d'hectares. »\\
Résumé condensé : « Le Cameroun, pays forestier, perd chaque année des milliers d'hectares. »\\
Deux phrases simples (deux verbes : \emph{est}, \emph{perd}) deviennent une seule phrase simple grâce à l'\emph{apposition} « pays forestier » : on gagne des mots sans perdre aucune idée.
\end{exemplebox}

\textbf{2. Alléger : découper une phrase composée trop longue en phrases simples.} « Les pluies diminuent et les récoltes baissent, mais les paysans ne reçoivent aucune aide, donc la pauvreté s'aggrave. » → « Les pluies diminuent. Les récoltes baissent. Or les paysans ne reçoivent aucune aide. La pauvreté s'aggrave donc. » Le découpage rend le propos plus clair, à condition de conserver les connecteurs qui portent la logique.

\begin{methode}[Vérifier la cohésion de son résumé]
\begin{enumerate}
\item Relisez votre résumé \textbf{à voix haute} : l'oreille détecte les ruptures de lien que l'œil ne voit pas.
\item Pour chaque phrase, demandez-vous : par quoi est-elle reliée à la précédente ? Il doit y avoir \emph{au moins} un connecteur, une reprise pronominale, un substitut lexical ou un mot du champ lexical.
\item Vérifiez chaque connecteur : exprime-t-il le bon rapport logique (cause, conséquence, opposition) ?
\item Comptez les verbes conjugués de vos phrases trop longues : au-delà de trois propositions, découpez.
\end{enumerate}
\end{methode}

\begin{exoresolu}[Améliorer la cohésion d'un paragraphe de résumé]
Énoncé. Voici un résumé mal articulé d'un texte sur la déforestation du bassin du Congo : « Le bassin du Congo possède la deuxième forêt tropicale du monde. Des milliers d'hectares disparaissent chaque année. L'agriculture itinérante détruit la forêt. L'exploitation minière détruit la forêt. Les espèces animales sont menacées. Le climat de la planète est perturbé. Il faut protéger cette forêt. » Améliorez la cohésion de ce paragraphe sans changer les idées.
\tcblower
Solution détaillée. Les idées sont toutes présentes (cohérence respectée), mais les phrases sont juxtaposées sans liens : ni connecteurs, ni reprises, et « détruit la forêt » est répété. On introduit des connecteurs logiques fidèles au sens (cause, addition, conséquence, conclusion), des reprises pronominales et un substitut lexical :\\
« Le bassin du Congo possède la deuxième forêt tropicale du monde. \emph{Pourtant}, ce massif forestier perd des milliers d'hectares chaque année. \emph{En effet}, l'agriculture itinérante \emph{et} l'exploitation minière le détruisent. \emph{Par conséquent}, les espèces animales sont menacées \emph{et} le climat de la planète s'en trouve perturbé. \emph{C'est pourquoi} il faut protéger cette forêt. »\\
Le paragraphe compte désormais cinq phrases reliées par « pourtant » (opposition), « en effet » (explication), « et » (addition), « par conséquent » (conséquence) et « c'est pourquoi » (conclusion) ; le substitut « ce massif forestier » et le pronom « le » assurent la cohésion lexicale et grammaticale.
\end{exoresolu}

\begin{exercice}[À vous de jouer]
1. Classez les connecteurs suivants par fonction logique : \emph{cependant, car, de plus, par conséquent, certes... mais, ainsi}.\\
2. Dites si chaque phrase est simple ou composée : a) « La forêt abrite des milliers d'espèces. » b) « Le bois est exporté, mais la forêt disparaît. » c) « Il pleut ; les rivières débordent. »\\
3. Condensez en une seule phrase : « Le bassin du Congo est une région immense. Il abrite la deuxième forêt tropicale du monde. »
\end{exercice}

\begin{corrige}[Exercice : cohésion et structure de la phrase]
1. Opposition : \emph{cependant} ; cause : \emph{car} ; addition : \emph{de plus} ; conséquence : \emph{par conséquent, ainsi} ; concession : \emph{certes... mais}.\\
2. a) Phrase simple (un verbe : \emph{abrite}). b) Phrase composée par coordination (\emph{est exporté}, \emph{disparaît} ; conjonction \emph{mais}). c) Phrase composée par juxtaposition (\emph{pleut}, \emph{débordent} ; point-virgule).\\
3. « Le bassin du Congo, région immense, abrite la deuxième forêt tropicale du monde » (condensation par apposition) ou « Le bassin du Congo, qui est une région immense, abrite la deuxième forêt tropicale du monde » (subordonnée relative).
\end{corrige}

\begin{aretenir}[L'essentiel]
La \cle{cohérence} garantit l'unité de sens et l'enchaînement logique des idées ; la \cle{cohésion} la rend visible grâce aux reprises pronominales, aux substituts lexicaux, aux champs lexicaux et aux connecteurs logiques. La \cle{phrase simple} (un verbe conjugué) et la \cle{phrase composée} (plusieurs propositions, par juxtaposition ou coordination avec \emph{mais, ou, et, donc, or, ni, car}) sont vos deux outils pour condenser ou alléger un résumé. Un connecteur mal choisi fausse le sens : choisissez-le toujours en fonction du rapport logique du texte de départ.
\end{aretenir}

\coursec{Techniques de réduction et de reformulation}

Dans la section précédente, nous avons appris à repérer la \cle{cohésion} et la \cle{cohérence} d'un texte, c'est-à-dire les fils invisibles (reprises, connecteurs, pronoms) qui relient les phrases entre elles, ainsi que la structure de la phrase simple et de la phrase composée. Ces outils nous servent maintenant directement : pour résumer un texte, il ne suffit pas de comprendre comment il est construit, il faut savoir le \emph{réduire} sans le déformer. C'est l'objet de cette section : maîtriser les quatre \cle{techniques de réduction} (suppression, généralisation, condensation, substitution) et l'art de la \cle{reformulation}, qui consiste à dire les idées de l'auteur avec ses propres mots.

\courssub{Le principe général : réduire sans trahir}

\begin{definition}[Résumé et contraction de texte]
La \cle{contraction de texte} est un exercice qui consiste à réduire un texte à une fraction de sa longueur (généralement le quart, par exemple 250 mots pour un texte de 1000 mots), en conservant \textbf{toutes les idées essentielles} et en éliminant tout ce qui est secondaire. Le résultat s'appelle un \cle{résumé}.
\end{definition}

L'intuition est simple : imaginez une mangue. Le texte complet, c'est le fruit entier, avec sa peau, sa pulpe et son noyau. Le résumé, c'est le noyau : il contient l'essentiel, ce qui permettrait de « replanter » l'idée, mais il est beaucoup plus petit. Notre travail consiste donc à enlever la peau et la pulpe (les exemples, les répétitions, les détails) sans jamais toucher au noyau (la thèse, les arguments, la logique).

\begin{aretenir}[La règle d'or du résumé]
Un bon résumé respecte trois interdictions absolues :
\begin{itemize}
\item ne \textbf{rien ajouter} au texte (ni idée nouvelle, ni jugement personnel, ni exemple de votre cru) ;
\item ne \textbf{rien trahir} : le sens, la thèse et le ton de l'auteur doivent être fidèlement respectés ;
\item ne \textbf{rien copier} : on reformule avec ses propres mots, sans citer le texte.
\end{itemize}
\end{aretenir}

\courssub{Technique 1 : la suppression}

\begin{definition}[Suppression]
La \cle{suppression} consiste à éliminer purement et simplement les éléments qui n'apportent rien à l'idée essentielle : les \textbf{répétitions}, les \textbf{exemples}, les \textbf{citations}, les \textbf{comparaisons}, les \textbf{énumérations} détaillées et les \textbf{détails chiffrés secondaires}.
\end{definition}

C'est la technique la plus simple et celle par laquelle on commence toujours. Observons un passage :

\begin{exemplebox}[Observer pour comprendre]
« Le sport est bénéfique pour la santé. Par exemple, la course à pied renforce le cœur. De même, la natation développe les muscles. Comme le disait un sage, "un esprit sain dans un corps sain". Le sport est donc vraiment bon pour la santé. »

Qu'est-ce qui est essentiel ? Une seule idée : \emph{le sport est bénéfique pour la santé}. Les deux exemples (course, natation) illustrent sans ajouter d'idée ; la citation décore ; la dernière phrase répète la première. On supprime tout cela.
\end{exemplebox}

\begin{attention}[Erreur fréquente : garder les exemples « pour faire la longueur »]
Beaucoup d'élèves conservent les exemples du texte parce que leur résumé leur semble « trop court ». C'est une faute grave : le résumé doit être \textbf{dense}, pas \textbf{rempli}. Un résumé qui atteint la longueur demandée avec des exemples recopiés est un mauvais résumé ; un résumé plus court mais fidèle aux idées essentielles est un bon résumé. Au Probatoire et au Baccalauréat technique, chaque exemple conservé est pénalisé.
\end{attention}

\courssub{Technique 2 : la généralisation}

\begin{definition}[Généralisation]
La \cle{généralisation} consiste à remplacer une énumération ou une suite de termes particuliers par un \textbf{terme générique} (un hyperonyme) qui les englobe tous.
\end{definition}

L'intuition : au lieu de nommer chaque membre d'une famille, on nomme la famille elle-même. Ainsi, « manguiers, palmiers, bananiers » devient « arbres fruitiers » ; « cousins, oncles, tantes » devient « parents » ou « membres de la famille ».

\begin{exemplebox}[Exemple chiffré]
« Les élèves de Première, de Terminale et de Seconde se sont rendus au stade » (13 mots) devient « Les élèves du lycée se sont rendus au stade » (8 mots). On a remplacé l'énumération des trois classes par le terme générique « du lycée » : gain de 5 mots, aucune perte de sens.
\end{exemplebox}

\begin{attention}[Piège : le terme générique doit être juste]
Le terme choisi doit recouvrir \emph{exactement} ce qu'il remplace. Dire « les élèves de l'école » pour des lycéens serait imprécis ; dire « les jeunes » serait trop large. Vérifiez toujours que votre terme générique n'est ni trop étroit (il oublie un élément) ni trop large (il ajoute ce que le texte ne disait pas).
\end{attention}

\courssub{Technique 3 : la condensation}

\begin{definition}[Condensation]
La \cle{condensation} consiste à réduire une proposition ou un groupe de mots à une unité plus courte, le plus souvent un \textbf{groupe nominal}, en supprimant le verbe et les connecteurs devenus inutiles.
\end{definition}

C'est la technique la plus rentable en nombre de mots. Quelques transformations types :

\begin{center}
\begin{tabularx}{\linewidth}{|X|X|l|}
\hline
\rowcolor{popblueL} \textbf{Forme longue} & \textbf{Forme condensée} & \textbf{Gain} \\
\hline
parce qu'il pleuvait & à cause de la pluie & 2 mots \\
\hline
lorsque les élèves arrivent & à l'arrivée des élèves & 2 mots \\
\hline
le gouvernement a décidé de construire des routes & la décision gouvernementale de construction routière & 4 mots \\
\hline
il faut que les citoyens respectent la loi & le respect de la loi & 4 mots \\
\hline
\end{tabularx}
\end{center}

\begin{savaistu}[La nominalisation, outil roi de la condensation]
La \cle{nominalisation} consiste à transformer un verbe (ou un adjectif) en nom : « on pollue » devient « la pollution » ; « les jeunes sont chômeurs » devient « le chômage des jeunes » ; « le pays se développe » devient « le développement du pays ». C'est l'outil le plus puissant de la contraction : une phrase entière de 8 ou 10 mots peut se réduire à un groupe nominal de 3 mots. Tous les résumés de qualité en font un usage massif.
\end{savaistu}

\courssub{Technique 4 : la substitution}

\begin{definition}[Substitution]
La \cle{substitution} consiste à remplacer une périphrase (un groupe de mots qui désigne une seule réalité) par un \textbf{mot unique} équivalent.
\end{definition}

Ainsi, « le chef de l'État » devient « le président » ; « l'ensemble des élèves d'une classe » devient « la promotion » ou « la classe » ; « le moyen de transport à deux roues mû par un moteur » devient « la moto ». La substitution est la sœur jumelle de la généralisation, mais elle ne porte pas sur une énumération : elle remplace une \emph{description} par le \emph{terme exact}.

\begin{exemplebox}[Substitution en contexte]
« Celui qui dirige l'établissement scolaire a annoncé que les cours reprendraient lundi » (14 mots) devient « Le proviseur a annoncé la reprise des cours lundi » (9 mots) : on combine ici substitution (« celui qui dirige l'établissement scolaire » $\rightarrow$ « le proviseur ») et condensation (« que les cours reprendraient » $\rightarrow$ « la reprise des cours »). Les techniques se combinent presque toujours.
\end{exemplebox}

\begin{popfigure}
\begin{center}
\begin{tikzpicture}[font=\small]
% Escalier des quatre techniques
\fill[popblueL] (0,0) rectangle (3.4,1);
\draw[popblue] (0,0) rectangle (3.4,1);
\node[align=center] at (1.7,0.5) {\textbf{1. Suppression}\\ barrer le superflu};
\fill[popgreenL] (1.2,1) rectangle (4.6,2);
\draw[popgreen] (1.2,1) rectangle (4.6,2);
\node[align=center] at (2.9,1.5) {\textbf{2. Généralisation}\\ énumération $\rightarrow$ terme générique};
\fill[poporangeL] (2.4,2) rectangle (5.8,3);
\draw[poporange] (2.4,2) rectangle (5.8,3);
\node[align=center] at (4.1,2.5) {\textbf{3. Condensation}\\ proposition $\rightarrow$ groupe nominal};
\fill[poppurpleL] (3.6,3) rectangle (7,4);
\draw[poppurple] (3.6,3) rectangle (7,4);
\node[align=center] at (5.3,3.5) {\textbf{4. Substitution}\\ périphrase $\rightarrow$ mot unique};
\draw[-{Stealth},thick,popdark] (0.4,1.1) -- (1.0,1.9);
\draw[-{Stealth},thick,popdark] (1.6,2.1) -- (2.2,2.9);
\draw[-{Stealth},thick,popdark] (2.8,3.1) -- (3.4,3.9);
\node[align=left,anchor=west] at (7.3,2) {Chaque marche\\ raccourcit le texte\\ sans le trahir.};
\end{tikzpicture}
\end{center}
\legende{Schéma en escalier des quatre techniques de réduction : on les applique dans cet ordre, de la plus simple (suppression) à la plus délicate (substitution), et on les combine librement.}
\end{popfigure}

\courssub{La reformulation : dire autrement sans trahir}

Les quatre techniques réduisent le texte ; la \cle{reformulation} le réécrit. Elle est obligatoire : un résumé fait de phrases recopiées du texte est sanctionné comme un plagiat.

\begin{definition}[Reformulation]
\cle{Reformuler}, c'est exprimer les idées de l'auteur \textbf{avec ses propres mots et ses propres constructions}, sans changer le sens, sans ajouter d'idée et sans en retirer d'essentielle.
\end{definition}

\begin{methode}[Comment reformuler une phrase]
\begin{enumerate}
\item \textbf{Comprendre} la phrase de l'auteur : identifier l'idée (qui ? fait quoi ? pourquoi ?).
\item \textbf{Changer les mots} : utiliser des synonymes (« important » $\rightarrow$ « essentiel » ; « diminuer » $\rightarrow$ « réduire »).
\item \textbf{Changer la construction} : passer de l'actif au passif (« les élèves ont construit le hangar » $\rightarrow$ « le hangar a été construit par les élèves »), ou nominaliser (« les habitants se plaignent » $\rightarrow$ « les plaintes des habitants »).
\item \textbf{Vérifier} : relire en comparant avec l'original. Le sens est-il identique ? Aucune idée ajoutée ? Aucun mot du texte recopié inutilement ?
\end{enumerate}
\end{methode}

\begin{exemplebox}[Reformulation réussie et reformulation fautive]
Phrase du texte : « Les pouvoirs publics doivent agir rapidement pour enrayer l'insécurité grandissante dans nos villes. »

Reformulation correcte : « L'État doit intervenir vite contre la montée de l'insécurité urbaine. » (synonymes + nominalisation, sens conservé.)

Reformulation fautive : « Le gouvernement, qui est incompétent, doit agir contre l'insécurité. » (ajout d'un jugement personnel : « incompétent » n'est pas dans le texte.)
\end{exemplebox}

\courssub{Ce qu'il faut conserver, ce qu'il faut proscrire}

\begin{aretenir}[À conserver absolument]
\begin{itemize}
\item la \textbf{thèse} de l'auteur (l'idée principale qu'il défend) ;
\item les \textbf{arguments} qui justifient cette thèse ;
\item les \textbf{articulations logiques} (d'abord, ensuite, cependant, donc...) qui montrent la progression du raisonnement ;
\item le \textbf{ton général} du texte (critique, admiratif, neutre, polémique).
\end{itemize}
\end{aretenir}

\begin{aretenir}[À proscrire absolument]
\begin{itemize}
\item les \textbf{citations} du texte (même entre guillemets) ;
\item la \textbf{première personne} : on n'écrit jamais « je pense que... » dans un résumé ; on écrit « l'auteur montre que... » ;
\item les \textbf{jugements personnels} : le résumé n'est pas un commentaire, on ne dit pas si l'on est d'accord ;
\item les \textbf{exemples} du texte, sauf si l'exemple est lui-même l'argument (cas rare).
\end{itemize}
\end{aretenir}

\begin{methode}[Méthode complète pour réduire un texte]
\begin{enumerate}
\item \textbf{Lire} le texte deux fois et dégager le plan (thèse, arguments).
\item \textbf{Barrer d'abord} tout ce qui illustre, répète ou détaille : exemples, citations, comparaisons, chiffres secondaires, répétitions.
\item \textbf{Condenser} ce qui reste : généraliser les énumérations, nominaliser les propositions, substituer les périphrases.
\item \textbf{Reformuler} chaque phrase avec ses propres mots, en reliant les idées par des connecteurs logiques.
\item \textbf{Compter les mots} et vérifier la fidélité : rien ajouté, rien trahi, rien copié.
\end{enumerate}
\end{methode}

\begin{popfigure}
\begin{center}
\begin{tikzpicture}[font=\footnotesize]
\node[draw=popblue,fill=popblueL,rounded corners,align=left,text width=5.6cm,minimum height=4.6cm,anchor=north west] (A) at (0,0) {%
\textbf{Passage du texte}\\[2pt]
« Les habitants de Douala se déplacent en taxis, en motos-taxis et en bus. Par exemple, un employé de Bonabéri met deux heures pour rejoindre Akwa. Les embouteillages, qui sont dus au grand nombre de véhicules, ralentissent la circulation. Il faut donc que la ville construise de nouvelles routes. »};
\node[draw=popgreen,fill=popgreenL,rounded corners,align=left,text width=5.6cm,minimum height=4.6cm,anchor=north west] (B) at (7.2,0) {%
\textbf{Passage réduit}\\[2pt]
« Les Doualais utilisent divers transports urbains, mais l'encombrement des véhicules ralentit la circulation ; la ville doit donc construire de nouvelles routes. »\\[4pt]
\textbf{Techniques :}\\
-- généralisation : « taxis, motos-taxis, bus » $\rightarrow$ « divers transports » ;\\
-- suppression : l'exemple de l'employé ;\\
-- condensation : « qui sont dus au grand nombre de véhicules » $\rightarrow$ « l'encombrement des véhicules ».};
\draw[-{Stealth},very thick,poporange] (5.8,-2.3) -- (7.0,-2.3) node[midway,above]{réduire};
\end{tikzpicture}
\end{center}
\legende{Tableau à deux colonnes : passage du texte (gauche) et passage réduit (droite), avec les techniques utilisées. Le passage de 55 mots tombe à 27 mots, soit environ la moitié, sans perte d'idée.}
\end{popfigure}

\courssub{Atelier guidé : réduire de moitié un paragraphe}

\begin{exoresolu}[Réduire un paragraphe de 80 mots sur le transport urbain à Douala]
Énoncé. Réduisez d'environ moitié le paragraphe suivant (80 mots) en appliquant les quatre techniques, puis justifiez chaque choix.

« À Douala, les déplacements sont un véritable cauchemar pour les habitants. Chaque matin, des milliers de travailleurs prennent les taxis, les motos-taxis ou encore les bus pour se rendre à leur lieu de travail. Par exemple, une vendeuse de Deïdo peut passer trois heures sur la route. Les embouteillages sont énormes parce que les routes sont trop étroites et trop peu nombreuses. Les autorités doivent absolument élargir les routes et construire de nouvelles voies. »
\tcblower
Solution détaillée.

Étape 1 — Suppression : on barre l'exemple de la vendeuse de Deïdo (illustration) et la répétition « élargir les routes et construire de nouvelles voies » qui dit deux fois la même chose.

Étape 2 — Généralisation : « les taxis, les motos-taxis ou encore les bus » devient « les transports en commun ».

Étape 3 — Condensation : « parce que les routes sont trop étroites et trop peu nombreuses » devient « à cause de l'étroitesse et de l'insuffisance des routes » ; « prendre les transports pour se rendre à leur lieu de travail » devient « rejoindre leur travail ».

Étape 4 — Reformulation : « un véritable cauchemar » devient « très difficiles » (ton neutre du résumé).

Résumé proposé (40 mots) : « À Douala, les déplacements sont très difficiles : chaque matin, des milliers de travailleurs empruntent les transports en commun pour rejoindre leur travail. L'étroitesse et l'insuffisance des routes provoquent d'énormes embouteillages ; les autorités doivent donc aménager de nouvelles voies. »

Vérification : thèse conservée (les déplacements sont difficiles), arguments conservés (transports saturés, routes insuffisantes, solution attendue), articulation logique conservée (cause $\rightarrow$ conséquence $\rightarrow$ solution), aucun exemple, aucune citation, aucun jugement personnel.
\end{exoresolu}

\begin{exercice}[À vous de jouer]
Réduisez de moitié le passage suivant (environ 60 mots) en appliquant au moins trois des quatre techniques, puis indiquez pour chaque modification la technique utilisée.

« Dans nos villages, les cultivateurs plantent du maïs, du manioc, des arachides et des ignames. Par exemple, mon grand-père cultive deux hectares de manioc chaque année. Mais les récoltes sont mauvaises parce que les pluies sont irrégulières et que le sol s'épuise. Il faut que les agriculteurs utilisent des engrais et pratiquent la rotation des cultures. »
\end{exercice}

\begin{corrige}[Exercice]
Résumé possible (environ 30 mots) : « Dans nos villages, les cultivateurs plantent diverses cultures vivrières, mais les récoltes sont mauvaises à cause de l'irrégularité des pluies et de l'épuisement du sol ; les agriculteurs doivent donc employer des engrais et pratiquer la rotation des cultures. »

Techniques utilisées : généralisation (« maïs, manioc, arachides, ignames » $\rightarrow$ « diverses cultures vivrières ») ; suppression (l'exemple du grand-père) ; condensation (« parce que les pluies sont irrégulières » $\rightarrow$ « à cause de l'irrégularité des pluies » ; « que le sol s'épuise » $\rightarrow$ « l'épuisement du sol »). La solution de l'auteur (engrais, rotation) est un argument : elle est conservée.
\end{corrige}

En résumé, réduire un texte revient à descendre l'escalier des quatre techniques — supprimer, généraliser, condenser, substituer — puis à reformuler le tout avec ses propres mots. La section suivante nous donnera des outils de style précieux pour cette reformulation : les \cle{figures d'atténuation} (litote, euphémisme, prétérition), qui permettent de dire les choses avec nuance.

\coursec{Les figures d'atténuation : litote, euphémisme, prétérition}

Après avoir maîtrisé les \emph{techniques de réduction et de reformulation} (suppression des détails, généralisation, nomination, changement de voix), nous abordons maintenant un niveau d'analyse plus fin : la compréhension des \textbf{figures d'atténuation}. Repérer ces figures dans un texte source est indispensable pour ne pas trahir la pensée de l'auteur lors de la contraction. Une litote mal comprise devient une affirmation faible ; un euphémisme mal restitué devient une approximation ; une prétérition ignorée fait perdre une argumentation implicite. Cette section vous donne les clés pour les identifier, les analyser et les reformuler avec justesse.

\courssub{La litote : dire moins pour faire entendre plus}

\begin{definition}[La litote]
La \cle{litote} est une figure d'atténuation qui consiste à \textbf{dire moins pour suggérer davantage}, en affirmant une idée par la \textbf{négation de son contraire}. Elle repose sur une structure négative (ne... pas, ne... point, sans, sans que) pour exprimer une réalité positive, souvent forte, voire hyperbolique.
\end{definition}

\begin{methode}[Repérer et analyser une litote]
\begin{enumerate}
    \item \textbf{Identifier la structure négative :} repérer « ne... pas », « ne... point », « sans », « nul(le) », « aucun(e) ».
    \item \textbf{Rétablir l'affirmation positive :} transformer « n'est pas triste » en « est joyeux » ; « n'est pas peu fier » en « est très fier ».
    \item \textbf{Évaluer l'intensité :} la litote n'est pas une simple sous-affirmation ; elle renforce souvent le propos par ironie, pudeur ou rhétorique.
    \item \textbf{Déterminer la fonction :} ironie (moquerie voilée), modestie (refus de l'emphase), émotion (pudeur), stratégie argumentative (faire adhérer sans imposer).
\end{enumerate}
\end{methode}

\begin{exemplebox}[Exemples classiques et camerounais]
\begin{itemize}
    \item \textbf{Littérature (Corneille, \emph{Le Cid}) :} « Va, je ne te hais point. » $\rightarrow$ \emph{Je t'aime} (amour fort, pudique, douloureux).
    \item \textbf{Quotidien :} « Ce n'est pas mal du tout. » $\rightarrow$ \emph{C'est très bien / excellent}.
    \item \textbf{Discours politique camerounais :} « La situation n'est pas dépourvue d'espoir. » $\rightarrow$ \emph{La situation est porteuse d'espoir / encourageante}. Le locuteur évite le triomphalisme tout en rassurant.
    \item \textbf{\emph{Le Lion et la perle} (Soyinka) :} Quand Baroka dit : « Je ne suis pas un vieillard tout à fait décati », il affirme par litote sa virilité et sa puissance intactes face à Sidi.
\end{itemize}
\end{exemplebox}

\begin{exoresolu}[Analyse d'une litote dans un extrait de presse]
\textbf{Énoncé :} Extrait : « Le bilan de la CAN 2021 \textbf{n'est pas négligeable} pour l'économie locale. »

\textbf{Analyse :}
\begin{enumerate}
    \item \textbf{Structure :} « n'est pas négligeable » (négation de l'adjectif « négligeable »).
    \item \textbf{Rétablissement :} « est considérable / significatif / important ».
    \item \textbf{Effet :} L'auteur évite l'exagération brute (« est énorme ») pour paraître objectif, mesuré, crédible. C'est une \textbf{stratégie argumentative de modération} qui renforce l'adhésion du lecteur.
\end{enumerate}
\tcblower
\textbf{Reformulation pour résumé :} « Le bilan de la CAN 2021 est significatif pour l'économie locale. » (On lève la litote pour la clarté du résumé, sauf si l'ironie ou la pudeur est le sujet même de l'étude.)
\end{exoresolu}

\courssub{L'euphémisme : adoucir la réalité}

\begin{definition}[L'euphémisme]
L'\cle{euphémisme} est une figure d'atténuation qui consiste à \textbf{remplacer une expression crue, brutale, taboue ou douloureuse par une expression plus douce, plus vague ou plus valorisante}. Il masque la réalité tout en la désignant.
\end{definition}

\begin{propriete}[Mécanismes principaux de l'euphémisme]
\begin{itemize}
    \item \textbf{La périphrase descriptive :} « personne âgée » (pour « vieux »), « malentendant » (pour « sourd »).
    \item \textbf{Le terme technique, médical ou administratif :} « défaillance cardiaque » (pour « crise cardiaque »), « interruption volontaire de grossesse » (pour « avortement »).
    \item \textbf{La métaphore adoucissante :} « s'éteindre », « dormir son dernier sommeil », « partir » (pour « mourir »).
    \item \textbf{Le recours au jargon sociopolitique :} « quartiers précaires / spontanés / d'habitat non planifié » (pour « bidonvilles / taudis »), « personnes en situation de rue » (pour « SDF »), « mesures d'ajustement structurel » (pour « austérité / coupes budgétaires »).
\end{itemize}
\end{propriete}

\begin{exemplebox}[Euphémismes dans le contexte camerounais et dans \emph{Le Lion et la perle}]
\begin{center}
\begin{tabularx}{\linewidth}{|l|X|X|}
\hline
\rowcolor{popblueL}\textbf{Domaine} & \textbf{Euphémisme (forme adoucie)} & \textbf{Réalité masquée (forme brute)} \\
\hline
Vie quotidienne / Santé & « Il a \textbf{perdu la vue} » / « \textbf{non-voyant} » & « Il est \textbf{aveugle} » \\
\hline
Mort / Deuil & « Il s'est \textbf{éteint} » / « Il a \textbf{rejoint ses ancêtres} » & « Il est \textbf{mort} » \\
\hline
Urbanisme / Social & « \textbf{Quartiers précaires} » / « \textbf{habitat spontané} » & « \textbf{Bidonvilles} » / « \textbf{taudis} » \\
\hline
Emploi / Économie & « \textbf{Libérer du personnel} » / « \textbf{mise en disponibilité} » & « \textbf{Licencier} » / « \textbf{renvoyer} » \\
\hline
\emph{Le Lion et la perle} & Baroka : « \textbf{prendre une nouvelle épouse} » / « \textbf{étendre mon toit} » & « \textbf{Se marier (polygamie)} » / « \textbf{conquérir Sidi} » \\
\hline
\emph{Le Lion et la perle} & Sidi (sur la photo) : « \textbf{Mon image a voyagé} » & « \textbf{J'ai été photographiée / exposée} » \\
\hline
\end{tabularx}
\end{center}
\end{exemplebox}

\begin{attention}[Piège majeur : litote ou euphémisme ?]
\textbf{Ne les confondez pas !}
\begin{itemize}
    \item \textbf{Litote :} \cle{négation du contraire} pour \textbf{affirmer fortement}. Exemple : « Il n'est pas bête » (= il est intelligent). Structure : \emph{ne... pas} + antonyme.
    \item \textbf{Euphémisme :} \cle{substitution lexicale} pour \textbf{adoucir ou voiler}. Exemple : « Il est non-voyant » (= il est aveugle). Structure : un \emph{mot doux} remplace un \emph{mot cru}.
\end{itemize}
\textbf{Test rapide :} si vous pouvez retirer la négation et mettre l'antonyme $\rightarrow$ litote. Si vous devez changer le vocabulaire pour retrouver le sens cru $\rightarrow$ euphémisme.
\end{attention}

\courssub{La prétérition : dire en feignant de taire}

\begin{definition}[La prétérition (ou paralipse)]
La \cle{prétérition} est une figure d'atténuation (souvent classée parmi les figures d'argumentation) qui consiste à \textbf{déclarer que l'on ne parlera pas de quelque chose, tout en le disant implicitement ou explicitement dans la même phrase}. Formules types : « Je ne parlerai pas de... », « Je passe sous silence... », « Sans mentionner... ».
\end{definition}

\begin{methode}[Analyser une prétérition]
\begin{enumerate}
    \item \textbf{Repérer la marque énonciative :} « Je ne dirai pas que... », « Il est inutile de rappeler... », « Passons sur... ».
    \item \textbf{Identifier le contenu « tu » :} quelle information est réellement transmise sous couvert de silence ?
    \item \textbf{Analyser la visée argumentative :} pourquoi feindre le silence ?
    \begin{itemize}
        \item \textbf{Insister sans agresser :} « Je ne rappellerai pas vos échecs passés » (les échecs sont rappelés).
        \item \textbf{Susciter la curiosité ou la complicité :} « Je ne vous ferai pas l'injure de penser que... » (flatterie manipulatoire).
        \item \textbf{Contourner la censure ou la bienséance :} dire l'indicible en prétendant le taire.
    \end{itemize}
\end{enumerate}
\end{methode}

\begin{exemplebox}[Prétérition dans le discours politique et littéraire]
\begin{itemize}
    \item \textbf{Discours politique (Cameroun) :} « \textbf{Je ne veux pas croire} que l'opposition souhaite le chaos. » $\rightarrow$ \emph{Réalité : l'orateur suggère fortement que l'opposition souhaite le chaos, tout en prenant une posture de « bon sens » et de modération.}
    \item \textbf{Baroka (\emph{Le Lion et la perle}, à Sidi) :} « \textbf{Je ne dirai pas} que tu es belle, car tes yeux le savent déjà. » $\rightarrow$ \emph{Il lui dit qu'elle est belle (compliment, séduction) tout en feignant de s'abstenir pour flatter sa vanité. Stratégie de séduction par l'ironie.}
    \item \textbf{Sadiku (à Sidi, sur Baroka) :} « \textbf{Je ne mentionnerai pas} sa virilité légendaire... » $\rightarrow$ \emph{Elle la mentionne explicitement pour appâter Sidi.}
\end{itemize}
\end{exemplebox}

\courssub{Fonctions communes de l'atténuation : pourquoi atténuer ?}

\begin{aretenir}[Les 4 fonctions majeures de l'atténuation (litote, euphémisme, prétérition)]
\begin{center}
\begin{tabularx}{\linewidth}{|l|X|X|}
\hline
\rowcolor{popblueL}\textbf{Fonction} & \textbf{Mécanisme} & \textbf{Effet sur le lecteur / destinataire} \\
\hline
\cle{Politesse / Pudeur} & Euphémisme (mort, maladie, corps), litote (compliments) & Respect du tabou, protection de l'image de soi et d'autrui, maintien du lien social. \\
\hline
\cle{Ironie / Moquerie} & Litote (« Ce n'est pas bête » pour une idée géniale), prétérition (« Je ne dirai pas qu'il est sot ») & Distance critique, complicité avec le lecteur averti, déstabilisation de l'adversaire. \\
\hline
\cle{Stratégie argumentative / Rhétorique} & Prétérition (insister en niant), litote (modérer pour convaincre), euphémisme (cacher la violence d'une mesure) & Persuasion douce, contournement de la résistance, manipulation (euphémisme administratif) ou rhétorique. \\
\hline
\cle{Esthétique / Style} & Litote (élégance classique), prétérition (effet de réel, oralité) & Plaisir du texte, caractérisation des personnages (Baroka = maître de l'atténuation). \\
\hline
\end{tabularx}
\end{center}
\end{aretenir}

\courssub{Repérage guidé : discours politique camerounais et \emph{Le Lion et la perle}}

\begin{experience}[Atelier de repérage : figures d'atténuation dans deux corpus]
\textbf{Consigne :} pour chaque énoncé, identifiez la figure (litote, euphémisme, prétérition), citez l'indice textuel, rétablissez le sens littéral et précisez la fonction.

\begin{center}
\begin{tabularx}{\linewidth}{|l|X|X|X|}
\hline
\rowcolor{popblueL}\textbf{Énoncé} & \textbf{Figure} & \textbf{Sens littéral (reformulation)} & \textbf{Fonction} \\
\hline
1. « Le gouvernement a \textbf{procédé à une rationalisation des effectifs} dans la fonction publique. » (communiqué officiel) & Euphémisme (jargon administratif) & « Licenciements massifs / compression budgétaire » & Adoucir la brutalité sociale ; légitimer par le vocabulaire technique. \\
\hline
2. « \textbf{Je ne soulignerai jamais assez} l'importance de la paix sociale. » (discours officiel) & Prétérition (« je ne soulignerai jamais assez » + énoncé de la chose) & « L'importance de la paix est capitale / vitale. » & Insister solennellement ; engagement personnel. \\
\hline
3. Baroka : « \textbf{Je ne suis pas un homme} à refuser un défi. » (\emph{Le Lion et la perle}) & Litote (négation + infinitif) & « Je suis un homme qui accepte / relève les défis. » & Affirmation virile, provocation mesurée, défi lancé à Lakunle. \\
\hline
4. « \textbf{Il s'est endormi} pour toujours. » (en évoquant un mort) & Euphémisme (métaphore) & « Il est mort. » & Pudeur, respect du défunt, atténuation de la douleur. \\
\hline
5. Lakunle : « \textbf{Je ne parlerai pas} de tes habitudes barbares. » (\emph{Le Lion et la perle}) & Prétérition (« je ne parlerai pas de... » + nomination) & « Tes habitudes sont barbares. » & Insulte voilée, posture morale hypocrite, provocation. \\
\hline
6. « La réforme constitutionnelle \textbf{n'est pas une mince affaire}. » (éditorial) & Litote (négation du contraire) & « C'est une affaire majeure / complexe / risquée. » & Modération lucide, alerte mesurée, crédibilité. \\
\hline
\end{tabularx}
\end{center}
\end{experience}

\courssub{Illustration synthétique : tableau comparatif des trois figures}

\begin{popfigure}
\begin{center}
\begin{tikzpicture}[
    header/.style={fill=popblue, text=white, font=\bfseries, minimum height=0.9cm, align=center},
    cell/.style={align=left, font=\small},
    cellc/.style={align=center, font=\small\bfseries},
]
% En-têtes (colonnes : 0-2.5 / 2.5-8 / 8-13.5)
\node[header, text width=2.3cm] at (1.25,0) {Figure};
\node[header, text width=5.3cm] at (5.25,0) {Définition et mécanisme};
\node[header, text width=5.3cm] at (10.75,0) {Exemple littéraire / camerounais};

% Ligne 1 : Litote
\node[cellc, text width=2.3cm] at (1.25,-1.1) {\textsc{Litote}};
\node[cell, text width=5.3cm, anchor=north west] at (2.6,-0.6) {Dire moins pour dire plus. \textbf{Négation du contraire}. Structure : \emph{ne... pas/point} + antonyme.};
\node[cell, text width=5.3cm, anchor=north west] at (8.1,-0.6) {\emph{Le Cid} : « Va, je ne te hais point. » $\rightarrow$ Je t'aime.\\[2pt] Cameroun : « Le résultat \textbf{n'est pas négligeable}. » $\rightarrow$ Il est important.};

% Ligne 2 : Euphémisme
\node[cellc, text width=2.3cm] at (1.25,-2.9) {\textsc{Euphémisme}};
\node[cell, text width=5.3cm, anchor=north west] at (2.6,-2.3) {Adoucir une réalité crue ou taboue. \textbf{Substitution lexicale}. Mécanismes : périphrase, jargon, métaphore.};
\node[cell, text width=5.3cm, anchor=north west] at (8.1,-2.3) {« \textbf{Quartiers précaires} » pour « bidonvilles ».\\[2pt] \emph{Le Lion et la perle} : « \textbf{prendre une nouvelle épouse} » pour polygamie / conquête.};

% Ligne 3 : Prétérition
\node[cellc, text width=2.3cm] at (1.25,-4.7) {\textsc{Prétérition}};
\node[cell, text width=5.3cm, anchor=north west] at (2.6,-4.1) {Dire en feignant de taire. \textbf{Annonce du silence} + énoncé de la chose tue. Structure : « Je ne dirai pas que... » / « Passons sur... ».};
\node[cell, text width=5.3cm, anchor=north west] at (8.1,-4.1) {Baroka : « \textbf{Je ne dirai pas} que tu es belle. » $\rightarrow$ Tu es belle.\\[2pt] Politique : « \textbf{Je ne veux pas croire} qu'il est corrompu. » $\rightarrow$ Il le suggère.};

% Ligne 4 : Fonctions et impact
\node[cellc, text width=2.3cm] at (1.25,-6.6) {\textsc{Fonctions}};
\node[cell, text width=5.3cm, anchor=north west] at (2.6,-5.9) {\textbf{Politesse / pudeur} (euphémisme de la mort).\\[2pt] \textbf{Ironie} (litote « pas bête »).\\[2pt] \textbf{Argumentation} (prétérition insidieuse, litote modérée).\\[2pt] \textbf{Esthétique} (style, caractérisation).};
\node[cell, text width=5.3cm, anchor=north west] at (8.1,-5.9) {\textbf{Impact pour le résumé :}\\[2pt] 1. Ne pas prendre la litote au pied de la lettre.\\[2pt] 2. Dévoiler l'euphémisme (précision factuelle).\\[2pt] 3. Expliciter la prétérition (restituer l'argument tu).};

% Cadre et filets
\draw[popline, thick] (0,0.5) rectangle (13.5,-8.1);
\draw[popline, thick] (0,-0.5) -- (13.5,-0.5);
\draw[popline] (0,-2.2) -- (13.5,-2.2);
\draw[popline] (0,-4.0) -- (13.5,-4.0);
\draw[popline] (0,-5.8) -- (13.5,-5.8);
\draw[popline, thick] (2.5,0.5) -- (2.5,-8.1);
\draw[popline, thick] (8,0.5) -- (8,-8.1);
\end{tikzpicture}
\end{center}
\legende{Tableau synthétique : litote, euphémisme, prétérition — définitions, mécanismes, exemples camerounais et littéraires (\emph{Le Lion et la perle}), fonctions communes et impact sur la rédaction du résumé.}
\end{popfigure}

\courssub{Enjeu pour la contraction de texte : de la détection à la reformulation}

\begin{propriete}[Règle d'or : lever l'atténuation pour la clarté, la garder pour l'analyse]
Dans un \textbf{résumé} (contraction objective), on \textbf{reformule le sens littéral explicite} :
\begin{itemize}
    \item Litote $\rightarrow$ affirmation positive directe. (« n'est pas négligeable » $\rightarrow$ « est important »).
    \item Euphémisme $\rightarrow$ terme précis, neutre, factuel. (« quartiers précaires » $\rightarrow$ « bidonvilles / habitats insalubres » ; « s'est éteint » $\rightarrow$ « est mort »).
    \item Prétérition $\rightarrow$ affirmation directe du contenu tu. (« Je ne dirai pas qu'il est coupable » $\rightarrow$ « Il est coupable » ou « L'auteur suggère sa culpabilité »).
\end{itemize}
Dans un \textbf{commentaire} ou une \textbf{dissertation}, on \textbf{analyse la figure} : on la cite, on la nomme, on explique sa fonction (ironie, stratégie, pudeur).
\end{propriete}

\begin{exoresolu}[Du repérage à la reformulation pour résumé]
\textbf{Texte source (extrait fictif d'éditorial) :}

« Face à la crise, le Premier Ministre a annoncé une \textbf{réorientation budgétaire}. Il a tenu à préciser que \textbf{l'effort demandé aux ménages ne serait pas insupportable}. \textbf{Je ne mentionnerai pas} les erreurs de gestion passées qui nous ont menés ici. »

\textbf{Analyse des figures :}
\begin{enumerate}
    \item « Réorientation budgétaire » = \textbf{euphémisme} (jargon) pour « coupes budgétaires / austérité / hausse d'impôts ».
    \item « Ne serait pas insupportable » = \textbf{litote} pour « sera supportable / modéré / acceptable ».
    \item « Je ne mentionnerai pas les erreurs... » = \textbf{prétérition} $\rightarrow$ l'orateur \textbf{mentionne et accuse} implicitement la gestion passée.
\end{enumerate}
\tcblower
\textbf{Proposition de phrase pour le résumé :}

« Le Premier Ministre a annoncé des \textbf{mesures d'austérité} (euphémisme levé), assurant qu'elles resteraient \textbf{supportables pour les ménages} (litote levée), tout en \textbf{imputant la crise à la gestion antérieure} (prétérition explicitée). »

\textbf{Pourquoi c'est réussi :} le résumé est factuel, clair, plus court, et ne trahit pas la charge argumentative réelle du texte (l'accusation implicite est rendue explicite).
\end{exoresolu}

\begin{attention}[Erreurs fatales en résumé (à ne pas commettre au Probatoire)]
\begin{enumerate}
    \item \textbf{Garder l'euphémisme :} écrire « Le gouvernement a procédé à une rationalisation des effectifs. » $\rightarrow$ \textbf{Sanction :} imprécision, absence de reformulation, langue de bois reproduite.
    \item \textbf{Inverser la litote :} écrire « Le résultat est négligeable » au lieu de « important ». $\rightarrow$ \textbf{Sanction :} contresens total.
    \item \textbf{Ignorer la prétérition :} ne pas rendre compte de l'accusation voilée. $\rightarrow$ \textbf{Sanction :} perte d'une idée essentielle (argumentative).
    \item \textbf{Confondre litote et euphémisme} dans l'analyse (question de cours) : « C'est un euphémisme de dire "je ne te hais point" ». $\rightarrow$ \textbf{Sanction :} erreur de terminologie.
\end{enumerate}
\end{attention}

\begin{savaistu}[\emph{Le Lion et la perle} : Baroka, maître de l'atténuation]
Dans \emph{Le Lion et la perle} de Wole Soyinka, \textbf{Baroka (le Bale)} est le personnage qui use le plus systématiquement des figures d'atténuation. C'est son \textbf{arme rhétorique principale} pour manipuler Sidi et conserver son pouvoir.
\begin{itemize}
    \item \textbf{Litote :} « Je ne suis pas un vieillard décati » (il affirme sa force et sa virilité).
    \item \textbf{Euphémisme :} « Étendre mon toit », « prendre une épouse » (il masque la conquête prédatrice sous le vocabulaire de la tradition et du devoir).
    \item \textbf{Prétérition :} « Je ne dirai pas que tu es belle... » (séduction par le déni feint).
    \item \textbf{Ironie voilée :} il présente sa ruse (sa fausse impuissance) comme une « confidence » ou une « faiblesse » partagée.
\end{itemize}
\textbf{Astuce Bac :} si le sujet de résumé ou de commentaire porte sur un extrait mettant en scène Baroka, \textbf{repérez systématiquement ses atténuations}. Elles sont la clé de sa psychologie (ruse, hypocrisie, pouvoir) et de la visée du texte (critique du pouvoir traditionnel manipulateur). Ne les « corrigez » pas bêtement dans le résumé ; \textbf{reformulez le sens réel de sa manipulation}.
\end{savaistu}

\begin{exercice}[Entraînement : repérage et reformulation]
\textbf{Consigne :} pour chaque phrase, identifiez la figure d'atténuation, expliquez le mécanisme en une ligne, et proposez une reformulation adaptée à un \textbf{résumé factuel}.
\begin{enumerate}
    \item « Le rapport de la Cour des Comptes \textbf{ne manque pas de sel} quant à la gestion du fonds Covid. »
    \item « La municipalité a décidé de \textbf{redéployer les agents} du service propreté. »
    \item « \textbf{Il est superflu de rappeler} que le respect du code de la route sauve des vies. »
    \item « Lakunle affirme que la dot est une \textbf{coutume barbare} qu'il \textbf{ne saurait cautionner}. » (Analyser « ne saurait cautionner ».)
    \item « Sidi, face au miroir, murmure : "\textbf{Je ne suis pas laide}". »
\end{enumerate}
\end{exercice}

\begin{corrige}[Exercice — Figures d'atténuation]
\begin{enumerate}
    \item \textbf{Litote.} « Ne manque pas de sel » = est très piquant / critique / ironique. \textbf{Résumé :} « Le rapport est très critique sur la gestion du fonds Covid. »
    \item \textbf{Euphémisme (jargon administratif).} « Redéployer » = muter, réaffecter, parfois licencier. \textbf{Résumé :} « La municipalité a muté / réaffecté les agents du service propreté. » (Choisir le terme précis selon le contexte.)
    \item \textbf{Prétérition.} « Il est superflu de rappeler » + énoncé de la chose = je le rappelle avec force. \textbf{Résumé :} « L'auteur insiste sur le fait que le respect du code de la route sauve des vies. »
    \item \textbf{Litote (verbe modal + négation).} « Ne saurait cautionner » = ne peut / ne veut pas accepter, refuse fermement. \textbf{Résumé :} « Lakunle refuse fermement la dot, qu'il juge barbare. »
    \item \textbf{Litote.} « Je ne suis pas laide » = je suis belle / je me trouve belle. \textbf{Résumé :} « Sidi se trouve belle et affirme sa beauté. » (Contexte : narcissisme, prise de conscience de son image.)
\end{enumerate}
\end{corrige}

\begin{aretenir}[Check-list « Figures d'atténuation » pour l'épreuve de contraction]
\cle{Avant d'écrire le résumé :} \textbf{relisez le texte source} à la chasse aux trois figures.
\begin{itemize}
    \item \textbf{Litote ?} (ne... pas + contraire) $\rightarrow$ \textbf{basculez en affirmation positive.}
    \item \textbf{Euphémisme ?} (mots doux : « partir », « rationaliser », « précaires ») $\rightarrow$ \textbf{remplacez par le terme factuel précis.}
    \item \textbf{Prétérition ?} (« je ne dirai pas », « passons sur ») $\rightarrow$ \textbf{explicitez l'idée tue.}
\end{itemize}
\cle{Dans l'analyse (commentaire / dissertation) :} citez $\rightarrow$ nommez $\rightarrow$ expliquez la fonction (ironie ? stratégie ? pudeur ?).

\cle{Pour \emph{Le Lion et la perle} :} Baroka = atténuation systématique. La repérer, c'est comprendre la manipulation.
\end{aretenir}

\coursec{Rédaction, confrontation et amélioration du résumé}

Dans la section précédente, nous avons étudié les figures d'atténuation — la litote, l'euphémisme et la prétérition — qui permettent à un auteur de dire moins pour suggérer davantage. Ces procédés appartiennent au style du texte de départ : dans un résumé, on ne les reproduit pas, on restitue l'idée qu'ils portent. Nous voici donc arrivés au moment décisif de la contraction de texte : la \cle{rédaction} proprement dite du résumé, puis sa \cle{confrontation} avec les productions des camarades et son \cle{amélioration}. C'est ici que se joue la réussite à l'épreuve de contraction de texte du Probatoire et du Baccalauréat technique.

\courssub{De l'intuition à la méthode : qu'est-ce que rédiger un résumé ?}

Imaginez que vous deviez raconter à un ami pressé le contenu d'un long documentaire d'une heure, en deux minutes seulement. Vous ne pouvez pas tout dire : vous choisissez l'essentiel, vous l'ordonnez, vous le dites avec vos propres mots, sans ajouter votre avis. Rédiger un résumé, c'est exactement cela, par écrit et avec des règles précises.

\begin{definition}[Le résumé rédigé]
Le \cle{résumé} est un texte bref, continu et rédigé, qui restitue fidèlement les idées essentielles d'un texte de départ, dans un nombre de mots imposé, sans exemple, sans citation, sans commentaire personnel. Il est rédigé avec les mots du résumateur (reformulation) et organisé selon le plan du texte de départ.
\end{definition}

\begin{attention}[L'erreur qui coûte le plus cher]
L'erreur la plus fréquente — et la plus sanctionnée — consiste à introduire son avis personnel : \emph{« je pense que l'auteur a raison »}, \emph{« à mon avis, ce texte est intéressant »}, \emph{« je suis d'accord avec cette idée »}. Le résumé doit rester \cle{neutre} et \cle{fidèle} : vous êtes le porte-parole de l'auteur, pas son juge. Toute marque de subjectivité (\emph{je}, \emph{selon moi}, \emph{heureusement}, \emph{malheureusement}) est une faute de méthode.
\end{attention}

\courssub{Les étapes de la rédaction : du brouillon à la copie}

Rédiger un résumé ne s'improvise pas : c'est un travail en étapes, comme un artisan qui prépare son ouvrage avant de le livrer. La figure ci-dessous présente l'enchaînement complet des opérations, de la lecture du texte jusqu'à la relecture finale.

\begin{popfigure}
\begin{center}
\begin{tikzpicture}[
  node distance=0.55cm,
  etape/.style={rectangle, rounded corners=3pt, draw=popblue, fill=popblueL, thick, text width=5.6cm, align=center, minimum height=0.85cm, font=\small},
  fleche/.style={-{Stealth[length=2.5mm]}, thick, popdark}
]
\node[etape, align=center] (lect){\textbf{1. Lecture analytique}\\ repérer la thèse et les grandes parties};
\node[etape, below=of lect, align=center] (plan){\textbf{2. Brouillon du plan}\\ une partie = une idée essentielle};
\node[etape, below=of plan, align=center] (reduc){\textbf{3. Réduction}\\ supprimer exemples, citations, répétitions};
\node[etape, below=of reduc, align=center] (redac){\textbf{4. Rédaction}\\ un paragraphe par partie, une idée par phrase};
\node[etape, below=of redac, align=center] (compte){\textbf{5. Comptage des mots}\\ ajuster à $\pm$ 10 \% de la longueur imposée};
\node[etape, below=of compte, align=center] (relect){\textbf{6. Relecture et copie}\\ orthographe, cohésion, nombre de mots indiqué};
\draw[fleche] (lect) -- (plan);
\draw[fleche] (plan) -- (reduc);
\draw[fleche] (reduc) -- (redac);
\draw[fleche] (redac) -- (compte);
\draw[fleche] (compte) -- (relect);
\draw[fleche, dashed, poporange] (compte.west) .. controls +(-1.3,0.4) and +(-1.3,-0.4) .. node[left, font=\footnotesize, poporange, align=center]{trop long ou\\ trop court} (redac.west);
\end{tikzpicture}
\end{center}
\legende{Organigramme des étapes de la rédaction d'un résumé. La flèche en pointillés indique la boucle d'ajustement : si le comptage révèle un écart, on retourne à la rédaction avant de recopier.}
\end{popfigure}

\begin{methode}[Rédiger le résumé, paragraphe par paragraphe]
\begin{enumerate}
\item \textbf{Établir le brouillon du plan} : à partir de votre lecture analytique, notez les deux, trois ou quatre grandes parties du texte et, sous chacune, l'idée essentielle en quelques mots.
\item \textbf{Rédiger une introduction} : présentez brièvement le texte (auteur, thème, thèse) en une ou deux phrases.
\item \textbf{Rédiger le développement} : un paragraphe par partie du plan ; dans chaque paragraphe, une idée essentielle par phrase ; reliez les phrases par des connecteurs logiques (\emph{d'abord, ensuite, en effet, cependant, enfin}).
\item \textbf{Supprimer impitoyablement} : aucun exemple, aucune citation, aucune comparaison ornementale, aucune répétition.
\item \textbf{Rédiger une conclusion} : restituez la position finale de l'auteur en une phrase, sans commentaire.
\item \textbf{Compter les mots}, ajuster, puis recopier au propre en indiquant le total à la fin.
\end{enumerate}
\end{methode}

\begin{methode}[Compter les mots et ajuster la longueur]
\begin{enumerate}
\item Comptez les mots du brouillon, ligne par ligne (comptez les mots d'une ligne pleine, multipliez par le nombre de lignes, puis ajustez les lignes incomplètes) : un comptage approximatif mais sérieux suffit.
\item Comparez au nombre imposé : si le résumé est \textbf{trop long}, supprimez des adjectifs, des subordonnées explicatives, des reprises d'idées ; s'il est \textbf{trop court}, réintégrez une idée secondaire importante ou développez légèrement une formulation.
\item Une fois la longueur correcte, \textbf{inscrivez le nombre de mots} à la fin du résumé, entre parenthèses : \emph{(120 mots)}.
\end{enumerate}
\end{methode}

\begin{exemplebox}[Un exemple chiffré pour bien comprendre la tolérance]
Un sujet propose un texte de 480 mots et demande un résumé en \textbf{120 mots}. La tolérance admise est de $\pm$ 10 \% :
\[ 10\% \text{ de } 120 = 12 \quad\Longrightarrow\quad 120 - 12 = 108 \quad\text{et}\quad 120 + 12 = 132. \]
Votre résumé doit donc contenir \textbf{entre 108 et 132 mots}. En dessous de 108 mots, vous avez sacrifié des idées essentielles ; au-dessus de 132 mots, vous n'avez pas assez contracté. Tout écart supérieur est \textbf{pénalisé à l'examen}. De même, l'oubli de l'indication du nombre de mots à la fin du résumé est sanctionné.
\end{exemplebox}

\courssub{L'introduction et la conclusion du résumé}

Le résumé est un texte complet : il possède donc une entrée en matière et une sortie, mais toutes deux sont extrêmement brèves.

\begin{definition}[Introduction et conclusion du résumé]
L'\cle{introduction} du résumé présente le texte de départ en une ou deux phrases : l'auteur (si le sujet le donne), le thème et la thèse défendue. La \cle{conclusion} restitue la position finale de l'auteur — son aboutissement, sa leçon, sa dernière idée — en une phrase, sans y ajouter aucun commentaire personnel.
\end{definition}

\begin{exemplebox}[Modèles d'introduction et de conclusion]
\textbf{Introduction :} \emph{Dans ce texte, l'auteur s'interroge sur la place de la lecture dans la formation du jeune technicien. Il soutient que lire régulièrement développe à la fois la maîtrise de la langue et l'esprit d'analyse.}

\textbf{Conclusion :} \emph{Pour l'auteur, la lecture n'est donc pas un loisir réservé aux littéraires, mais un outil indispensable de la réussite professionnelle.}

Remarquez : nulle part le résumateur ne dit ce qu'il en pense. Il se contente de rapporter.
\end{exemplebox}

\courssub{La grille d'auto-évaluation : se corriger avant d'être corrigé}

Avant de remettre votre résumé, évaluez-le vous-même à l'aide d'une grille. C'est ce que font les correcteurs à l'examen : autant utiliser leurs critères avant eux.

\begin{popfigure}
\begin{center}
\begin{tikzpicture}
\fill[popblue] (0,0) rectangle (12.6,0.8);
\node[white, font=\bfseries\small, anchor=west] at (0.15,0.4) {Critère d'évaluation};
\node[white, font=\bfseries\small, anchor=west] at (8.2,0.4) {Ce que le correcteur vérifie};
\node[white, font=\bfseries\small, anchor=east] at (12.45,0.4) {Barème};
\foreach \y/\c/\d/\b in {
  -0.8/{Fidélité au texte}/{Aucune idée inventée, aucun avis personnel, thèse respectée}/{5},
  -1.6/{Idées essentielles}/{Toutes les grandes parties du plan sont restituées}/{4},
  -2.4/{Respect de la longueur}/{Nombre de mots dans la tolérance, total indiqué à la fin}/{3},
  -3.2/{Cohésion du texte}/{Paragraphes liés, connecteurs logiques, reformulation fluide}/{4},
  -4.0/{Correction de la langue}/{Orthographe, grammaire, ponctuation}/{3},
  -4.8/{Présentation}/{Copie lisible, alinéas, soin}/{1}}{
  \fill[popcream] (0,\y-0.8) rectangle (12.6,\y);
  \node[font=\small\bfseries, anchor=west, popdark] at (0.15,\y-0.4) {\c};
  \node[font=\footnotesize, anchor=west, align=left, text width=4.1cm] at (8.2,\y-0.4) {\d};
  \node[font=\small\bfseries, anchor=east, poporange] at (12.45,\y-0.4) {\b\ / \b};
  \draw[popline] (0,\y-0.8) -- (12.6,\y-0.8);
}
\fill[popgreenL] (0,-5.6) rectangle (12.6,-4.8);
\node[font=\small\bfseries, anchor=west, popdark] at (0.15,-5.2) {TOTAL};
\node[font=\small\bfseries, anchor=east, popdark] at (12.45,-5.2) {20 / 20};
\draw[popdark, thick] (0,0) rectangle (12.6,-5.6);
\draw[popdark] (8.05,0) -- (8.05,-5.6);
\draw[popdark] (11.9,0) -- (11.9,-5.6);
\end{tikzpicture}
\end{center}
\legende{Grille d'évaluation du résumé sur 20 points. Utilisez-la en auto-évaluation avant de rendre votre copie, et en évaluation croisée lors de la confrontation en binômes.}
\end{popfigure}

\begin{aretenir}[Les cinq questions de l'auto-évaluation]
Avant de recopier, posez-vous cinq questions :
\begin{enumerate}
\item \textbf{Fidélité} : ai-je dit uniquement ce que dit l'auteur, sans avis personnel ?
\item \textbf{Longueur} : mon nombre de mots est-il dans la fourchette $\pm$ 10 \%, et l'ai-je indiqué ?
\item \textbf{Cohésion} : mes paragraphes suivent-ils le plan du texte, reliés par des connecteurs ?
\item \textbf{Langue} : ai-je relu l'orthographe, les accords, la ponctuation ?
\item \textbf{Présentation} : ma copie est-elle lisible, avec des alinéas nets ?
\end{enumerate}
\end{aretenir}

\courssub{La confrontation des productions : apprendre en binômes}

Le résumé ne s'améliore pas seulement dans la solitude du brouillon : il progresse par la \cle{lecture croisée}. En binôme, chacun lit le résumé de l'autre, l'annote selon la grille, puis formule des conseils précis.

\begin{methode}[Conduire une lecture croisée efficace]
\begin{enumerate}
\item \textbf{Échangez} les brouillons en silence ; chacun lit le texte de départ puis le résumé de son camarade.
\item \textbf{Annotez} avec la grille : soulignez en vert les passages fidèles et bien formulés, en rouge les avis personnels, les exemples oubliés, les idées manquantes.
\item \textbf{Notez chaque critère} sur la grille et justifiez la note oralement.
\item \textbf{Formulez deux conseils d'amélioration précis} : pas de \emph{« c'est bien »} vague, mais \emph{« ta deuxième phrase donne ton avis, supprime-la »} ou \emph{« il te manque l'idée de la troisième partie »}.
\item \textbf{Réécrivez} : chacun intègre les remarques reçues et produit une \cle{version améliorée} de son résumé.
\end{enumerate}
\end{methode}

\begin{savaistu}[L'évaluation formative : une pédagogie qui fait ses preuves]
Les recherches en sciences de l'éducation montrent que la confrontation des productions entre pairs fait progresser plus vite que la seule correction du professeur. Pourquoi ? Parce que corriger le texte d'un autre oblige à formuler les critères de réussite, donc à les comprendre soi-même ; et recevoir un conseil d'un camarade est souvent mieux accepté qu'une remarque venue « d'en haut ». Cette démarche s'appelle la \cle{pédagogie de l'évaluation formative} : l'évaluation n'est pas une sanction finale, mais un outil au service de l'apprentissage. Dans les classes techniques du Cameroun, les ateliers d'écriture en binômes appliquent exactement ce principe.
\end{savaistu}

\begin{exoresolu}[Corriger et améliorer un résumé fautif]
\textbf{Énoncé.} Voici le résumé produit par un élève, à partir d'un texte de 400 mots sur la formation professionnelle des jeunes (résumé demandé en 100 mots) :

\emph{« Ce texte parle de la formation professionnelle. Je pense que l'auteur a tout à fait raison de dire qu'elle est importante. Il explique que les jeunes doivent apprendre un métier pour trouver du travail, par exemple mon cousin a appris la mécanique et il a trouvé un emploi à Douala. L'auteur ajoute que l'État doit créer des centres de formation. C'est une très bonne idée. En conclusion, ce texte est très intéressant et je le recommande à tous les jeunes. » (environ 85 mots)}

Relevez toutes les fautes de méthode, puis proposez une version améliorée.

\tcblower
\textbf{Solution détaillée.} Relevé des fautes selon la grille :
\begin{itemize}
\item \textbf{Avis personnel} (faute majeure) : \emph{« je pense que l'auteur a tout à fait raison »}, \emph{« c'est une très bonne idée »}, \emph{« je le recommande »}.
\item \textbf{Exemple interdit} : l'histoire du cousin mécanicien à Douala est un exemple ajouté, qui n'appartient même pas au texte.
\item \textbf{Idée essentielle manquante} : la thèse de l'auteur n'est pas clairement énoncée dès l'introduction.
\item \textbf{Conclusion fautive} : \emph{« ce texte est très intéressant »} juge le texte au lieu de restituer la position finale de l'auteur.
\item \textbf{Longueur} : 85 mots, en dessous de la fourchette autorisée ($100 \pm 10$, soit 90 à 110 mots).
\item \textbf{Nombre de mots non indiqué} à la fin.
\end{itemize}

\textbf{Version améliorée :} \emph{Ce texte traite de la formation professionnelle des jeunes et soutient qu'elle est la clé de leur insertion dans l'emploi. L'auteur montre d'abord que l'apprentissage d'un métier donne aux jeunes une compétence recherchée par les entreprises. Il explique ensuite que la formation réduit le chômage en adaptant les savoir-faire aux besoins réels du marché du travail. Enfin, il affirme que les pouvoirs publics doivent multiplier les centres de formation accessibles à tous. Pour lui, former les jeunes, c'est bâtir l'avenir économique du pays.} (96 mots)

La version améliorée est fidèle, neutre, sans exemple ni citation, organisée selon le plan du texte, et sa longueur (96 mots) respecte la fourchette imposée.
\end{exoresolu}

\courssub{Compte rendu et remédiation : le bilan collectif}

Après la confrontation en binômes, la classe procède à un \cle{compte rendu} collectif : le professeur relève au tableau les difficultés les plus fréquentes observées dans les productions, puis propose une \cle{remédiation}, c'est-à-dire des exercices ciblés pour corriger chaque difficulté.

\begin{center}
\begin{tabularx}{\linewidth}{|l|X|X|}\hline
\rowcolor{popblueL} \textbf{Difficulté fréquente} & \textbf{Cause probable} & \textbf{Exercice de remédiation} \\ \hline
Avis personnel dans le résumé & Confusion entre résumé et commentaire & Transformer cinq phrases subjectives en phrases neutres \\ \hline
Résumé trop long & Exemples et citations non supprimés & Barrer tout ce qui est exemple ou citation dans un texte donné \\ \hline
Idées essentielles oubliées & Plan du texte mal dégagé & Retrouver les trois grandes parties d'un texte argumentatif \\ \hline
Manque de cohésion & Absence de connecteurs logiques & Relier des phrases par \emph{d'abord, ensuite, cependant, enfin} \\ \hline
Fautes de langue & Relecture négligée & Dictée de phrases du résumé, correction des accords \\ \hline
\end{tabularx}
\end{center}

\begin{exercice}[Rédaction complète d'un résumé]
Lisez attentivement le texte suivant, puis rédigez-en un résumé en \textbf{80 mots} (tolérance : $\pm$ 10 \%). Indiquez le nombre de mots à la fin.

\emph{« Le téléphone portable a transformé la vie des jeunes Camerounais. D'abord, il facilite la communication : en quelques secondes, un élève de Garoua échange avec son ami de Bafoussam. Ensuite, il est devenu un outil d'apprentissage : grâce à lui, on consulte des leçons, on révise ses cours, on apprend même de nouveaux métiers. Cependant, son usage excessif présente des dangers : il distrait pendant les études et isole parfois celui qui s'y enferme. Il appartient donc à chaque jeune d'en faire un usage maîtrisé, car le portable est un bon serviteur mais un mauvais maître. »}
\end{exercice}

\begin{corrige}[Exercice : rédaction complète d'un résumé]
\textbf{Éléments de correction.} Plan du texte : 1) le portable facilite la communication ; 2) il sert d'outil d'apprentissage ; 3) son usage excessif est dangereux ; 4) conclusion : il faut un usage maîtrisé.

\textbf{Résumé possible :} \emph{Ce texte examine les effets du téléphone portable sur la vie des jeunes Camerounais. L'auteur montre d'abord qu'il facilite grandement la communication à distance entre villes, puis qu'il est devenu un véritable outil d'apprentissage grâce aux ressources numériques. Il souligne cependant les dangers d'un usage excessif : distraction pendant les études et isolement relationnel. Il invite donc chaque jeune à faire un usage maîtrisé et raisonné de cet appareil, car le portable doit rester un bon serviteur et non devenir un mauvais maître.} (81 mots)

La fourchette autorisée était de 72 à 88 mots : le résumé est conforme. Il est fidèle, neutre, sans exemple précis ni citation, et suit le plan du texte.
\end{corrige}

\begin{exercice}[Auto-évaluation et conseils]
Évaluez le résumé que vous venez de produire à l'aide de la grille sur 20 points présentée dans cette leçon. Notez chaque critère, puis rédigez deux conseils d'amélioration que vous vous adressez à vous-même, comme le ferait votre binôme.
\end{exercice}

\begin{corrige}[Exercice : auto-évaluation et conseils]
\textbf{Éléments de correction.} La démarche attendue comporte trois temps : 1) reporter une note justifiée pour chacun des six critères de la grille (fidélité, idées essentielles, longueur, cohésion, langue, présentation) ; 2) repérer les points faibles (par exemple une idée secondaire oubliée, un connecteur manquant, un mot de trop) ; 3) formuler deux conseils précis et applicables, du type : \emph{« ajouter un connecteur entre la deuxième et la troisième phrase »}, \emph{« vérifier le comptage des mots avant de recopier »}. L'important n'est pas la note obtenue, mais la capacité à se corriger soi-même : c'est toute la logique de l'évaluation formative.
\end{corrige}

\begin{aretenir}[L'essentiel de la section]
Rédiger un résumé suit un parcours en six étapes : lecture analytique, brouillon du plan, réduction, rédaction paragraphe par paragraphe, comptage des mots, relecture. La longueur imposée s'accompagne d'une tolérance de $\pm$ 10 \% et le nombre de mots doit figurer à la fin. Le résumé reste toujours neutre : jamais d'avis personnel, jamais d'exemple, jamais de citation. Enfin, la confrontation en binômes, guidée par une grille d'évaluation, puis la réécriture, sont les moyens les plus sûrs de progresser : évaluer, c'est déjà apprendre.
\end{aretenir}

\coursec{Ouverture littéraire : Le Lion et la perle — bilan de l'œuvre}

Dans la section précédente, nous avons appris à rédiger, confronter et améliorer un résumé : repérer les idées essentielles, les reformuler avec nos propres mots, respecter la cohésion et la cohérence du texte. Ces compétences ne servent pas seulement à l'exercice scolaire : elles nous aident aussi à \emph{lire} une œuvre littéraire, c'est-à-dire à en dégager l'essentiel pour en parler avec justesse. Nous allons maintenant mettre ces acquis au service d'un bilan complet de l'œuvre étudiée cette année : \emph{Le Lion et la perle} (\emph{The Lion and the Jewel}) de Wole Soyinka. Ce bilan prépare directement à l'évaluation sur œuvre intégrale prévue au Probatoire et au Baccalauréat technique.

\courssub{Wole Soyinka et le contexte de l'œuvre}

\textbf{Un dramaturge engagé}\par

Wole Soyinka est un écrivain nigérian né en 1934 à Abeokuta, dans le sud-ouest du Nigeria, en pays yoruba. Dramaturge, poète, romancier et essayiste, il écrit en anglais mais puise sa matière dans la culture yoruba : mythes, masques, chants, danses et proverbes. En 1986, il devient le \cle{premier Africain} à recevoir le prix Nobel de littérature, récompensé pour une œuvre qui, selon le jury, « façonne le drame de l'existence avec une large perspective culturelle ».

\begin{savaistu}[Un écrivain qui paie de sa personne]
Wole Soyinka n'est pas seulement un homme de théâtre : c'est un citoyen engagé. En 1967, pendant la guerre du Biafra, il est arrêté et emprisonné pendant près de deux ans, en grande partie au secret, pour avoir dénoncé la guerre et appelé à la paix. Son théâtre engage autant qu'il divertit : derrière le rire, il y a toujours une réflexion sérieuse sur le pouvoir, la liberté et l'identité africaine.
\end{savaistu}

\textbf{Le Nigeria postcolonial : tradition contre modernité}\par

La pièce est écrite en 1959, à la veille de l'indépendance du Nigeria (1\textsuperscript{er} octobre 1960) — la même année 1960 où le Cameroun accède aussi à l'indépendance (1\textsuperscript{er} janvier 1960). C'est un moment charnière : les sociétés africaines, libérées de la colonisation, se posent une question brûlante : \emph{que garder de la tradition, que prendre à la modernité occidentale ?} Soyinka met en scène ce débat à travers un village imaginaire, Ilujinle, où s'affrontent deux visions du monde.

\begin{popfigure}
\begin{tikzpicture}[x=1cm,y=1cm]
\draw[->,thick,popdark] (0,0) -- (11,0);
\foreach \x/\date/\evt in {1.2/1959/{Écriture de la pièce},5.5/1960/{Indépendances du Nigeria et du Cameroun},9.8/1986/{Prix Nobel de Soyinka}}{
  \draw[thick,popblue] (\x,0.15) -- (\x,-0.15);
  \node[above,font=\bfseries\small,text=popblue] at (\x,0.2) {\date};
  \node[below,font=\footnotesize,text width=3cm,align=center] at (\x,-0.25) {\evt};
}
\end{tikzpicture}
\legende{Frise chronologique : l'œuvre de Soyinka s'inscrit dans le grand mouvement des indépendances africaines.}
\end{popfigure}

\courssub{Bilan des thèmes}

\textbf{Le choc entre tradition et modernité}\par

Le conflit central de la pièce oppose deux hommes qui veulent épouser Sidi, la plus belle fille du village :

\begin{itemize}
\item \cle{Baroka}, le « Lion », chef du village, âgé de soixante-deux ans, polygame, rusé : il incarne la \cle{tradition}, la sagesse ancienne, mais aussi l'habileté politique.
\item \cle{Lakunle}, jeune instituteur, qui rêve de progrès, refuse la dot, veut une femme « moderne » et libérée : il incarne la \cle{modernité} — mais une modernité souvent superficielle, faite de mots étrangers et de mépris pour les coutumes.
\end{itemize}

\begin{exemplebox}[Un titre symbolique]
Le titre même de l'œuvre est symbolique : le \emph{Lion} (force, autorité, tradition) et la \emph{perle} (la beauté convoitée, Sidi) annoncent déjà l'issue du conflit. Dès la couverture du livre, Soyinka suggère que c'est la tradition, incarnée par Baroka, qui remportera le « joyau ». Le titre est donc une première clé d'interprétation : à l'examen, pensez toujours à commenter le titre d'une œuvre.
\end{exemplebox}

Attention : Soyinka ne dit pas simplement « la tradition a raison, la modernité a tort ». Il critique les deux camps : Baroka est manipulateur, Lakunle est ridicule et lâche. La pièce invite à une réflexion nuancée sur ce que l'Afrique doit conserver et transformer.

\textbf{Le statut de la femme}\par

Sidi, la « perle », est au cœur du débat. D'abord simple fille du village, elle devient orgueilleuse lorsqu'un magazine publie ses photographies : sa beauté lui donne le sentiment d'être supérieure à tous, même au chef. À travers elle, Soyinka interroge la place de la femme : objet de désir et de conquête, la femme peut-elle devenir sujet de son propre destin ?

\textbf{La ruse et le pouvoir}\par

Baroka feint d'être devenu impuissant pour attirer Sidi dans son palais et la séduire. Cette ruse montre que le pouvoir traditionnel ne repose pas seulement sur la force : il repose sur l'intelligence, la parole et la connaissance des hommes. La pièce est aussi une leçon de politique.

\begin{definition}[La satire]
La \cle{satire} est une œuvre qui critique les défauts d'une société, d'un personnage ou d'une institution par le rire, l'ironie et l'exagération. Au théâtre, la satire fait rire le spectateur tout en l'obligeant à réfléchir sur ses propres travers.
\end{definition}

\emph{Le Lion et la perle} est une satire : Lakunle, avec ses grands mots anglais et ses idées mal digérées, tourne en ridicule la modernité mal comprise ; Baroka, lui, révèle les travers d'un pouvoir traditionnel qui use de ruse. Le rire devient une arme critique.

\courssub{Bilan des personnages et de leur évolution}

\begin{popfigure}
\begin{tikzpicture}[scale=1.05]
\node[draw,fill=popgoldL,rounded corners,font=\bfseries,text width=2.6cm,align=center] (baroka) at (0,0) {Baroka\\ \footnotesize le Lion : tradition, ruse};
\node[draw,fill=poppinkL,rounded corners,font=\bfseries,text width=2.6cm,align=center] (sidi) at (4,3.4) {Sidi\\ \footnotesize la perle : beauté, orgueil};
\node[draw,fill=popblueL,rounded corners,font=\bfseries,text width=2.6cm,align=center] (lakunle) at (8,0) {Lakunle\\ \footnotesize l'instituteur : modernité};
\draw[<->,thick,popdark] (baroka) -- (sidi) node[midway,left,font=\footnotesize,align=center,text width=2.2cm]{séduit par ruse};
\draw[<->,thick,popdark] (sidi) -- (lakunle) node[midway,right,font=\footnotesize,align=center,text width=2.2cm]{refuse la dot};
\draw[<->,thick,poporange] (baroka) -- (lakunle) node[midway,below,font=\footnotesize]{rivalité : tradition / modernité};
\end{tikzpicture}
\legende{Triangle des relations : Baroka et Lakunle s'affrontent pour conquérir Sidi, la perle du village.}
\end{popfigure}

\textbf{La transformation de Sidi}\par

Sidi évolue tout au long de la pièce. Au début, elle est une jeune fille fière mais simple. Les photographies du magazine la rendent vaniteuse : elle se croit désormais au-dessus des lois du village. Moquant Baroka qu'elle croit impuissant, elle se rend à son palais pour l'humilier — et tombe dans son piège. À la fin, contre toute attente, elle choisit d'épouser Baroka et non Lakunle. Sa transformation est double : elle perd ses illusions de vanité, mais elle gagne une forme de maturité en choisissant, elle-même, son destin.

\textbf{La victoire ambiguë de Baroka}\par

Baroka triomphe : il épouse la perle du village. Mais sa victoire est \emph{ambiguë} : elle est obtenue par le mensonge. Soyinka laisse le spectateur juge : faut-il admirer l'intelligence du vieux chef ou condamner sa manipulation ? Cette ambiguïté est la richesse de l'œuvre : elle interdit toute lecture simpliste.

\courssub{Évaluation de l'étude de l'œuvre}

\begin{methode}[Réussir l'évaluation sur une œuvre intégrale]
\begin{enumerate}
\item \textbf{Réviser le contexte} : auteur, date, pays, situation historique (ici : Soyinka, 1959, Nigeria à la veille de l'indépendance).
\item \textbf{Réviser les thèmes} : dresser la liste des thèmes majeurs (tradition/modernité, femme, pouvoir, ruse) avec pour chacun un exemple précis tiré de la pièce.
\item \textbf{Réviser les personnages} : fiche par personnage (identité, rôle, évolution, rapport aux autres).
\item \textbf{Savoir citer un passage précis} : retenir deux ou trois répliques ou scènes clés (par exemple la scène de la ruse de Baroka) pour appuyer l'analyse.
\item \textbf{Structurer la réponse} : introduction (présentation de l'œuvre et de la problématique), développement en deux ou trois parties, conclusion.
\end{enumerate}
\end{methode}

\begin{attention}[Résumer n'est pas analyser]
L'erreur la plus fréquente à l'examen consiste à \emph{raconter l'intrigue} au lieu de l'\emph{interpréter}. Dire « Baroka attire Sidi dans son palais et l'épouse » n'apporte rien : l'examinateur connaît l'histoire. Il faut expliquer \emph{ce que cela signifie} : « En triomphant par la ruse, Baroka montre que la tradition survit grâce à son intelligence et à sa capacité d'adaptation, ce qui nuance le portrait d'une modernité triomphante. » Analyser, c'est répondre à la question \emph{pourquoi} et \emph{qu'est-ce que l'auteur veut dire}, pas \emph{que se passe-t-il}.
\end{attention}

\begin{exoresolu}[Quiz de synthèse sur l'œuvre]
\textbf{Énoncé.} Répondez brièvement :
\begin{enumerate}
\item Qui est l'auteur du \emph{Lion et la perle} et quelle distinction majeure a-t-il reçue ?
\item Quel événement historique de 1960 éclaire le thème central de la pièce ?
\item Pourquoi Sidi refuse-t-elle d'abord d'épouser Lakunle ?
\item Par quel procédé Baroka parvient-il à séduire Sidi ?
\item Pourquoi dit-on que la victoire de Baroka est ambiguë ?
\end{enumerate}
\tcblower
\textbf{Solution détaillée.}
\begin{enumerate}
\item Wole Soyinka, dramaturge nigérian ; il a reçu le prix Nobel de littérature en 1986, premier Africain ainsi récompensé.
\item L'indépendance du Nigeria (et du Cameroun) en 1960 : la pièce, écrite en 1959, reflète le débat de l'époque entre tradition et modernité dans les sociétés africaines nouvellement libérées.
\item Parce que Lakunle refuse de payer la dot, ce qui, aux yeux de Sidi et du village, déshonorerait la mariée ; elle refuse aussi d'être traitée comme une épouse « moderne » sans valeur.
\item Par la ruse : il fait croire qu'il est devenu impuissant, ce qui pousse Sidi, moqueuse, à venir le narguer dans son palais, où il la séduit.
\item Parce qu'elle est obtenue par le mensonge et la manipulation : on peut admirer l'intelligence du chef tout en condamnant son procédé. Soyinka laisse le jugement ouvert.
\end{enumerate}
\end{exoresolu}

\begin{exercice}[Questions de synthèse et commentaire de passage]
\begin{enumerate}
\item Montrez, en vous appuyant sur deux scènes précises, comment Soyinka critique à la fois la tradition et la modernité.
\item Le choix final de Sidi (épouser Baroka) est-il une victoire ou une défaite pour la femme africaine ? Rédigez un paragraphe argumenté de dix à quinze lignes.
\item Commentaire de passage : relisez la scène où Baroka révèle sa ruse à Sidi. Étudiez le ton (ironie, comique), les enjeux du dialogue et ce que cette scène révèle du pouvoir.
\end{enumerate}
\end{exercice}

\begin{corrige}[Exercice : question 2, éléments de réponse]
Une bonne réponse nuancée peut soutenir les deux thèses. \emph{Victoire} : Sidi choisit elle-même son époux, elle n'est plus un simple objet convoité ; elle accède au statut d'épouse du chef, position de prestige. \emph{Défaite} : son choix est le résultat d'une manipulation ; elle reste prise dans un système polygame où la femme est un trophée (la « perle »). La conclusion la plus riche : Soyinka montre une liberté réelle mais limitée par les structures sociales, invitant le spectateur à réfléchir au statut de la femme dans la société africaine moderne.
\end{corrige}

\courssub{Lien avec la contraction de texte : résumé d'un article de presse}

Les techniques de résumé apprises dans cette unité s'appliquent directement aux sujets de société que l'œuvre de Soyinka nous a fait découvrir. Au Probatoire et au Baccalauréat technique, l'épreuve de contraction porte souvent sur un article de presse traitant, par exemple, de la place de la tradition dans le Cameroun moderne : rôle des chefferies traditionnelles, persistance des rites (funérailles, mariages, cérémonies d'initiation), tension avec la modernité urbaine (villes, école, nouveaux modes de vie).

\begin{methode}[Résumer un article sur tradition et modernité]
\begin{enumerate}
\item \textbf{Lire deux fois} l'article en soulignant la thèse de l'auteur et les idées principales.
\item \textbf{Dégager le plan} : par exemple « constat (les traditions persistent) / causes (identité, cohésion sociale) / limites ou tensions (conflits avec la vie moderne) ».
\item \textbf{Supprimer} les exemples secondaires, les anecdotes, les répétitions, les citations longues.
\item \textbf{Reformuler} chaque idée essentielle avec ses propres mots, en phrases simples ou composées correctement construites.
\item \textbf{Lier} les idées par des connecteurs logiques (d'abord, ensuite, en effet, cependant, enfin) pour garantir cohésion et cohérence.
\item \textbf{Compter les mots} et vérifier que la consigne (par exemple le quart du texte) est respectée.
\end{enumerate}
\end{methode}

\begin{exoresolu}[Résumé d'un court passage de presse]
\textbf{Énoncé.} Résumez en environ 40 mots le passage suivant (environ 120 mots) :

« Dans les grandes villes camerounaises comme Douala et Yaoundé, les jeunes générations semblent parfois oublier les coutumes de leurs villages. Pourtant, dès qu'un décès survient, c'est vers le village que tout le monde retourne : les funérailles traditionnelles rassemblent des centaines de personnes, avec leurs danses, leurs chants et leurs rites séculaires. Les chefs traditionnels, loin d'avoir disparu, jouent encore un rôle de médiation dans les conflits fonciers et familiaux. Certains intellectuels s'inquiètent néanmoins de la folklorisation de ces traditions, réduites parfois à des spectacles pour touristes. La question reste ouverte : comment concilier la modernité urbaine et l'héritage des ancêtres ? »
\tcblower
\textbf{Solution (environ 40 mots).} Malgré la modernité urbaine qui éloigne les jeunes des coutumes, les traditions restent vivaces au Cameroun : rites funéraires et médiation des chefs perdurent. Reste à éviter leur folklorisation et à concilier héritage ancestral et vie moderne.
\end{exoresolu}

\begin{exercice}[À vous de résumer]
À partir d'un article de presse récent (fourni par l'enseignant) sur les chefferies traditionnelles ou les rites au Cameroun, rédigez un résumé respectant le quart du texte. Confrontez ensuite votre production avec celle d'un camarade : vérifiez ensemble que toutes les idées essentielles sont présentes, qu'aucune phrase n'est recopiée, et que les connecteurs logiques assurent la cohérence. Améliorez votre version.
\end{exercice}

\courssub{Pour aller plus loin (complément utile au Bac, hors strict programme)}

La satire que nous avons découverte chez Soyinka est une arme majeure du théâtre africain, anglophone comme francophone. Au Cameroun, des auteurs comme Guillaume Oyono-Mbia (\emph{Trois prétendants… un mari}) utilisent le même procédé : le rire pour dénoncer les travers de la société — ici encore le conflit entre coutumes et vie moderne. Retenir ce parallèle permet, à l'examen, d'élargir intelligemment une conclusion et de montrer une culture littéraire africaine.

\begin{aretenir}[L'essentiel du bilan]
\begin{itemize}
\item \emph{Le Lion et la perle} (1959) de Wole Soyinka, dramaturge nigérian, prix Nobel 1986, met en scène le conflit tradition (Baroka) / modernité (Lakunle) autour de Sidi, la perle.
\item Thèmes majeurs : choc des cultures, statut de la femme, ruse et pouvoir ; le titre est symbolique.
\item La pièce est une \cle{satire} : elle critique par le rire les deux camps.
\item À l'examen : analyser, interpréter, citer — jamais se contenter de raconter.
\item Les techniques de contraction (plan, réduction, reformulation, cohésion) s'appliquent aux sujets de société liés à l'œuvre, comme la tradition dans le Cameroun moderne.
\end{itemize}
\end{aretenir}

\newpage

\coursec{Activité d'intégration}

\courssub{Objectifs de l'activité}

À l'issue de cette activité d'intégration, l'élève sera capable de :

\begin{itemize}
\item repérer la \cle{structure argumentative} d'un article de presse et en construire le \cle{plan} en trois ou quatre parties ;
\item appliquer les \cle{techniques de réduction} (suppression, généralisation, condensation, substitution) à un texte authentique ;
\item \cle{reformuler} les idées essentielles avec ses propres mots, sans recopier le texte ;
\item rédiger un \cle{résumé} cohérent et cohésif respectant une contrainte de longueur ($125$ mots $\pm$ 10 \%) ;
\item évaluer la production d'un camarade à l'aide d'une \cle{grille de notation} et améliorer sa propre copie après confrontation.
\end{itemize}

\courssub{Situation}

Dans le cadre de la semaine de la presse organisée par le club journalistique de ton lycée technique de Douala, le chef d'établissement souhaite publier, sur le panneau d'affichage, une sélection d'articles d'actualité résumés par les élèves de Première. Tu as été choisi(e) pour résumer l'article ci-dessous, paru dans un quotidien camerounais. Ton résumé devra faire exactement \textbf{125 mots ($\pm$ 10 \%)}, soit entre \textbf{113 et 137 mots}, titre non compris.

\begin{attention}[Bien calculer la marge autorisée]
10 \% de 125 mots représentent 12,5 mots. La fourchette admise est donc de $125 - 12{,}5 = 112{,}5$ à $125 + 12{,}5 = 137{,}5$ mots, arrondie à \textbf{113--137 mots}. Tout dépassement, même d'une dizaine de mots, est sanctionné à l'examen : compte tes mots avant de rendre la copie.
\end{attention}

\textbf{Article à résumer (environ 500 mots) :}

\begin{quote}
\textbf{Les motos-taxis dans les villes camerounaises : entre nécessité économique et insécurité}

\emph{Paru dans Cameroon Tribune, rubrique « Société ».}

À Douala, Yaoundé, Bafoussam ou Garoua, impossible de traverser la ville sans croiser les « benskineurs », ces conducteurs de motos-taxis qui fendent la circulation à toute allure. Devenus en quelques années un acteur incontournable du transport urbain, ils suscitent à la fois l'admiration et l'inquiétude des populations.

D'abord, il faut le reconnaître : les motos-taxis répondent à un besoin réel. Dans les quartiers non desservis par les bus, dans les rues en terre battue impraticables aux voitures pendant la saison des pluies, la moto est souvent le seul moyen de transport disponible. Elle est rapide, accessible et bon marché : pour 200 à 500 francs CFA, un habitant d'Ekounou ou de Bonabéri peut rejoindre son lieu de travail en vingt minutes, là où un taxi collectif mettrait plus d'une heure dans les embouteillages. De plus, ce secteur fait vivre des milliers de familles. Selon les syndicats du secteur, on compterait plus de 100 000 conducteurs de motos-taxis dans les seules villes de Douala et Yaoundé. Pour beaucoup de jeunes diplômés sans emploi, la moto achetée à crédit représente une chance de gagner dignement sa vie, de payer les études des cadets et de subvenir aux besoins du foyer.

Cependant, cette activité a un visage sombre. Les accidents impliquant des motos-taxis se comptent par centaines chaque mois. Les services d'urgence des hôpitaux de référence reçoivent quotidiennement des victimes de collisions, de chutes et de renversements. Les causes sont connues : excès de vitesse, non-port du casque, surcharge (il n'est pas rare de voir trois passagers sur une même moto), méconnaissance du code de la route. De nombreux conducteurs n'ont jamais passé le permis de conduire et ont appris le métier « sur le tas ». À cela s'ajoute l'insécurité : certaines motos sont utilisées par des malfaiteurs pour des vols à l'arraché, ce qui jette la suspicion sur toute la profession.

Face à cette situation, les autorités ont pris des mesures. L'immatriculation des motos est désormais obligatoire, le port du casque est exigé pour le conducteur et le passager, et des opérations de contrôle sont régulièrement menées par la police et la gendarmerie. Des mairies ont créé des parcs de stationnement et tenté d'organiser les conducteurs en associations reconnues. Mais l'application de ces règlements reste inégale : les contrôles sont souvent épisodiques, et les amendes ne découragent guère les contrevenants.

Que faire alors ? Interdire purement et simplement les motos-taxis, comme le réclament certains, serait une erreur : cela priverait des milliers de jeunes de leur gagne-pain et compliquerait la mobilité des habitants des quartiers périphériques. La solution passe plutôt par la formation professionnelle des conducteurs, la généralisation du permis de conduire, le renforcement des contrôles et la sensibilisation des usagers. Des expériences menées dans d'autres pays africains montrent qu'un secteur organisé et encadré peut être à la fois rentable et sûr.

En définitive, les motos-taxis sont un mal nécessaire qu'il faut apprivoiser plutôt que combattre. Entre la nécessité économique et l'exigence de sécurité, le Cameroun doit trouver un équilibre : professionnaliser le métier pour que la moto cesse d'être un danger quotidien et demeure ce qu'elle est d'abord, un service rendu à la population.
\end{quote}

\courssub{Consignes}

\begin{enumerate}
\item Lis attentivement l'article deux fois. À la deuxième lecture, souligne les \cle{idées essentielles} et barre les exemples chiffrés, les répétitions et les détails secondaires.
\item Construis le \cle{plan du texte} en quatre parties : pour chaque partie, indique les numéros des paragraphes concernés et formule l'idée principale en une phrase.
\item Applique à chaque partie au moins deux \cle{techniques de réduction} différentes (suppression, généralisation, condensation, substitution). Note sur ton brouillon la technique utilisée.
\item Rédige individuellement un résumé de \textbf{125 mots ($\pm$ 10 \%)}, en respectant les règles du genre : phrases complètes, présent de vérité générale, aucune opinion personnelle, aucune phrase recopiée du texte, connecteurs logiques assurant la \cle{cohésion}.
\item Échange ta copie avec celle de ton binôme. Note la production reçue à l'aide de la grille ci-dessous et rédige un commentaire justificatif de deux lignes par critère.
\item À partir des remarques reçues, réécris une \cle{version améliorée} de ton résumé. La classe construira ensuite un corrigé collectif au tableau à partir des meilleures productions.
\end{enumerate}

\begin{center}
\begin{tabularx}{\linewidth}{|l|X|c|}
\hline
\rowcolor{popblueL} \textbf{Critère} & \textbf{Attendus} & \textbf{Note} \\
\hline
Fidélité au texte & Toutes les idées essentielles sont présentes, sans déformation ni ajout personnel. & /5 \\
\hline
Respect de la longueur & 113 à 137 mots ; phrases complètes ; aucune citation recopiée. & /5 \\
\hline
Cohésion et cohérence & Enchaînement logique des idées, connecteurs variés et correctement employés. & /5 \\
\hline
Qualité de la langue & Orthographe, grammaire, ponctuation, vocabulaire précis. & /5 \\
\hline
\end{tabularx}
\end{center}

\courssub{Ressources à mobiliser}

\begin{aretenir}[Boîte à outils de la contraction]
\begin{itemize}
\item \textbf{Plan du texte} : regrouper les paragraphes par idée directrice ; un article argumentatif suit souvent le schéma \emph{constat — avantages — inconvénients — solutions — conclusion}.
\item \textbf{Techniques de réduction} : \cle{suppression} (exemples, chiffres, répétitions), \cle{généralisation} (« 200 à 500 FCFA » devient « à bas prix »), \cle{condensation} (une phrase pour deux), \cle{substitution} (un mot pour une expression).
\item \textbf{Reformulation} : changer les mots et la construction des phrases sans changer le sens ; transformer une phrase composée en phrase simple quand c'est possible.
\item \textbf{Cohésion} : connecteurs (\emph{d'abord, cependant, en effet, c'est pourquoi, en définitive}), reprises pronominales, champs lexicaux.
\item \textbf{Interdits du résumé} : pas de « je », pas de jugement personnel, pas de citation, pas de titre inventé qui trahit le texte.
\end{itemize}
\end{aretenir}

\begin{methode}[Les cinq étapes du résumé]
\begin{enumerate}
\item \textbf{Lire} deux fois le texte : d'abord pour comprendre, ensuite pour sélectionner.
\item \textbf{Trier} : distinguer les idées essentielles (à garder) des exemples, chiffres et répétitions (à éliminer).
\item \textbf{Planifier} : regrouper les idées en trois ou quatre parties fidèles à la progression du texte.
\item \textbf{Réduire et reformuler} : appliquer les quatre techniques de réduction, avec ses propres mots.
\item \textbf{Rédiger puis vérifier} : rédiger des phrases complètes et liées, compter les mots, relire pour corriger la langue.
\end{enumerate}
\end{methode}

\courssub{Démarche et corrigé commenté}

\begin{exoresolu}[Résumé de l'article sur les motos-taxis]
Énoncé : résumer l'article en 125 mots ($\pm$ 10 \%), en suivant les étapes de la contraction de texte.
\tcblower
\textbf{Étape 1 — Lecture et repérage des idées essentielles.} On barre les détails : les chiffres précis (200 à 500 FCFA, 100 000 conducteurs), les noms de quartiers (Ekounou, Bonabéri), l'exemple des vols à l'arraché. On souligne les idées directrices : (a) les motos-taxis sont partout et utiles ; (b) ils répondent à un besoin de transport et font vivre des milliers de jeunes ; (c) ils provoquent de nombreux accidents et posent un problème d'insécurité ; (d) des mesures existent mais sont mal appliquées ; (e) il faut professionnaliser le secteur plutôt que l'interdire.

\textbf{Étape 2 — Plan du texte.}
\begin{itemize}
\item Partie 1 (§1) : les motos-taxis, un phénomène urbain incontournable.
\item Partie 2 (§2) : leur utilité économique et sociale.
\item Partie 3 (§3) : leurs dangers (accidents, insécurité).
\item Partie 4 (§4 à §6) : les mesures prises, insuffisantes, et la solution : encadrer et professionnaliser.
\end{itemize}

\textbf{Étape 3 — Réduction.} Suppression des exemples chiffrés ; généralisation : « 200 à 500 francs » devient « à moindre coût » ; condensation : les §4 et §5 sont fondus en une seule idée (« les mesures restent mal appliquées, il faut former et contrôler ») ; substitution : « gagner dignement sa vie » devient « vivre de ce métier ».

\textbf{Étape 4 — Rédaction du résumé (126 mots) :}

\emph{Présentes dans toutes les villes du Cameroun, les motos-taxis sont devenues indispensables. D'une part, elles répondent à un besoin réel de transport, rapide et à moindre coût, surtout dans les quartiers mal desservis ; elles permettent en outre à des milliers de jeunes de vivre de ce métier. D'autre part, elles provoquent de nombreux accidents, car les conducteurs roulent vite, sans casque ni permis, et certaines motos servent même à des vols. Les autorités ont bien imposé l'immatriculation et le port du casque, mais ces mesures restent mal appliquées. Plutôt que d'interdire cette activité, il faut donc former les conducteurs et renforcer les contrôles, afin de concilier nécessité économique et sécurité des populations.}

\textbf{Étape 5 — Vérification.} 126 mots : conforme à la fourchette 113--137. Aucune phrase recopiée, aucun « je ». La cohésion est assurée par \emph{d'une part / d'autre part, en outre, car, mais, plutôt que, donc, afin de}. Les quatre parties du plan sont représentées : la fidélité au texte est respectée.
\end{exoresolu}

\begin{attention}[Les trois fautes qui coûtent cher au Probatoire technique]
\begin{itemize}
\item \textbf{Recopier} des phrases du texte : le correcteur les repère immédiatement et sanctionne la copie.
\item \textbf{Oublier une idée essentielle} : le plus souvent la partie « solutions », située à la fin du texte, qu'on néglige par fatigue.
\item \textbf{Dépasser la longueur} imposée : compter les mots est une étape obligatoire, pas une option.
\end{itemize}
\end{attention}

\courssub{Bilan de l'activité}

Cette activité a permis de mobiliser, dans une situation proche de l'examen du Probatoire technique, l'ensemble des savoir-faire de l'unité : repérer la structure argumentative d'un texte explicatif, construire un plan, réduire et reformuler sans trahir, rédiger un texte cohérent dans une longueur imposée, puis évaluer et améliorer une production. La notation croisée a montré que les erreurs les plus fréquentes sont la recopie de phrases du texte, l'oubli d'une idée essentielle (souvent la partie « solutions ») et le dépassement de la longueur autorisée. La réécriture finale et le corrigé collectif au tableau ont enfin permis à chacun de mesurer ses progrès et de retenir qu'un bon résumé est un texte \textbf{fidèle, bref, clair et personnel dans sa formulation}.

\newpage

\coursec{Méthodes et exercices résolus}

\coursec{Découverte : le texte explicatif et la contraction}

\begin{definition}[Texte explicatif]
Un \cle{texte explicatif} a pour fonction d'informer, d'expliquer un phénomène, un processus ou une idée. Il s'appuie sur des faits, des données, des exemples et suit une progression logique (cause $\rightarrow$ conséquence, problème $\rightarrow$ solution, général $\rightarrow$ particulier). Ses marques linguistiques : connecteurs logiques (\emph{d'abord, car, donc, cependant}), temps de l'indicatif (présent de vérité générale, passé composé pour les faits), vocabulaire précis et dénombré.
\end{definition}

\begin{definition}[Contraction de texte (résumé)]
La \cle{contraction de texte} (ou résumé) consiste à restituer les idées essentielles d'un texte source dans un nombre de mots imposé (généralement le quart du texte original), en reformulant avec ses propres mots, sans trahir la pensée de l'auteur, ni ajouter d'idées personnelles, ni recopier des phrases entières.
\end{definition}

\begin{methode}[Repérer la structure argumentative d'un texte]
Pour résumer un texte, il faut d'abord comprendre comment il est construit. La structure argumentative est l'organisation des idées qui soutiennent une thèse.
\begin{enumerate}
\item \textbf{Lire le texte en entier}, sans rien noter, pour saisir le sujet et la thèse de l'auteur (ce qu'il veut démontrer).
\item \textbf{Découper le texte en parties} : repérer les paragraphes et les connecteurs logiques (\emph{d'abord, ensuite, enfin, car, cependant, donc}).
\item \textbf{Dégager l'idée principale de chaque partie} : une idée par paragraphe en général. Souligner la phrase qui la contient.
\item \textbf{Construire le plan du texte} : rédiger chaque idée sous forme d'une phrase courte, dans l'ordre du texte.
\item \textbf{Vérifier la progression} : le texte progresse-t-il de la cause à la conséquence, du problème à la solution, de la thèse à la preuve ?
\end{enumerate}
\textbf{Réflexe} : le plan du texte est le squelette du résumé. Sans plan, pas de résumé fidèle.
\end{methode}

\begin{exoresolu}[Construire le plan d'un texte explicatif]
\textbf{Énoncé.} Lis le texte suivant et établis son plan en trois idées principales.

\emph{« L'eau est une richesse fragile. D'abord, sa répartition est inégale : certaines régions du monde manquent cruellement d'eau potable pendant que d'autres la gaspillent. Ensuite, la pollution des nappes phréatiques et des rivières diminue chaque année les réserves disponibles. Enfin, le changement climatique modifie le cycle des pluies et aggrave les sécheresses. Il est donc urgent de gérer l'eau comme un bien commun à protéger. »}
\tcblower
\textbf{Solution détaillée.}
\begin{enumerate}
\item \textbf{Lecture globale} : le sujet est l'eau ; la thèse de l'auteur est que l'eau est une richesse fragile qu'il faut protéger.
\item \textbf{Repérage des connecteurs} : \emph{d'abord}, \emph{ensuite}, \emph{enfin} annoncent trois arguments ; \emph{donc} introduit la conclusion. La progression est linéaire : trois causes qui justifient une conclusion.
\item \textbf{Dégagement des idées} :
\begin{itemize}
\item Idée 1 : la répartition de l'eau est inégale dans le monde.
\item Idée 2 : la pollution réduit les réserves d'eau disponibles.
\item Idée 3 : le changement climatique aggrave les sécheresses.
\end{itemize}
\item \textbf{Plan du texte} :
\begin{itemize}
\item I. L'eau est inégalement répartie.
\item II. La pollution diminue les réserves.
\item III. Le changement climatique aggrave les sécheresses.
\item Conclusion : il faut gérer l'eau comme un bien commun.
\end{itemize}
\end{enumerate}
\textbf{Conclusion} : le texte suit une progression argumentative en trois étapes (inégalité, pollution, climat) qui aboutissent à une conclusion : la protection de l'eau. Ce plan servira de base au résumé.
\end{exoresolu}

\coursec{La structure argumentative et le plan du texte}

\begin{propriete}[Types de progression dans le texte explicatif]
On distingue trois types principaux de progression :
\begin{itemize}
\item \textbf{Progression linéaire (ou en chaîne)} : chaque idée appelle la suivante (A $\rightarrow$ B $\rightarrow$ C). Exemple : cause $\rightarrow$ conséquence $\rightarrow$ solution.
\item \textbf{Progression parallèle (ou par aspects)} : plusieurs aspects d'un même sujet sont traités successivement sans dépendance forte (A ; B ; C). Exemple : avantages économiques ; avantages sociaux ; avantages écologiques.
\item \textbf{Progression dialectique} : thèse $\rightarrow$ antithèse $\rightarrow$ synthèse. Le texte présente un point de vue, son opposé, puis une réconciliation.
\end{itemize}
\textbf{Au Probatoire} : on vous demandera souvent d'identifier le type de progression pour justifier l'ordre des idées dans votre résumé.
\end{propriete}

\begin{methode}[Établir le plan détaillé d'un texte]
\begin{enumerate}
\item Numéroter les paragraphes.
\item Repérer les \cle{connecteurs logiques} (marqueurs de relation : cause, conséquence, opposition, concession, énumération, conclusion).
\item Formuler l'idée-force de chaque paragraphe en une phrase nominale ou verbale courte.
\item Grouper les paragraphes qui traitent d'une même idée principale (sous-parties).
\item Rédiger le plan en chiffres romains (I, II, III) pour les grandes parties, en lettres (A, B) pour les sous-parties.
\end{enumerate}
\end{methode}

\begin{exoresolu}[Identifier le type de progression]
\textbf{Énoncé.} Quel type de progression suit ce texte ? Justifiez.

\emph{« Le téléphone portable a transformé la vie des Camerounais. Grâce à lui, les commerçants passent leurs commandes sans se déplacer, les familles restent en contact avec leurs proches éloignés, et le transfert d'argent par mobile money facilite les paiements quotidiens. Cependant, cet outil merveilleux, ce petit bijou de technologie, crée aussi une dépendance inquiétante chez les jeunes, qui y passent des heures entières, au détriment de leurs études et de leur sommeil. »}
\tcblower
\textbf{Solution.}
Le texte suit une \textbf{progression dialectique} (thèse / antithèse) :
\begin{itemize}
\item \textbf{Thèse (I)} : le portable facilite la vie quotidienne (trois exemples parallèles : commerce, communication, paiement).
\item \textbf{Antithèse (II)} : mais il crée une dépendance néfaste chez les jeunes (études, sommeil).
\item \textbf{Synthèse implicite} : l'outil est ambivalent, il faut en user avec modération.
\end{itemize}
Le connecteur \emph{cependant} marque le basculement. Ce type de progression impose au résumé de garder l'opposition.
\end{exoresolu}

\coursec{Cohésion, cohérence et structure de la phrase}

\begin{definition}[Cohérence et cohésion]
\begin{itemize}
\item \textbf{Cohérence} : logique sémantique du texte ; les idées s'enchaînent selon une progression rationnelle (cause $\rightarrow$ effet, général $\rightarrow$ particulier, thèse $\rightarrow$ arguments).
\item \textbf{Cohésion} : ensemble des moyens linguistiques qui assurent la liaison entre les énoncés (reprises nominales, pronoms, connecteurs, champs lexicaux, temps verbaux).
\end{itemize}
Un résumé réussi allie les deux : il est \emph{cohérent} (fidèle à la logique de l'auteur) et \emph{cohesif} (lu comme un texte uni, pas une liste de phrases).
\end{definition}

\begin{propriete}[Moyens de cohésion grammaticale et lexicale]
\begin{center}
\begin{tabularx}{\linewidth}{|l|X|X|}\hline
\textbf{Type} & \textbf{Moyens} & \textbf{Exemples} \\ \hline
Reprise nominale & Nom + déterminant démonstratif / possessif & \emph{cette situation, ce phénomène, son origine} \\ \hline
Reprise pronominale & Pronoms personnels, démonstratifs, relatifs & \emph{il, elle, cela, celui-ci, qui, dont} \\ \hline
Substitution lexicale & Synonymes, hyperonymes, mots du champ lexical & \emph{l'eau $\rightarrow$ cette ressource, le liquide vital} \\ \hline
Connecteurs logiques & Conjonctions, adverbes, locutions & \emph{donc, cependant, par conséquent, en outre} \\ \hline
Chaînage tenseur & Maintien ou progression des temps verbaux & Présent de vérité $\rightarrow$ passé composé (faits) \\ \hline
\end{tabularx}
\end{center}
\end{propriete}

\begin{methode}[Assurer la cohésion et la cohérence du résumé]
\begin{enumerate}
\item \textbf{Cohérence} : les idées doivent suivre l'ordre logique du texte (cause avant conséquence, problème avant solution).
\item \textbf{Cohésion grammaticale} : utiliser les reprises nominales (\emph{ce phénomène, cette situation}), les pronoms (\emph{il, elle, cela}) et les démonstratifs pour éviter les répétitions.
\item \textbf{Cohésion lexicale} : employer des synonymes et des mots du même champ lexical pour enchaîner les idées.
\item \textbf{Connecteurs logiques} : relier les phrases par \emph{d'abord, ensuite, enfin, ainsi, cependant, c'est pourquoi}, selon la relation logique du texte.
\item \textbf{Phrase simple ou composée} : varier les constructions mais rester clair ; une phrase trop longue brouille le résumé.
\end{enumerate}
\textbf{Réflexe} : relire le résumé à voix haute ; si l'enchaînement sonne faux, ajouter ou corriger un connecteur.
\end{methode}

\begin{exoresolu}[Améliorer la cohésion d'un résumé]
\textbf{Énoncé.} Le résumé suivant manque de cohésion. Réécris-le en un texte fluide et cohérent.

\emph{« L'eau est inégalement répartie. La pollution diminue les réserves. Le changement climatique aggrave les sécheresses. Il faut protéger l'eau. »}
\tcblower
\textbf{Solution détaillée.}
\begin{enumerate}
\item \textbf{Diagnostic} : quatre phrases simples juxtaposées, sans connecteur, sans reprise ; le lien logique (trois causes $\rightarrow$ une conclusion) n'apparaît pas.
\item \textbf{Ajout des connecteurs} : \emph{d'abord} pour la première cause, \emph{ensuite} pour la deuxième, \emph{enfin} pour la troisième, \emph{c'est pourquoi} pour la conclusion.
\item \textbf{Reprises nominales et pronominales} : reprendre \emph{l'eau} par \emph{cette ressource} ou \emph{elle} pour éviter la répétition.
\item \textbf{Réécriture proposée} : « L'eau est d'abord inégalement répartie dans le monde ; ensuite, la pollution diminue les réserves disponibles ; enfin, le changement climatique aggrave les sécheresses. C'est pourquoi il est urgent de protéger cette ressource. »
\item \textbf{Vérification} : l'ordre des idées est respecté, les relations logiques sont explicites, le texte se lit d'une traite.
\end{enumerate}
\textbf{Conclusion} : grâce aux connecteurs et aux reprises, le résumé est devenu un texte cohérent et cohésif, fidèle à la progression du texte source.
\end{exoresolu}

\begin{propriete}[Structure de la phrase : simple vs composée]
\begin{itemize}
\item \textbf{Phrase simple} : une seule proposition (un seul verbe conjugué). Privilégiée pour la clarté du résumé. Ex. : \emph{La pollution réduit les réserves.}
\item \textbf{Phrase composée} : plusieurs propositions liées par coordination (\emph{mais, ou, et, donc, or, ni, car}) ou subordination (\emph{que, quand, comme, si, bien que}). Utile pour marquer les relations logiques (cause, conséquence, opposition). Ex. : \emph{Parce que l'eau est polluée, les réserves diminuent.}
\end{itemize}
\textbf{Conseil} : dans un résumé court, alternez phrases simples (pour les idées fortes) et phrases composées courtes (pour les liens logiques). Évitez les phrases à plus de deux propositions.
\end{propriete}

\coursec{Techniques de réduction et de reformulation}

\begin{methode}[Réduire un texte : les techniques de contraction]
Le résumé doit conserver les idées essentielles en respectant le quart du texte environ (ou la limite imposée par le sujet). Trois techniques de réduction :
\begin{enumerate}
\item \textbf{La suppression} : éliminer les exemples, les répétitions, les citations, les détails, les figures de style, les énumérations longues.
\item \textbf{La généralisation} : remplacer une énumération par un terme générique (\emph{mangues, bananes, ananas} $\rightarrow$ \emph{fruits} ; \emph{Douala, Yaoundé, Bafoussam} $\rightarrow$ \emph{les grandes villes camerounaises}).
\item \textbf{La condensation} : remplacer plusieurs mots par un seul (\emph{celui qui enseigne} $\rightarrow$ \emph{l'enseignant} ; \emph{il ne faut pas jeter les ordures n'importe où} $\rightarrow$ \emph{il faut éviter le dépôt sauvage} ; \emph{le fait que les jeunes n'aient pas d'emploi} $\rightarrow$ \emph{le chômage des jeunes}).
\end{enumerate}
\textbf{Réflexes} : ne jamais recopier de phrases entières du texte ; compter les mots à chaque étape ; conserver les mots-clés (noms propres, termes techniques, notions centrales, connecteurs logiques porteurs de sens).
\end{methode}

\begin{exoresolu}[Réduire un paragraphe]
\textbf{Énoncé.} Réduis le passage suivant (52 mots) à environ 13 mots, sans recopier de phrase du texte.

\emph{« Dans les villes camerounaises comme Douala et Yaoundé, on voit partout, dans les rues, sur les trottoirs et même dans les caniveaux, des sachets plastiques, des bouteilles vides et des papiers sales qui s'accumulent jour après jour. »}
\tcblower
\textbf{Solution détaillée.}
\begin{enumerate}
\item \textbf{Idée essentielle} : les déchets s'accumulent dans les villes camerounaises.
\item \textbf{Suppressions} : les exemples de villes (\emph{Douala et Yaoundé}), l'énumération des lieux (\emph{rues, trottoirs, caniveaux}), l'énumération des déchets (\emph{sachets, bouteilles, papiers}), la précision \emph{jour après jour}.
\item \textbf{Généralisation} : \emph{sachets plastiques, bouteilles vides, papiers sales} $\rightarrow$ \emph{les déchets}.
\item \textbf{Condensation} : \emph{on voit partout... qui s'accumulent} $\rightarrow$ \emph{s'accumulent}.
\item \textbf{Résumé proposé} : « Les déchets s'accumulent dans les villes camerounaises. » (8 mots, conforme à la limite).
\end{enumerate}
\textbf{Conclusion} : le passage est réduit à son idée essentielle, sans phrase recopiée ni détail inutile.
\end{exoresolu}

\begin{methode}[Reformuler les idées essentielles]
Reformuler, c'est dire la même chose avec ses propres mots, sans trahir la pensée de l'auteur.
\begin{enumerate}
\item \textbf{Remplacer les mots par des synonymes} : \emph{fragile} $\rightarrow$ \emph{menacée} ; \emph{urgent} $\rightarrow$ \emph{nécessaire rapidement} ; \emph{gaspillent} $\rightarrow$ \emph{dilapident}.
\item \textbf{Changer la construction de la phrase} : transformer une phrase passive en phrase active, une phrase composée en phrase simple (\emph{Comme il pleuvait, le match fut reporté} $\rightarrow$ \emph{La pluie fit reporter le match} ; \emph{Les jeunes qui sont au chômage quittent le village} $\rightarrow$ \emph{Le chômage pousse les jeunes à quitter le village}).
\item \textbf{Conserver les mots-clés irremplaçables} : noms propres, termes techniques, notions centrales (ex. \emph{mobile money, nappes phréatiques, exode rural}).
\item \textbf{Garder les relations logiques} : si le texte dit \emph{parce que}, le résumé doit garder la cause ; s'il dit \emph{mais}, garder l'opposition.
\item \textbf{Vérifier la fidélité} : relire le texte et le résumé côte à côte ; aucune idée nouvelle ne doit apparaître, aucune idée essentielle ne doit disparaître.
\end{enumerate}
\textbf{À retenir} : reformuler $\neq$ paraphraser mot à mot ; reformuler $\neq$ commenter ou donner son avis.
\end{methode}

\begin{exoresolu}[Reformuler une idée sans la trahir]
\textbf{Énoncé.} Reformule la phrase suivante en changeant les mots et la construction, sans en changer le sens.

\emph{« Parce que la jeunesse manque d'emplois, beaucoup de jeunes diplômés quittent les villages pour tenter leur chance dans les grandes villes. »}
\tcblower
\textbf{Solution détaillée.}
\begin{enumerate}
\item \textbf{Analyse de la phrase} : phrase composée avec une subordonnée de cause (\emph{parce que...}) et une principale. Idées : (a) chômage des jeunes, (b) exode rural des diplômés.
\item \textbf{Synonymes} : \emph{manque d'emplois} $\rightarrow$ \emph{chômage} ; \emph{quittent les villages} $\rightarrow$ \emph{abandonnent les campagnes} ; \emph{tenter leur chance} $\rightarrow$ \emph{chercher du travail}.
\item \textbf{Changement de construction} : transformer la subordonnée de cause en complément de nom (\emph{en raison de}, \emph{faute de}).
\item \textbf{Reformulation proposée} : « Faute d'emplois, de nombreux jeunes diplômés abandonnent les campagnes pour chercher du travail en ville. »
\item \textbf{Vérification} : la cause (chômage) et la conséquence (exode rural) sont conservées ; aucune idée ajoutée ; aucune phrase recopiée.
\end{enumerate}
\textbf{Conclusion} : la reformulation est fidèle au sens du texte tout en étant rédigée dans une langue personnelle.
\end{exoresolu}

\coursec{Les figures d'atténuation : litote, euphémisme, prétérition}

\begin{definition}[Figures d'atténuation]
Les \cle{figures d'atténuation} (ou figures de sous-énonciation) permettent à l'énonciateur de dire moins fort, d'adoucir ou de voiler son propos, tout en laissant deviner le sens réel. Le résumeur doit les identifier pour ne pas durcir ni trahir la pensée de l'auteur.
\end{definition}

\begin{propriete}[Les trois figures d'atténuation principales]
\begin{center}
\begin{tabularx}{\linewidth}{|l|X|X|X|}\hline
\textbf{Figure} & \textbf{Définition} & \textbf{Exemple} & \textbf{Sens réel} \\ \hline
\textbf{Litote} & Dire moins pour faire entendre plus ; affirmation par une négation atténuée. & \emph{« Ce n'est pas bête »} ; \emph{« Il n'est pas inintelligent »} ; \emph{« Va, je ne te hais point »} (Corneille) & C'est intelligent ; Je t'aime. \\ \hline
\textbf{Euphémisme} & Adoucir une réalité brutale, taboue ou triste par un terme plus doux. & \emph{« Il s'est éteint »} ; \emph{« Il nous a quittés »} ; \emph{« Personne âgée »} (au lieu de vieux) ; \emph{« Malentendant »} (au lieu de sourd) & Il est mort ; Il est vieux ; Il est sourd. \\ \hline
\textbf{Prétérition} & Faire semblant de ne pas dire ce qu'on dit en le signalant. & \emph{« Je ne parlerai pas de ses erreurs... »} ; \emph{« Passons sous silence ses dettes. »} & Je dénonce ses erreurs ; Je signale ses dettes. \\ \hline
\end{tabularx}
\end{center}
\end{propriete}

\begin{methode}[Identifier et exploiter les figures d'atténuation dans un résumé]
\begin{enumerate}
\item \textbf{Repérer} la tournure atténuée (négation litotique, terme euphémistique, annonce de silence).
\item \textbf{Interpréter} : « Que veut dire réellement l'auteur derrière cette atténuation ? »
\item \textbf{Restituer} le sens réel dans le résumé, sans exagérer la force de l'énoncé si le ton compte pour la compréhension (ex. : garder une nuance si l'auteur use d'ironie).
\item \textbf{Ne pas inventer} : ne pas transformer une litote en insultes, un euphémisme en terme cru, une prétérition en accusation non fondée.
\end{enumerate}
\textbf{Réflexe} : demander toujours « que veut dire réellement l'auteur ? » avant de reformuler.
\end{methode}

\begin{exoresolu}[Reconnaître les figures d'atténuation]
\textbf{Énoncé.} Identifie la figure d'atténuation employée dans chaque phrase et explique ce qu'elle signifie réellement.

a) \emph{« Ce candidat n'est pas un génie, mais il travaille sérieusement. »}

b) \emph{« Le vieux chef du village nous a quittés cette nuit. »}

c) \emph{« Je passerai sous silence les erreurs de gestion du maire. »}
\tcblower
\textbf{Solution détaillée.}
\begin{enumerate}
\item \textbf{Phrase a)} : la tournure négative \emph{n'est pas un génie} dit moins pour suggérer plus : c'est une \textbf{litote}. Sens réel : le candidat est plutôt médiocre intellectuellement, mais sérieux.
\item \textbf{Phrase b)} : \emph{nous a quittés} adoucit la réalité de la mort : c'est un \textbf{euphémisme}. Sens réel : le chef du village est mort.
\item \textbf{Phrase c)} : l'énonciateur déclare ne pas parler des erreurs tout en les signalant : c'est une \textbf{prétérition}. Sens réel : le maire a commis des erreurs de gestion, et l'orateur veut qu'on le sache.
\item \textbf{Conséquence pour le résumé} : dans un résumé, on rendrait (a) par « le candidat est médiocre mais travailleur », (b) par « le chef est mort », (c) par « l'orateur dénonce les erreurs du maire », en veillant à ne pas inventer de jugement que le texte ne contient pas.
\end{enumerate}
\textbf{Conclusion} : litote, euphémisme et prétérition atténuent la forme mais non le fond ; le résumeur doit restituer le sens réel sans exagérer.
\end{exoresolu}

\begin{exoresolu}[Figures d'atténuation dans un texte à résumer]
\textbf{Énoncé.} Le texte suivant contient des figures d'atténuation. Identifiez-les et proposez une reformulation pour le résumé.

\emph{« Le directeur n'est pas incompétent, loin de là. Il a simplement omis de mentionner, je passe sous silence les détails, que les comptes ne sont pas tout à fait équilibrés. Le trésorier, lui, a rejoint ses ancêtres l'an dernier. »}
\tcblower
\textbf{Solution.}
\begin{itemize}
\item \emph{n'est pas incompétent} $\rightarrow$ \textbf{litote} (il est compétent).
\item \emph{omission de mentionner / je passe sous silence} $\rightarrow$ \textbf{prétérition} (il signale volontairement l'omission).
\item \emph{pas tout à fait équilibrés} $\rightarrow$ \textbf{litote} (ils sont en déficit).
\item \emph{a rejoint ses ancêtres} $\rightarrow$ \textbf{euphémisme} (il est mort).
\end{itemize}
\textbf{Reformulation pour résumé} : « Le directeur est compétent mais a tu le déficit des comptes ; le trésorier est décédé l'an dernier. »
\end{exoresolu}

\coursec{Rédaction, confrontation et amélioration du résumé}

\begin{methode}[Rédiger et justifier une démarche, une réponse ou une conclusion]
Rédiger le résumé est une démarche en cinq temps, qu'il faut savoir justifier (au compte rendu ou à l'oral).
\begin{enumerate}
\item \textbf{Brouillon} : rédiger un premier jet à partir du plan, sans se soucier du nombre de mots.
\item \textbf{Comptage} : compter les mots et réduire jusqu'à la limite imposée (souvent le quart du texte).
\item \textbf{Confrontation} : comparer son résumé avec celui d'un camarade ou avec le texte ; noter les écarts (idées manquantes, idées en trop, phrases recopiées, connecteurs absents).
\item \textbf{Amélioration} : corriger la langue (orthographe, accords, ponctuation), renforcer la cohésion, supprimer les restes de recopie, vérifier la fidélité.
\item \textbf{Justification} : être capable d'expliquer ses choix : « J'ai supprimé les exemples car ils illustrent sans ajouter d'idée ; j'ai généralisé l'énumération ; j'ai gardé le connecteur \emph{donc} car il porte la conclusion. »
\end{enumerate}
\textbf{Formule à retenir} : un bon résumé est \textbf{fidèle} (mêmes idées), \textbf{bref} (limite respectée), \textbf{personnel} (aucune phrase recopiée), \textbf{cohérent} (texte fluide).
\end{methode}

\begin{exoresolu}[Rédiger et justifier un résumé complet]
\textbf{Énoncé.} Résume le texte suivant (environ 80 mots) en 20 mots environ, puis justifie deux de tes choix de réduction.

\emph{« Le téléphone portable a transformé la vie des Camerounais. Grâce à lui, les commerçants passent leurs commandes sans se déplacer, les familles restent en contact avec leurs proches éloignés, et le transfert d'argent par mobile money facilite les paiements quotidiens. Cependant, cet outil merveilleux, ce petit bijou de technologie, crée aussi une dépendance inquiétante chez les jeunes, qui y passent des heures entières, au détriment de leurs études et de leur sommeil. »}
\tcblower
\textbf{Solution détaillée.}
\begin{enumerate}
\item \textbf{Plan du texte} : I. Le portable facilite la vie (commerce, communication, paiements). II. Mais il crée une dépendance néfaste chez les jeunes.
\item \textbf{Suppressions} : les trois exemples détaillés (commandes, proches, mobile money) ; les figures de style (\emph{outil merveilleux, petit bijou}) ; la précision \emph{des heures entières}.
\item \textbf{Généralisation} : les trois exemples $\rightarrow$ \emph{les activités quotidiennes}.
\item \textbf{Résumé proposé} : « Le téléphone portable facilite la vie quotidienne des Camerounais, mais il rend les jeunes dépendants, au détriment de leurs études. » (19 mots, limite respectée).
\item \textbf{Justification des choix} :
\begin{itemize}
\item J'ai supprimé l'énumération des usages car elle illustre une seule idée (la facilitation), que j'ai généralisée par « facilite la vie quotidienne ».
\item J'ai conservé le connecteur d'opposition \emph{mais} car il porte la structure argumentative du texte (avantages \emph{contre} inconvénients) ; le supprimer aurait trahi la pensée de l'auteur.
\end{itemize}
\end{enumerate}
\textbf{Conclusion} : le résumé est fidèle (les deux idées et l'opposition sont gardées), bref (19 mots), personnel (aucune phrase recopiée) et cohérent.
\end{exoresolu}

\begin{methode}[Confrontation et amélioration par les pairs]
En classe, échangez vos résumés et appliquez cette grille d'analyse :
\begin{center}
\begin{tabularx}{\linewidth}{|l|X|}\hline
\textbf{Critère} & \textbf{Questions à se poser} \\ \hline
Fidélité & Toutes les idées essentielles sont-elles présentes ? Aucune idée ajoutée ? \\ \hline
Concision & Le nombre de mots est-il respecté ? Y a-t-il des phrases recopiées ? \\ \hline
Reformulation & Les idées sont-elles exprimées avec ses propres mots (synonymes, changement de structure) ? \\ \hline
Cohérence & L'ordre des idées suit-il la progression du texte source ? \\ \hline
Cohésion & Les phrases sont-elles reliées par des connecteurs et des reprises ? \\ \hline
Langue & Orthographe, accords, ponctuation, syntaxe sont-ils corrects ? \\ \hline
\end{tabularx}
\end{center}
Notez les points forts et les points à améliorer ; réécrivez une version finale corrigée.
\end{methode}

\begin{attention}[Les 5 erreurs qui coûtent des points au Bac / Probatoire]
\begin{enumerate}
\item \textbf{Recopier des phrases du texte.} Un résumé avec des phrases reprises mot pour mot est sanctionné : reformule toujours avec tes propres mots, en ne gardant que les mots-clés irremplaçables.
\item \textbf{Dépasser la limite de mots imposée.} Compte tes mots et annonce le nombre sous le résumé ; un dépassement important fait perdre des points. Vise légèrement en dessous de la limite.
\item \textbf{Donner son avis ou ajouter des idées.} Le résumé restitue la pensée de l'auteur, pas la tienne : aucune appréciation (\emph{à mon avis, heureusement, malheureusement}), aucun exemple nouveau.
\item \textbf{Perdre les relations logiques.} Supprimer les connecteurs (\emph{car, mais, donc}) détruit la structure argumentative : garde ou reformule les liens de cause, d'opposition et de conséquence.
\item \textbf{Négliger la langue et la cohésion.} Un résumé en phrases télégraphiques, avec fautes d'accord et sans reprises, est mal noté : rédige un texte fluide, relie les phrases, relis-toi pour corriger orthographe et ponctuation.
\end{enumerate}
\end{attention}

\coursec{Ouverture littéraire : Le Lion et la perle — bilan de l'œuvre}

\begin{savaistu}[Wole Soyinka et \emph{Le Lion et la perle}]
\emph{Le Lion et la perle} (\emph{The Lion and the Jewel}, 1959) est une comédie de l'écrivain nigérian \textbf{Wole Soyinka}, prix Nobel de littérature 1986. La pièce se déroule à \textbf{Ilujinle}, village yoruba imaginaire, et met en scène le conflit entre tradition et modernité à travers trois personnages principaux :
\begin{itemize}
\item \textbf{Baroka} (le Bale, « le Lion ») : chef traditionnel, rusé, polygame, incarnation de la tradition africaine vivante et adaptable.
\item \textbf{Lakunle} (l'instituteur) : jeune homme « moderne », occidentalisé, refusant la dot, parlant un anglais ampoulé, caricature de la modernité superficielle.
\item \textbf{Sidi} (la « Perle ») : belle villageoise, convoitée par les deux hommes, symbole de la jeunesse africaine tiraillée entre les deux modèles.
\end{itemize}
\end{savaistu}

\begin{propriete}[Thèmes majeurs de la pièce]
\begin{itemize}
\item \textbf{Tradition vs Modernité} : non pas un choc frontal mais une négociation ; Baroka utilise la modernité (timbre-poste, machine) pour asseoir son pouvoir traditionnel.
\item \textbf{La condition féminine} : Sidi n'est pas un objet passif ; elle use de sa beauté comme pouvoir, mais finit piégée par les codes patriarcaux (dot, polygamie).
\item \textbf{Le pouvoir et la ruse} : Baroka triomphe par l'intelligence et la connaissance des hommes, pas par la force.
\item \textbf{Le langage et l'image} : la photographie (le magazine) devient un enjeu de pouvoir ; le langage de Lakunle (anglais prétentieux) contraste avec la verve proverbiale de Baroka.
\item \textbf{Le théâtre total} : chants, danses, mimes, proverbes — Soyinka intègre les formes traditionnelles yoruba (egungun, gelede) dans la structure dramatique occidentale.
\end{itemize}
\end{propriete}

\begin{methode}[Bilan de l'étude de l'œuvre : points à maîtriser pour l'évaluation]
\begin{enumerate}
\item \textbf{Contexte} : situer l'auteur (Nigéria, Yoruba, indépendance 1960), le genre (comédie, théâtre total), la date (1959, pré-indépendance).
\item \textbf{Personnages} : caractériser Baroka, Lakunle, Sidi, Sadiku ; expliquer leur fonction symbolique.
\item \textbf{Structure} : trois parties (Matin, Midi, Soir) = temps de la journée et temps de la séduction ; unité de lieu (Ilujinle).
\item \textbf{Scènes clés} : la danse du voyageur perdu (mime/théâtre dans le théâtre), la scène du timbre-poste (modernité détournée), la scène finale (triomphe de Baroka, résignation de Sidi).
\item \textbf{Style} : proverbes yoruba, images animalières (lion, perle, renard), ironie, humour, mélange anglais/yoruba.
\item \textbf{Portée} : la pièce ne dit pas « tradition bonne / modernité mauvaise » ; elle montre une Afrique qui négocie son avenir.
\end{enumerate}
\end{methode}

\begin{exoresolu}[Question type Probatoire sur l'œuvre]
\textbf{Énoncé.} « Dans \emph{Le Lion et la perle}, la tradition n'est pas figée : elle sait user de la modernité pour se perpétuer. » Discutez cette affirmation en vous appuyant sur le personnage de Baroka.
\tcblower
\textbf{Proposition de réponse structurée.}
\begin{enumerate}
\item \textbf{Introduction} : présenter l'œuvre, l'auteur, le personnage de Baroka (Bale d'Ilujinle), l'affirmation à discuter.
\item \textbf{Argument 1 : Baroka incarne la tradition vivante.} Il respecte les coutumes (polygamie, dot, proverbes, danse), il est le gardien de l'ordre social.
\item \textbf{Argument 2 : Il détourne la modernité à son profit.} Exemple du timbre-poste : il veut effigier Sidi pour « moderniser » le village, mais c'est un piège pour la séduire. Il utilise la machine (imprimerie) pour asseoir son prestige personnel.
\item \textbf{Argument 3 : Sa ruse est une intelligence adaptative.} Il feint l'impuissance pour tester Sidi ; il connaît la psychologie de Lakunle et de Sidi. La tradition n'est pas rigide, elle est stratège.
\item \textbf{Nuance} : cette adaptation a un coût : Sidi est manipulée, la condition féminine n'évolue pas (dot, polygamie maintenues). La modernité n'apporte pas l'émancipation attendue.
\item \textbf{Conclusion} : Baroka prouve que la tradition yoruba est dynamique, capable d'intégrer le nouveau sans perdre son essence, mais au prix d'une conservation des rapports de pouvoir patriarchaux.
\end{enumerate}
\end{exoresolu}

\begin{exercice}[Entraînement au résumé — Sujet type Probatoire]
\textbf{Texte support (environ 100 mots) :}

\emph{« L'agriculture camerounaise fait face à de multiples défis. D'une part, les exploitations familiales, qui représentent plus de 80\% de la production, manquent d'intrants améliorés, de matériel adapté et d'accès au crédit. D'autre part, les pistes rurales sont souvent impraticables en saison des pluies, ce qui isole les zones de production des marchés de consommation. Par ailleurs, le changement climatique perturbe les calendriers culturales : les pluies arrivent en retard ou s'arrêtent trop tôt, provoquant des pertes de récoltes. Face à cette situation, l'État a lancé le Programme d'Appui à la Compétitivité de l'Agriculture Familiale (PACA) pour financer des infrastructures et former les producteurs. »}

\textbf{Consigne :} Résumez ce texte en 25 mots maximum. Indiquez le nombre de mots de votre résumé. Justifiez brièvement deux choix de réduction opérés.
\end{exercice}

\begin{corrige}[Exercice n°1 — Proposition de correction]
\textbf{Résumé proposé (24 mots) :}
« L'agriculture familiale camerounaise, freinée par le manque de moyens, l'enclavement et le climat, bénéficie du programme PACA pour moderniser les exploitations. »

\textbf{Justification de deux choix :}
\begin{enumerate}
\item \textbf{Suppression du chiffre « 80\% » et de l'énumération des manques (intrants, matériel, crédit)} : généralisés par « manque de moyens » pour gagner de la place sans perdre l'idée essentielle.
\item \textbf{Condensation de la cause climatique} : « le changement climatique perturbe les calendriers culturales... » $\rightarrow$ « le climat » (hyperonyme) ; la conséquence (pertes de récoltes) est implicite dans « freinée ».
\item \textbf{Conservation de « PACA »} : sigle technique irremplaçable, nom propre du dispositif étatique.
\item \textbf{Connecteur « Face à cette situation » $\rightarrow$ « pour »} : marque la finalité de l'action étatique (solution aux problèmes énoncés).
\end{enumerate}
\end{corrige}

\begin{exercice}[Entraînement figures d'atténuation]
Identifiez la figure d'atténuation dans chaque phrase et donnez le sens réel pour un résumé.

1. \emph{« Ce n'est pas un mauvais élève, simplement un peu paresseux. »}

2. \emph{« Le ministre a dû s'absenter pour raisons de santé ; on dit qu'il a rejoint les siens. »}

3. \emph{« Je ne mentionnerai pas les détournements de fonds dont on parle dans la presse. »}
\end{exercice}

\begin{corrige}[Exercice n°2 — Correction]
\begin{enumerate}
\item \textbf{Litote} (\emph{pas un mauvais élève}) $\rightarrow$ Sens : c'est un élève moyen/médiocre. Résumé : « L'élève est moyen, freiné par la paresse. »
\item \textbf{Euphémisme} (\emph{a rejoint les siens}) $\rightarrow$ Sens : il est mort. Résumé : « Le ministre est décédé. »
\item \textbf{Prétérition} (\emph{je ne mentionnerai pas...}) $\rightarrow$ Sens : l'orateur dénonce les détournements. Résumé : « L'orateur signale des détournements de fonds présumés. »
\end{enumerate}
\end{corrige}

\newpage

\coursec{Exercices}

\courssub{Je vérifie mes connaissances}

\begin{exercice}[QCM : la contraction de texte]
Pour chaque question, une seule réponse est correcte. Entoure la lettre correspondante.
\begin{enumerate}
\item Un résumé de texte doit contenir :
\begin{enumerate}
\item[a)] les idées essentielles de l'auteur, sans opinion personnelle ;
\item[b)] les idées essentielles et mon avis sur le sujet ;
\item[c)] tous les exemples du texte de départ.
\end{enumerate}
\item Si un texte compte 480 mots et que la consigne demande un résumé au quart, la longueur attendue est :
\begin{enumerate}
\item[a)] 240 mots ; \item[b)] 120 mots ; \item[c)] 60 mots.
\end{enumerate}
\item La technique qui consiste à remplacer « le téléphone portable, la tablette et l'ordinateur » par « les outils numériques » s'appelle :
\begin{enumerate}
\item[a)] la suppression ; \item[b)] la généralisation ; \item[c)] la nominalisation.
\end{enumerate}
\item Une phrase simple est une phrase qui contient :
\begin{enumerate}
\item[a)] plusieurs verbes conjugués ; \item[b)] un seul verbe conjugué ; \item[c)] aucun verbe.
\end{enumerate}
\end{enumerate}
\end{exercice}

\begin{exercice}[Vrai ou faux ? Justifie ta réponse]
Réponds par \emph{vrai} ou \emph{faux}, puis justifie en une ou deux phrases.
\begin{enumerate}
\item Dans un résumé, on peut conserver les exemples de l'auteur pour rendre le texte plus vivant.
\item La cohésion d'un texte repose notamment sur les reprises nominales, les pronoms et les connecteurs logiques.
\item Une progression thématique reprend le même thème du début à la fin du texte.
\item La litote est une figure qui consiste à exagérer une idée pour frapper l'esprit du lecteur.
\end{enumerate}
\end{exercice}

\begin{exercice}[Définitions]
Définis en deux ou trois phrases chacun des termes suivants, puis illustre chacun par un exemple de ton invention :
\begin{enumerate}
\item la contraction de texte ;
\item la thèse d'un texte argumentatif ;
\item la cohérence textuelle ;
\item l'euphémisme.
\end{enumerate}
\end{exercice}

\begin{exercice}[Calculs directs]
\begin{enumerate}
\item Un texte compte 600 mots. Calcule la longueur d'un résumé au quart, puis celle d'un résumé au tiers.
\item La consigne impose un résumé de 120 mots, avec une marge de tolérance de 10 \%. Calcule le nombre minimal et le nombre maximal de mots acceptés.
\item Un élève a rédigé un résumé de 180 mots pour un texte de 450 mots. Quelle fraction du texte représente son résumé ? Cette longueur est-elle conforme à une consigne « au quart » ?
\end{enumerate}
\end{exercice}

\courssub{Je m'entraîne}

\begin{exercice}[Repérage dans un éditorial]
Lis attentivement l'extrait suivant (environ 300 mots) :

\emph{« Dans nos villes comme dans nos villages, le téléphone portable s'est imposé comme un compagnon quotidien des élèves. Faut-il pour autant l'autoriser en classe ? Je réponds non, et je l'affirme sans détour. D'abord, le portable détourne l'attention : combien d'élèves, au lieu d'écouter l'explication du professeur, consultent leurs messages sous la table ? Un enseignant de Yaoundé raconte avoir confisqué douze appareils en une seule journée. Ensuite, cet appareil favorise la triche : pendant les évaluations, des réponses circulent par messagerie, comme l'a révélé une enquête menée dans un lycée de Douala en 2024. Enfin, l'usage excessif du portable nuit à la santé des adolescents : troubles du sommeil, fatigue visuelle, isolement. Certes, dira-t-on, le téléphone peut servir à chercher des informations. Mais cet avantage est mince comparé aux dégâts observés. L'école doit rester un lieu de concentration et d'échange humain : bannissons le portable de la salle de classe. »}

\begin{enumerate}
\item Relève la thèse défendue par l'auteur et recopie la phrase qui l'exprime.
\item Identifie les trois arguments développés et les deux exemples qui les illustrent.
\item Établis le plan du texte en trois ou quatre parties.
\item Précise le type de progression du texte (thématique, argumentative ou autre) et justifie ta réponse.
\end{enumerate}
\end{exercice}

\begin{exercice}[Réduction guidée]
Réduis chacune des phrases suivantes d'environ la moitié, en appliquant à chaque fois la technique indiquée entre parenthèses. Nomme ensuite la technique utilisée.
\begin{enumerate}
\item « Les élèves qui utilisent leur téléphone portable pendant les cours ne suivent pas attentivement les explications du professeur. » (suppression)
\item « Les motos, les voitures, les bus et les camions encombrent les rues de nos villes. » (généralisation)
\item « Il faut que les jeunes prennent conscience du fait que la lecture est une activité qui enrichit l'esprit et qui développe l'imagination. » (condensation)
\item « Le directeur de l'école a annoncé aux parents d'élèves que l'établissement fermerait ses portes pendant toute la durée des travaux de rénovation. » (substitution par un pronom)
\item « Les élèves travaillent avec sérieux, et cela entraîne la réussite aux examens. » (nominalisation)
\end{enumerate}
\end{exercice}

\begin{exercice}[Phrases simples, phrases composées et connecteurs]
\textbf{A.} Transforme les phrases composées suivantes en phrases simples, puis les phrases simples en phrases composées.
\begin{enumerate}
\item « Sidi entra dans le village et tous les regards se tournèrent vers elle. » (à transformer en phrase simple)
\item « Baroka, bien qu'il soit âgé, demeure redoutable. » (à transformer en phrase simple)
\item « Le photographe publia l'image de Sidi. » (à développer en phrase composée)
\item « Lakunle refusa de payer la dot. » (à développer en phrase composée)
\item « Le Bale trompa Sadiku, ce qui provoqua la colère de Sidi. » (à transformer en phrase simple)
\end{enumerate}
\textbf{B.} Classe les dix connecteurs suivants selon leur fonction logique (addition, opposition, cause, conséquence, concession) : \emph{car ; cependant ; de plus ; par conséquent ; bien que ; en effet ; ensuite ; néanmoins ; c'est pourquoi ; en outre}.
\end{exercice}

\begin{exercice}[Figures d'atténuation]
Pour chacune des phrases suivantes, identifie la figure d'atténuation employée (litote, euphémisme ou prétérition) et explique son effet sur le lecteur.
\begin{enumerate}
\item « Ce n'est pas un imbécile, ce Lakunle : il sait manier les mots. »
\item « Le vieux Bale a rejoint ses ancêtres cette nuit. »
\item « Je ne te parle pas de ses dettes, ni de ses querelles avec les voisins. »
\item « Va, tu n'es pas sans savoir que la tradition compte ici. »
\item « Notre ami est un peu fatigué ce matin », dit-il du malade en très mauvais état.
\item « Je passe sous silence les erreurs de jeunesse du notable. »
\item « Ce repas n'était pas mauvais. »
\item « Il nous a quittés trop tôt », annonça le crieur public.
\item « Ce n'est pas tous les jours qu'un étranger honore notre village. »
\item « Je ne mentionnerai pas le prix exorbitant de cette cérémonie. »
\end{enumerate}
\end{exercice}

\courssub{Je me prépare au Bac}

\begin{exercice}[Contraction de texte type Baccalauréat technique]
Lis le texte suivant (environ 480 mots), puis réponds aux questions.

\emph{« Le téléphone portable est aujourd'hui présent dans toutes les poches, y compris celles des élèves camerounais. Cet outil, utile pour communiquer et s'informer, pose pourtant de sérieux problèmes dans le milieu scolaire. Il convient donc d'en encadrer strictement l'usage chez les élèves.}

\emph{D'abord, le téléphone portable constitue une source permanente de distraction. Pendant les cours, beaucoup d'élèves consultent leurs réseaux sociaux, échangent des messages ou jouent à des jeux vidéo au lieu de suivre les explications. Un professeur de mathématiques d'un lycée de Bafoussam affirme que la moitié de ses élèves possède un téléphone allumé en classe. Cette distraction se traduit par une baisse évidente des résultats scolaires : les élèves absorbés par leur écran assimilent mal les leçons et réussissent moins bien aux évaluations.}

\emph{Ensuite, le portable favorise des comportements contraires à l'honnêteté. Lors des compositions, certains élèves reçoivent des réponses par messagerie ou photographient les sujets pour les envoyer à des complices. Plusieurs cas de fraude liés au téléphone ont été signalés lors des examens officiels ces dernières années. De plus, l'appareil peut servir à diffuser des images humiliantes de camarades, ce qui alimente le harcèlement entre élèves.}

\emph{Enfin, l'usage excessif du téléphone nuit à la santé et à la vie sociale des adolescents. Les jeunes qui passent leurs nuits sur leur écran souffrent de fatigue et de troubles du sommeil. Rivés à leur appareil, ils délaissent le sport, la lecture et les conversations en famille. Un médecin scolaire de Douala observe que de nombreux élèves se plaignent de maux de tête et de douleurs aux yeux.}

\emph{Certes, le téléphone portable présente des avantages : il permet de contacter rapidement les parents en cas d'urgence et d'accéder à des ressources éducatives. Toutefois, ces avantages peuvent être préservés sans laisser l'appareil envahir la salle de classe. La solution consiste à interdire le portable pendant les heures de cours, tout en sensibilisant les élèves à un usage responsable. Parents, enseignants et élèves doivent s'entendre sur des règles claires. L'école forme l'intelligence ; elle ne saurait devenir une annexe des réseaux sociaux. »}

\begin{enumerate}
\item Dégage la thèse de l'auteur et les trois arguments principaux qui la soutiennent.
\item Établis le plan du texte en indiquant la fonction de chaque partie.
\item Rédige un résumé de ce texte en 120 mots (plus ou moins 10 \%), en respectant la neutralité (aucune opinion personnelle), la cohésion (connecteurs et reprises) et la reformulation des idées essentielles.
\item Indique à la fin de ton résumé le nombre exact de mots utilisés.
\end{enumerate}
\end{exercice}

\begin{exercice}[Remédiation : corriger un résumé fautif]
Voici le résumé rédigé par un élève à partir du texte sur le téléphone portable (consigne : 120 mots au maximum, neutralité exigée). Ce résumé compte 165 mots.

\emph{« Moi, je trouve que ce texte a totalement raison : les téléphones portables sont un vrai fléau ! L'auteur explique que le portable distrait les élèves en classe. Par exemple, un professeur de mathématiques d'un lycée de Bafoussam affirme que la moitié de ses élèves possède un téléphone allumé en classe. Ensuite, il dit que le portable favorise la triche pendant les compositions, avec des réponses envoyées par messagerie. Il parle aussi du harcèlement. Enfin, le téléphone abîme la santé : un médecin scolaire de Douala observe des maux de tête chez les élèves. Personnellement, je connais même un camarade qui a échoué à cause de son téléphone. L'auteur reconnaît quand même que le portable permet de contacter les parents en urgence. En conclusion, je pense que tous les téléphones devraient être interdits partout au Cameroun, et même confisqués définitivement. »}

\begin{enumerate}
\item Releve les cinq erreurs commises par cet élève (opinion personnelle, exemples conservés, longueur non respectée, conclusion inventée, reprise mot à mot, etc.) et justifie chaque relevé.
\item Pour chaque erreur, propose une correction précise.
\item Rédige une version corrigée du résumé, conforme à la consigne initiale.
\end{enumerate}
\end{exercice}

\begin{exercice}[Question de synthèse sur \emph{Le Lion et la perle}]
Sujet : \textbf{« Baroka est-il le vainqueur ou le fossoyeur de la tradition ? »}

En t'appuyant sur ta connaissance de l'œuvre de Wole Soyinka et sur le contexte dans lequel elle s'inscrit (confrontation entre tradition africaine et modernité au Nigeria), réponds aux questions suivantes :
\begin{enumerate}
\item Rappelle en quelques lignes qui est Baroka et quelle place il occupe dans le village d'Ilujinlè.
\item Présente deux arguments montrant que Baroka défend et incarne la tradition (ruses, statut de Bale, rapport à la dot et à la polygamie).
\item Présente deux arguments montrant que Baroka, par ses manœuvres, trahit ou fragilise cette même tradition (mensonge sur son impuissance, instrumentalisation de Sadiku).
\item Rédige un développement construit d'environ vingt lignes, avec une introduction, des arguments étayés par des références précises au texte, et une conclusion personnelle nuancée.
\end{enumerate}
\end{exercice}

\newpage

\coursec{Corrigés des exercices}

\coursec{Corrigés des exercices — Contraction de texte : rédiger un résumé}

\begin{corrige}[Exercice 1 — QCM : la contraction de texte]
\begin{enumerate}
\item \textbf{Réponse a)} : le résumé ne retient que les idées essentielles de l'auteur, sans aucune opinion personnelle. Les réponses b) et c) sont fausses : le résumé exclut le point de vue du rédacteur et supprime les exemples.
\item \textbf{Réponse b)} : le quart de 480 mots est $480 \div 4 = 120$ mots.
\item \textbf{Réponse b)} : remplacer une énumération de termes précis par un terme plus général est la \cle{généralisation}.
\item \textbf{Réponse b)} : une \cle{phrase simple} ne contient qu'un seul verbe conjugué (une seule proposition).
\end{enumerate}
\end{corrige}

\begin{corrige}[Exercice 2 — Vrai ou faux ?]
\begin{enumerate}
\item \textbf{Faux.} Dans un résumé, les exemples, anecdotes, chiffres et citations sont supprimés : seules les idées essentielles sont conservées.
\item \textbf{Vrai.} La \cle{cohésion} est assurée par les procédés de reprise (pronoms, reprises nominales, synonymes) et par les connecteurs logiques qui relient les phrases entre elles.
\item \textbf{Vrai.} Dans une \cle{progression thématique}, le même thème est repris et approfondi tout au long du texte, contrairement à la progression linéaire où chaque phrase apporte un élément nouveau.
\item \textbf{Faux.} La \cle{litote} consiste à dire moins pour faire entendre plus (ex. « ce n'est pas mal » pour « c'est bien »). L'exagération relève de l'\emph{hyperbole}.
\end{enumerate}
\end{corrige}

\begin{corrige}[Exercice 3 — Définitions]
\begin{enumerate}
\item \textbf{La contraction de texte} est un exercice qui consiste à réduire un texte à une fraction donnée de sa longueur (le quart, le tiers) en ne conservant que ses idées essentielles, sans ajouter d'opinion personnelle ni reprendre mot à mot les formulations de l'auteur. \emph{Exemple : résumer en 100 mots un article de 400 mots sur la lecture.}
\item \textbf{La thèse} d'un texte argumentatif est l'idée principale que l'auteur défend et cherche à faire admettre au lecteur. Elle est soutenue par des arguments et des exemples. \emph{Exemple : dans un article contre le portable en classe, la thèse est : « il faut interdire le téléphone portable en classe ».}
\item \textbf{La cohérence textuelle} est la qualité d'un texte dont toutes les parties s'organisent logiquement autour d'une même idée, sans contradiction. Elle repose sur l'unité de thème, l'ordre des idées et l'accord entre les temps et les personnes. \emph{Exemple : un texte qui annonce trois arguments et les développe tous dans l'ordre annoncé est cohérent.}
\item \textbf{L'euphémisme} est une figure d'atténuation qui consiste à exprimer une réalité brutale ou délicate avec des termes adoucis. \emph{Exemple : « il nous a quittés » pour dire « il est mort ».}
\end{enumerate}
\end{corrige}

\begin{corrige}[Exercice 4 — Calculs directs]
\begin{enumerate}
\item Résumé au quart : $600 \div 4 = 150$ mots. Résumé au tiers : $600 \div 3 = 200$ mots.
\item Marge de 10 \% de 120 mots : $120 \times \dfrac{10}{100} = 12$ mots. Minimum : $120 - 12 = 108$ mots ; maximum : $120 + 12 = 132$ mots. Le résumé doit donc contenir entre 108 et 132 mots.
\item Fraction représentée : $\dfrac{180}{450} = \dfrac{2}{5}$. Or le quart de 450 est $450 \div 4 = 112,5$ mots. Le résumé de 180 mots (les deux cinquièmes) est donc \textbf{trop long} : il n'est pas conforme à la consigne « au quart ».
\end{enumerate}
\end{corrige}

\begin{corrige}[Exercice 5 — Repérage dans un éditorial]
\begin{enumerate}
\item \textbf{Thèse} : l'auteur est contre le téléphone portable en classe. Phrase qui l'exprime : « Faut-il pour autant l'autoriser en classe ? Je réponds non, et je l'affirme sans détour. » (reprise dans la conclusion : « bannissons le portable de la salle de classe »).
\item \textbf{Arguments} : (1) le portable détourne l'attention ; (2) il favorise la triche ; (3) il nuit à la santé des adolescents. \textbf{Exemples} : l'enseignant de Yaoundé qui a confisqué douze appareils en une journée ; l'enquête de 2024 dans un lycée de Douala sur les réponses qui circulent par messagerie.
\item \textbf{Plan du texte} : I. Introduction : situation et question posée, annonce de la thèse. II. Développement : trois arguments (distraction, triche, santé). III. Réfutation de l'objection adverse (« Certes... Mais... »). IV. Conclusion : rappel de la thèse sous forme d'appel.
\item \textbf{Progression argumentative} : le texte est organisé pour défendre une thèse à l'aide d'arguments ordonnés par des connecteurs (« d'abord », « ensuite », « enfin ») et s'achève sur une conclusion qui reprend la thèse. Ce n'est pas une simple progression thématique, car l'auteur ne se contente pas de décrire le portable : il prend position et cherche à convaincre.
\end{enumerate}
\end{corrige}

\begin{corrige}[Exercice 6 — Réduction guidée]
\begin{enumerate}
\item « Les élèves qui téléphonent en cours ne suivent pas les explications. » (suppression des expansions non essentielles : « leur », « portable », « pendant les cours » devient « en cours », « attentivement » et « du professeur » sont supprimés). Technique : \cle{suppression}.
\item « Les véhicules encombrent les rues de nos villes. » Technique : \cle{généralisation} (l'énumération est remplacée par le terme générique « véhicules »).
\item « Les jeunes doivent comprendre que la lecture enrichit l'esprit. » (ou : « La lecture enrichit l'esprit et l'imagination. ») Technique : \cle{condensation} (plusieurs tournures lourdes sont resserrées en une formulation brève).
\item « Le directeur a annoncé aux parents que l'école fermerait pendant les travaux. » (ou : « Il leur a annoncé que l'établissement fermerait pendant les travaux. ») Technique : \cle{substitution} par un pronom ou une reprise nominale.
\item « Le sérieux des élèves entraîne la réussite aux examens. » Technique : \cle{nominalisation} (le verbe « travaillent » et le groupe « avec sérieux » deviennent le nom « le sérieux »).
\end{enumerate}
\end{corrige}

\begin{corrige}[Exercice 7 — Phrases simples, phrases composées et connecteurs]
\textbf{A. Transformations}
\begin{enumerate}
\item Phrase simple : « À l'entrée de Sidi dans le village, tous les regards se tournèrent vers elle. » (un seul verbe conjugué : « se tournèrent »).
\item Phrase simple : « Malgré son âge, Baroka demeure redoutable. » (un seul verbe conjugué : « demeure »).
\item Phrase composée (propositions possibles) : « Le photographe publia l'image de Sidi, et cette image fit le tour du pays. » ou « Le photographe publia l'image de Sidi, qui devint célèbre dans tout le village. »
\item Phrase composée : « Lakunle refusa de payer la dot, car il jugeait cette coutume dépassée. » ou « Lakunle refusa de payer la dot, ce qui scandalisa le village. »
\item Phrase simple : « En trompant Sadiku, le Bale provoqua la colère de Sidi. » (un seul verbe conjugué : « provoqua »).
\end{enumerate}
\textbf{B. Classement des connecteurs}
\begin{center}
\begin{tabularx}{\linewidth}{|l|X|}\hline
\rowcolor{popblueL} \textbf{Fonction logique} & \textbf{Connecteurs} \\ \hline
Addition & de plus ; en outre \\ \hline
Opposition & cependant ; néanmoins \\ \hline
Cause & car ; en effet \\ \hline
Conséquence & par conséquent ; c'est pourquoi \\ \hline
Concession & bien que \\ \hline
\end{tabularx}
\end{center}
\emph{Remarque} : « ensuite » exprime l'ordre chronologique (énumération) ; il n'entre dans aucune des cinq fonctions demandées, mais il est très utile pour organiser le plan d'un résumé.
\end{corrige}

\begin{corrige}[Exercice 8 — Figures d'atténuation]
\begin{enumerate}
\item \textbf{Litote} (« ce n'est pas un imbécile ») : dire moins (« pas imbécile ») pour faire entendre plus (« il est intelligent ») ; l'effet est une louange prudente et ironique.
\item \textbf{Euphémisme} (« a rejoint ses ancêtres ») : adoucir l'annonce de la mort pour ménager les auditeurs.
\item \textbf{Prétérition} (« je ne te parle pas de... ») : annoncer qu'on tait les dettes et les querelles, c'est en réalité les révéler ; effet accusateur déguisé.
\item \textbf{Litote} (« tu n'es pas sans savoir ») : la double négation affirme avec force que l'interlocuteur sait parfaitement ; effet d'insistance polie.
\item \textbf{Euphémisme} (« un peu fatigué » pour un malade en très mauvais état) : atténuer la gravité de la maladie pour rassurer ou ménager l'entourage.
\item \textbf{Prétérition} (« je passe sous silence... ») : feindre de taire les erreurs du notable tout en les signalant ; effet critique voilé.
\item \textbf{Litote} (« n'était pas mauvais ») : dire moins pour signifier que le repas était bon ; compliment mesuré.
\item \textbf{Euphémisme} (« il nous a quittés ») : formule adoucie pour annoncer un décès.
\item \textbf{Litote} (« ce n'est pas tous les jours ») : souligner l'exceptionnalité de la visite en niant sa fréquence.
\item \textbf{Prétérition} (« je ne mentionnerai pas le prix exorbitant ») : dénoncer le prix tout en prétendant le taire.
\end{enumerate}
\end{corrige}

\begin{corrige}[Exercice 9 — Contraction de texte type Baccalauréat technique]
\textbf{1. Thèse et arguments.} Thèse : il faut encadrer strictement l'usage du téléphone portable chez les élèves (« Il convient donc d'en encadrer strictement l'usage chez les élèves »). Arguments : (1) le portable est une source permanente de distraction qui fait baisser les résultats ; (2) il favorise la triche et le harcèlement ; (3) il nuit à la santé et à la vie sociale des adolescents.

\textbf{2. Plan du texte.} I. Introduction : constat et annonce de la thèse. II. Premier argument : la distraction (fonction : prouver). III. Deuxième argument : la malhonnêteté (fonction : prouver). IV. Troisième argument : la santé et la vie sociale (fonction : prouver). V. Concession et réfutation des avantages, proposition de solution, conclusion (fonction : nuancer puis convaincre).

\textbf{3. Résumé proposé (121 mots).}

\emph{Le téléphone portable, devenu quotidien chez les élèves camerounais, exige selon l'auteur un encadrement strict. D'abord, il distrait les élèves pendant les cours, ce qui nuit à l'assimilation des leçons et fait baisser les résultats scolaires. Ensuite, il encourage la malhonnêteté : il facilite la fraude aux évaluations et alimente le harcèlement entre camarades. Enfin, son usage excessif fatigue les adolescents, trouble leur sommeil et les éloigne du sport, de la lecture et de la vie familiale. Certes, le portable permet de joindre les parents en urgence et d'accéder à des ressources éducatives. Toutefois, ces avantages ne justifient pas sa présence en classe. L'auteur propose donc de l'interdire pendant les cours et de sensibiliser les jeunes à un usage responsable, afin que l'école reste un lieu de formation.}

\textbf{4.} Nombre de mots : 121 (conforme : $120 \pm 12$, donc entre 108 et 132).

\textbf{Barème indicatif (sur 20)} : thèse et arguments correctement dégagés (4 pts) ; plan avec fonctions (4 pts) ; respect de la longueur (2 pts) ; neutralité et reformulation sans reprise mot à mot (4 pts) ; cohésion (connecteurs, reprises) (3 pts) ; correction de la langue (3 pts).

\textbf{Points de vigilance} : citer la phrase exacte de la thèse ; ne conserver aucun exemple (Bafoussam, Douala) ; ne pas donner son avis ; indiquer le nombre de mots ; vérifier que la marge de 10 \% est respectée.
\end{corrige}

\begin{corrige}[Exercice 10 — Remédiation : corriger un résumé fautif]
\textbf{1. Les erreurs relevées.}
\begin{enumerate}
\item \textbf{Opinions personnelles} : « Moi, je trouve que ce texte a totalement raison », « Personnellement, je connais même un camarade... », « je pense que... ». Un résumé doit rester neutre : il rapporte les idées de l'auteur, pas celles du rédacteur.
\item \textbf{Exemples conservés} : le professeur de Bafoussam et le médecin de Douala sont des exemples ; ils doivent être supprimés.
\item \textbf{Longueur non respectée} : 165 mots au lieu de 120 maximum (soit 45 mots de trop).
\item \textbf{Conclusion inventée} : « interdits partout au Cameroun, et même confisqués définitivement » n'apparaît pas dans le texte ; l'auteur propose seulement l'interdiction pendant les cours et la sensibilisation.
\item \textbf{Reprise mot à mot} : « un professeur de mathématiques d'un lycée de Bafoussam affirme que la moitié de ses élèves possède un téléphone allumé en classe » est recopiée sans reformulation.
\end{enumerate}

\textbf{2. Corrections proposées.}
\begin{enumerate}
\item Supprimer toutes les marques de subjectivité (« moi », « je pense », « personnellement ») et rapporter les idées à la troisième personne (« l'auteur explique que... »).
\item Remplacer les exemples par l'idée générale qu'ils illustrent (la distraction, les troubles de la santé).
\item Réduire à 120 mots en supprimant les détails et en condensant les formulations.
\item Remplacer la conclusion inventée par la solution réelle du texte : interdiction en classe et sensibilisation à un usage responsable.
\item Reformuler avec ses propres mots (synonymes, nominalisations, changement de construction).
\end{enumerate}

\textbf{3. Version corrigée (environ 115 mots).}

\emph{Selon l'auteur, l'usage du téléphone portable chez les élèves doit être strictement encadré. D'abord, cet appareil distrait les élèves en classe et provoque une baisse des résultats scolaires. Ensuite, il favorise la fraude pendant les évaluations et alimente le harcèlement entre camarades. Enfin, son usage excessif nuit à la santé des adolescents et les éloigne du sport, de la lecture et de la vie familiale. L'auteur reconnaît toutefois que le portable permet de contacter les parents en urgence et d'accéder à des ressources éducatives. Il propose donc de l'interdire pendant les heures de cours tout en sensibilisant les élèves à un usage responsable, afin que l'école demeure un lieu de formation.}

\textbf{Points de vigilance} : vérifier systématiquement les quatre critères d'un bon résumé : fidélité à l'auteur, neutralité, longueur conforme, reformulation personnelle.
\end{corrige}

\begin{corrige}[Exercice 11 — Question de synthèse sur \emph{Le Lion et la perle}]
\textbf{1. Qui est Baroka ?} Baroka est le Bale, c'est-à-dire le chef du village d'Ilujinlè, au Nigeria. Surnommé « le Lion », il est âgé de soixante-deux ans, polygame, rusé et redouté. Il incarne l'autorité traditionnelle face à Lakunle, le jeune instituteur moderniste.

\textbf{2. Baroka défenseur de la tradition.} (a) Il défend la coutume de la dot, que Lakunle refuse de payer, et épouse Sidi selon les rites du village. (b) Il assume pleinement son statut de Bale polygame : ses nombreuses épouses, dont Sadiku la première, perpétuent l'organisation traditionnelle du lignage, et il affirme vouloir préserver les coutumes face à la modernité incarnée par le chemin de fer et l'école.

\textbf{3. Baroka fossoyeur de la tradition.} (a) Il ment sur sa prétendue impuissance pour piéger Sidi : en trahissant la parole sacrée du chef, il corrompt les valeurs qu'il prétend défendre. (b) Il instrumentalise Sadiku, sa première épouse, en lui confiant un faux secret, transformant la confiance conjugale traditionnelle en instrument de ruse égoïste.

\textbf{4. Développement construit (éléments attendus).}

\emph{Introduction} : dans \emph{Le Lion et la perle} (1959), Wole Soyinka met en scène la confrontation entre tradition et modernité au Nigeria. Baroka, le Bale d'Ilujinlè, apparaît comme le champion de la tradition ; mais ses manœuvres pour conquérir Sidi invitent à se demander s'il la défend ou s'il la trahit.

\emph{Développement} : d'un côté, Baroka incarne la tradition : chef polygame respecté, il défend la dot, les rites et la sagesse des anciens, et résiste habilement à la modernité symbolisée par le photographe et le projet de route. De l'autre, il fragilise cette tradition : son mensonge sur son impuissance, savamment exploité grâce à Sadiku, déshonore la parole du chef ; sa conquête de Sidi relève de la ruse égoïste plus que de l'honneur. La tradition devient entre ses mains un instrument de pouvoir personnel.

\emph{Conclusion nuancée} : Baroka est à la fois le vainqueur et le fossoyeur de la tradition : il la fait triompher contre Lakunle, mais en la vidant de sa morale. Soyinka suggère ainsi qu'une tradition qui ne se renouvelle pas et qui sert les appétits individuels risque de se détruire elle-même.

\textbf{Barème indicatif (sur 20)} : présentation de Baroka et du contexte (3 pts) ; arguments « défenseur » étayés (4 pts) ; arguments « fossoyeur » étayés (4 pts) ; construction du développement (introduction, articulations, conclusion nuancée) (5 pts) ; correction de la langue et qualité de l'expression (4 pts).

\textbf{Points de vigilance} : citer des faits précis de l'œuvre (la dot refusée par Lakunle, le faux secret confié à Sadiku, le surnom « le Lion ») ; éviter le hors-sujet sur la modernité seule ; la conclusion personnelle doit rester nuancée et s'appuyer sur l'analyse menée.
\end{corrige}

\newpage

\coursec{Fiche bilan}

\courssub{Carte mentale de la leçon}
\begin{popfigure}
\begin{center}\tcbox[colback=popink,colframe=popink,coltext=white,arc=9pt,boxsep=4pt,left=6pt,right=6pt]{\bfseries\small Contraction de texte \& Résumé}\end{center}
\vspace{-2pt}
\begin{tcbraster}[raster columns=3,raster equal height=rows,raster column skip=6pt,raster row skip=6pt,enhanced,blankest]
\begin{tcolorbox}[enhanced,colback=poporange,colframe=poporange,coltext=white,boxrule=0.9pt,arc=3pt,left=4pt,right=4pt,top=3pt,bottom=3pt,boxsep=1pt,halign=left]\bfseries\scriptsize Texte explicatif (structure, progression)\end{tcolorbox}
\begin{tcolorbox}[enhanced,colback=popgreen,colframe=popgreen,coltext=white,boxrule=0.9pt,arc=3pt,left=4pt,right=4pt,top=3pt,bottom=3pt,boxsep=1pt,halign=left]\bfseries\scriptsize Structure argumentative \& Plan du texte\end{tcolorbox}
\begin{tcolorbox}[enhanced,colback=poppurple,colframe=poppurple,coltext=white,boxrule=0.9pt,arc=3pt,left=4pt,right=4pt,top=3pt,bottom=3pt,boxsep=1pt,halign=left]\bfseries\scriptsize Cohésion \& Cohérence Phrase simple/composée\end{tcolorbox}
\begin{tcolorbox}[enhanced,colback=poppink,colframe=poppink,coltext=white,boxrule=0.9pt,arc=3pt,left=4pt,right=4pt,top=3pt,bottom=3pt,boxsep=1pt,halign=left]\bfseries\scriptsize Techniques de réduction \& Reformulation\end{tcolorbox}
\begin{tcolorbox}[enhanced,colback=popgold,colframe=popgold,coltext=white,boxrule=0.9pt,arc=3pt,left=4pt,right=4pt,top=3pt,bottom=3pt,boxsep=1pt,halign=left]\bfseries\scriptsize Figures d'atténuation Litote, Euphémisme, Prétérition\end{tcolorbox}
\begin{tcolorbox}[enhanced,colback=popblue,colframe=popblue,coltext=white,boxrule=0.9pt,arc=3pt,left=4pt,right=4pt,top=3pt,bottom=3pt,boxsep=1pt,halign=left]\bfseries\scriptsize Rédaction, Confrontation Amélioration du résumé\end{tcolorbox}
\begin{tcolorbox}[enhanced,colback=popmuted,colframe=popmuted,coltext=white,boxrule=0.9pt,arc=3pt,left=4pt,right=4pt,top=3pt,bottom=3pt,boxsep=1pt,halign=left]\bfseries\scriptsize Le Lion et la Perle Bilan de l'œuvre\end{tcolorbox}
\end{tcbraster}

\legende{Carte mentale récapitulative de la leçon 07 : les 7 piliers de la contraction de texte en Première technique.}
\end{popfigure}

\courssub{Bilan synthétique}

\begin{aretenir}[Bilan synthétique]

\textbf{Définitions clés}
\begin{itemize}
  \item \cle{Texte explicatif} : Texte dont la finalité est d'informer, d'expliquer un phénomène, un concept ou une procédure. Il s'appuie sur l'objectivité, la clarté, la structure logique (définition, classification, cause/effet, problème/solution).
  \item \cle{Contraction de texte (Résumé)} : Rédaction d'un texte plus court qui rend compte fidèlement des idées essentielles du texte source, sans en modifier le sens, en respectant l'ordre logique et en supprimant les détails, exemples, répétitions et commentaires subjectifs.
  \item \cle{Cohésion} : Ensemble des moyens linguistiques (connecteurs, anaphores, répétitions, champs lexicaux, temps verbaux) qui assurent la liaison \emph{formelle} entre les phrases.
  \item \cle{Cohérence} : Logique \emph{sémantique} du texte : organisation des idées, progression thématique, respect de la structure argumentative (thèses, arguments, exemples).
  \item \cle{Figures d'atténuation} : Procédés rhétoriques qui adoucissent l'énoncé : \textbf{Litote} (dire moins pour suggérer plus), \textbf{Euphémisme} (terme doux pour réalité dure), \textbf{Prétérition} (faire semblant de passer sous silence ce qu'on dit).
\end{itemize}

\textbf{Structure argumentative \& Plan du texte (à identifier avant de résumer)}
\begin{center}
\begin{tabularx}{\linewidth}{|l|X|X|}\hline
\rowcolor{popblueL}
\textbf{Type de plan} & \textbf{Caractéristiques} & \textbf{Mots-clés / Progression} \\ \hline
\cle{Thématique} & Enchaînement d'aspects d'un même thème (description, énumération). & D'abord, ensuite, enfin ; d'une part, d'autre part. \\ \hline
\cle{Dialectique} & Thèse -- Antithèse -- Synthèse (argumentation contradictoire). & Certes... mais... ; cependant ; en définitive. \\ \hline
\cle{Analytique} & Problème -- Causes -- Conséquences -- Solutions (explicatif). & Pourquoi ? Parce que ; causes : ... ; solutions : ... \\ \hline
\cle{Chronologique} & Suites d'événements dans le temps (récit, procédure). & D'abord, puis, finalement ; avant/après. \\ \hline
\end{tabularx}
\end{center}

\textbf{Méthode express : Les 5 étapes de la contraction}
\begin{enumerate}
  \item \textbf{Lecture active} : Identifier le \cle{sujet}, la \cle{thèse} principale, le \cle{type de plan} et les \cle{connecteurs} logiques.
  \item \textbf{Découpage \& Titrage} : Segmenter le texte en paragraphes logiques (macro-structure) et donner un titre à chaque séquence.
  \item \textbf{Sélection des idées essentielles} : Éliminer les exemples, détails, illustrations, reformulations, digressions. Ne garder que les \cle{arguments nucléaires}.
  \item \textbf{Reformulation \& Réduction} : Changer la forme (nominalisation, subordination, suppression de l'apostrophe, passage au discours indirect) sans trahir le sens. Utiliser les \cle{connecteurs} de la cohésion.
  \item \textbf{Rédaction finale \& Contrôle} : Rédiger le résumé en respectant la \cle{limite de mots} ($\pm$10\%). Vérifier : fidélité, cohérence, cohésion, correction linguistique, absence de « je ».
\end{enumerate}

\textbf{Points de langue : Phrase simple / Phrase composée}
\begin{itemize}
  \item \cle{Phrase simple} : Une seule proposition (un seul verbe conjugué). Ex : \emph{Le résumé exige de la rigueur.}
  \item \cle{Phrase composée} : Au moins deux propositions. Deux types :
  \begin{itemize}
    \item \textbf{Coordination} : Propositions indépendantes liées par \emph{mais, ou, et, donc, or, ni, car}. Ex : \emph{Le texte est long, mais le résumé est court.}
    \item \textbf{Subordination} : Proposition principale + proposition(s) subordonnée(s) (conjonctive, relative, infinitive, participiale). Ex : \emph{Il faut résumer le texte [pour réussir l'épreuve].}
  \end{itemize}
  \item \textbf{Astuce contraction} : Transformer des phrases coordonnées en subordonnées (gain de place, fluidité). Ex : \emph{Le texte explique X. Il argue Y.} $\rightarrow$ \emph{Le texte explique X en argumentant Y.}
\end{itemize}

\textbf{Le Lion et la Perle (Wole Soyinka) — Bilan}
\begin{itemize}
  \item \cle{Contexte} : Nigéria, village d'Ilujinle, confrontation Tradition (Baroka) vs Modernité (Lakunlé), Sidi comme enjeu.
  \item \cle{Thèmes majeurs} : \textbf{Pouvoir}/stratégie, condition féminine, choc des cultures, ruse vs naïveté, oralité/théâtre total.
  \item \cle{Structure} : 3 parties (Matin, Midi, Soir) = temps dramatique et symbolique. Langage : proverbes, chants, danse, métaphores yoruba.
\end{itemize}

\end{aretenir}

\courssub{Checklist avant le Bac}

\begin{aretenir}[Checklist avant le Bac — Contraction de texte]
\begin{itemize}
  \item $\square$ Je sais définir le texte explicatif et en reconnaître la structure (plan thématique, dialectique, analytique).
  \item $\square$ Je distingue \textbf{cohésion} (moyens formels : connecteurs, anaphores) et \textbf{cohérence} (logique des idées).
  \item $\square$ J'identifie les \textbf{progrèsions} (thématique, argumentative, narrative) dans un texte.
  \item $\square$ Je maîtrise la \textbf{méthode en 5 étapes} : Lecture $\rightarrow$ Découpage $\rightarrow$ Sélection $\rightarrow$ Reformulation $\rightarrow$ Contrôle.
  \item $\square$ J'applique les \textbf{techniques de réduction} : suppression (exemples, détails), généralisation, nominalisation, subordination, discours indirect.
  \item $\square$ Je repère et comprends les \textbf{figures d'atténuation} : litote, euphémisme, prétérition.
  \item $\square$ Je transforme des \textbf{phrases coordonnées en subordonnées} pour condenser l'information.
  \item $\square$ Je respecte la \textbf{contrainte de longueur} (nombre de mots imposé $\pm$ 10\%) et l'absence de subjectivité (« je », jugements de valeur).
  \item $\square$ Je relis mon résumé pour vérifier : \textbf{fidélité au sens}, \textbf{cohérence interne}, \textbf{cohésion linguistique}, \textbf{correction orthographique/grammaticale}.
  \item $\square$ Je connais le \textbf{bilan de \emph{Le Lion et la Perle}} : contexte, personnages, thèmes, structure en 3 temps, spécificités stylistiques (oralité).
  \item $\square$ Je suis capable de \textbf{confronter ma production} à une grille d'évaluation (contenu, forme, langue) et de l'améliorer.
\end{itemize}
\end{aretenir}

\findecours

\end{document}
